% ============================================================
%  IKT448 ICT Research Methods -- Course script, Autumn 2026
%  University of Agder -- Timo Amling
%  Built from the twelve session decks by tools/build_script.py
%  Compile: pdflatex x3
% ============================================================
\documentclass[11pt,a4paper,british,twoside=false,headings=small,chapterprefix=false,numbers=noenddot]{scrbook}
\usepackage[T1]{fontenc}
\usepackage[utf8]{inputenc}
\usepackage[british]{babel}
\usepackage{lmodern}
\usepackage{microtype}
\usepackage[a4paper,margin=2.6cm,top=2.8cm,bottom=2.8cm]{geometry}
\usepackage{graphicx}
\usepackage{xcolor}
\usepackage{tikz}
\usetikzlibrary{arrows.meta,positioning,shapes.geometric,calc,decorations.pathreplacing}
\usepackage{booktabs}
\usepackage{array}
\usepackage{tabularx}
\usepackage{enumitem}
\usepackage{amsmath,amssymb}
\usepackage{pgffor}
\usepackage[most]{tcolorbox}
\usepackage[hidelinks]{hyperref}
\usepackage{url}
\usepackage{scrlayer-scrpage}
\usepackage{caption}

% ---- colours (UiA brand book) ---------------------------------------
\definecolor{uiared}{HTML}{C8102E}
\definecolor{uiablue}{HTML}{131E29}
\definecolor{uiagrey}{HTML}{595959}
\definecolor{uiamutedred}{HTML}{FEE7E8}
\definecolor{uiamutedblue}{HTML}{E1EDF9}
\colorlet{kristiania}{uiared}
\colorlet{kgrey}{uiagrey}
\colorlet{ink}{uiablue}

% ---- headings -------------------------------------------------------
\addtokomafont{disposition}{\color{uiared}}
\addtokomafont{chapter}{\Huge}
\setkomafont{pageheadfoot}{\small\color{uiagrey}}
\pagestyle{scrheadings}
\clearpairofpagestyles
\ihead{IKT448 ICT Research Methods}
\ohead{\headmark}
\ofoot{\pagemark}
\automark[section]{chapter}

% ---- course constants ------------------------------------------------
\newcommand{\coursecode}{IKT448}
\newcommand{\coursename}{ICT Research Methods}
\newcommand{\coursesemester}{Autumn 2026}
\newcommand{\lecturername}{Timo Amling}
\newcommand{\lecturertitle}{Assistant Professor}
\newcommand{\lecturermail}{timo.d.amling@uia.no}
\newcommand{\institution}{University of Agder}
\newcommand{\faculty}{Faculty of Engineering and Science}

% ---- beamer compatibility -------------------------------------------
\newif\ifcolopen
\newenvironment{columns}[1][]{\par\noindent\colopenfalse\ignorespaces}%
  {\ifcolopen\end{minipage}\fi\par\vspace{0.4em}}
\newcommand{\column}[1]{\ifcolopen\end{minipage}\hfill\fi\begin{minipage}[t]{#1}\colopentrue}
\def\endcolumn{\ifcolopen\end{minipage}\hfill\colopenfalse\fi}  % \begin{column}{w}...\end{column} form
\tcbset{scriptbox/.style={enhanced, colback=gray!5, colframe=uiared, boxrule=0.5pt,
        arc=1.5pt, left=5pt, right=5pt, top=4pt, bottom=4pt, fonttitle=\bfseries,
        coltitle=white, colbacktitle=uiared, before skip=6pt, after skip=8pt}}
\newenvironment{block}[1]{\begin{tcolorbox}[scriptbox, title={#1}]}{\end{tcolorbox}}
\newenvironment{alertblock}[1]{\begin{tcolorbox}[scriptbox, colframe=uiared, colback=uiamutedred, title={#1}]}{\end{tcolorbox}}
\newenvironment{exampleblock}[1]{\begin{tcolorbox}[scriptbox, colframe=uiablue, colbacktitle=uiablue, colback=uiamutedblue, title={#1}]}{\end{tcolorbox}}
\newcommand{\answer}[1]{\par\vspace{0.2em}{\small\itshape\color{uiagrey}#1}\par\vspace{0.2em}}
\newcommand{\src}[1]{{\scriptsize\color{uiagrey}#1}}
\newcommand{\onionloc}[7]{}
\newcolumntype{L}[1]{>{\raggedright\arraybackslash}p{#1}}
\newcommand{\onionopt}{\fontsize{5.0pt}{5.9pt}\selectfont}
\newcommand{\onionlay}{\fontsize{5.6pt}{6.4pt}\selectfont\bfseries}
\newcommand{\onionsub}{\fontsize{4.6pt}{5.4pt}\selectfont}
\newcommand{\ldr}[1]{\draw[black!65, line width=0.4pt, dash pattern=on 1.5pt off 1.3pt] #1;}
\tikzset{
  kbox/.style={draw=uiared, thick, rounded corners=2pt, fill=uiared!6, align=center, inner sep=5pt, font=\small},
  gbox/.style={draw=uiagrey!70, thick, rounded corners=2pt, fill=gray!8, align=center, inner sep=5pt, font=\small},
  karrow/.style={-{Stealth[length=2.5mm]}, thick, uiared},
  garrow/.style={-{Stealth[length=2.5mm]}, thick, uiagrey},
}
\setlist[itemize]{itemsep=2pt, topsep=3pt, leftmargin=1.4em}
\setlist[enumerate]{itemsep=2pt, topsep=3pt, leftmargin=1.8em}
\renewcommand{\labelitemi}{\textcolor{uiared}{$\blacktriangleright$}}
\renewcommand{\labelitemii}{\textcolor{uiagrey}{$\bullet$}}

% ---- session chapter helpers ----------------------------------------
\newcommand{\sessionhead}[2]{\begin{tcolorbox}[enhanced, colback=uiamutedblue, colframe=uiamutedblue, arc=1.5pt, left=6pt, right=6pt]
  \small \textbf{Session #1} \;$\cdot$\; #2\end{tcolorbox}}
\newcommand{\lead}[1]{{\itshape\color{uiagrey}#1}\par\medskip}

\title{IKT448 ICT Research Methods}
\author{Timo Amling}
\begin{document}


\begin{titlepage}
  \begin{tikzpicture}[remember picture,overlay]
    \fill[uiared] (current page.north west) rectangle ([yshift=-9cm]current page.north east);
  \end{tikzpicture}
  \vspace*{1.2cm}
  {\color{white}\Huge\bfseries IKT448 ICT Research Methods\par}
  \vspace{0.6em}
  {\color{white}\Large Course script \;$\cdot$\; Autumn 2026\par}
  \vspace{0.6em}
  {\color{white}\large Master's programmes in Cyber Security and Cybersecurity Engineering\par}
  \vspace{5.5cm}
  {\Large Timo Amling\par}
  {\normalsize Assistant Professor, University of Agder\par}
  {\normalsize\texttt{timo.d.amling@uia.no}\par}
  \vfill
  \IfFileExists{uia-logo.png}{\includegraphics[height=1.4cm]{uia-logo}}{}
  \vspace{0.8cm}
  {\small\color{uiagrey}Version 1.0 \;$\cdot$\; 14 September 2026\par}
\end{titlepage}

\frontmatter


\chapter*{Preface}
\addcontentsline{toc}{chapter}{Preface}

This script accompanies the course IKT448 ICT Research Methods at the University of Agder in the autumn of 2026. It is not a substitute for the textbook --- Oates, Griffiths and McLean's \emph{Researching Information Systems and Computing} (2nd edition, 2022) remains the curriculum --- and it is not a transcript of the sessions. It sits between the two: it records what the sessions add to the book, in the order in which the sessions add it, so that the material can be read before a session and revisited afterwards.

The course is built around one piece of work, the essay, and one habit, the question \emph{how do they know?} Every chapter contributes a building block to the essay and ends with the questions we discuss at the end of the session, together with possible answers. The answers are not the only ones; the reasoning behind them is what matters.

Three conventions are used throughout. Definitions appear in shaded boxes and are stated once, before the term is used. Worked examples come from published studies, most of them recent and several of them from cyber security; the sources are listed at the end of each chapter, in APA style, and the papers are available in the course's reading folder. Numbers below ten are spelled out; British spelling is used.

Where a claim rests on a source, the source is named. Where I express a view --- on what examiners look for, on what a good essay does --- it is marked as such.

\medskip\noindent Timo Amling\\ Oslo, September 2026


\chapter*{How to use this script}
\addcontentsline{toc}{chapter}{How to use this script}

\section*{The course in twelve Thursdays}

Each chapter corresponds to one session. Read the assigned chapters of the textbook first, then the chapter of this script, then bring your questions to the session. The table lists the sessions, their dates and the reading.

\begin{center}\small
\renewcommand{\arraystretch}{1.15}
\begin{tabular}{@{}rlll@{}}
\toprule
& \textbf{Date} & \textbf{Session} & \textbf{Reading (Oates)}\\
\midrule
1 & 10.09 & Introduction & Ch.\ 1\\
2 & 17.09 & Approaches to research & Ch.\ 19--20\\
3 & 24.09 & Purpose, products and process of research & Ch.\ 1--3\\
4 & 01.10 & Formulating research questions & Ch.\ 2\\
5 & 08.10 & Reviewing the literature & Ch.\ 6\\
6 & 15.10 & Research strategies 1: survey, design and creation, experiment & Ch.\ 7--9\\
7 & 22.10 & Research strategies 2: case study, action research, ethnography & Ch.\ 10--12\\
8 & 29.10 & Quantitative data and analysis & Ch.\ 15, 17\\
9 & 05.11 & Qualitative data and analysis & Ch.\ 13--14, 16, 18\\
10 & 12.11 & Ethics: privacy and social responsibility & Ch.\ 4--5\\
11 & 19.11 & Producing the essay & Ch.\ 21\\
12 & 26.11 & Summing-up and supervision & everything\\
\bottomrule
\end{tabular}
\end{center}

\section*{Learning outcomes}

After completing the course, the student should (LO1) have an overview of qualitative and quantitative research methods within ICT research, with an emphasis on methods used within cyber security; (LO2) have an understanding of the use of formal methods and models, and the validity of these methods; (LO3) be able to assess research-ethical aspects, with an emphasis on privacy and social responsibility; (LO4) have insight into the chosen project and how cyber security fits into an overall ICT setting; (LO5) have an understanding of the structure of scientific investigations within ICT and cyber security, the relationship between breadth and depth, and what limitations such investigations have; and (LO6) be able to choose a suitable method in relation to the topic of a scientific investigation and problem solving within cyber security, such as a master's thesis.

\section*{Assessment}

The grade rests on three parts: presence and participation in discussions (20\,\%), a small project presented in a group (20\,\%), and the essay (60\,\%: written report 40\,\%, presentation and discussion in the final oral exam 20\,\%). Approved mandatory submissions and presentations are the condition for taking the exam; details and dates are announced separately. The essay has four building blocks --- introduction and question (Part I), relevant work (Part II), choice of methods (Part III), data, analysis, ethics and limitations (Part IV) --- and Chapter 11 presents it as a paper with six sections.

\section*{The book and the papers}

The textbook is the only required book. Chapter references follow the second edition (2022); a third edition (2025/26) is almost identical and mainly adds material on AI in research. The papers discussed in the sessions are listed at the end of each chapter and collected in the reading folder, numbered by the session in which they first appear.

\tableofcontents

\mainmatter

\part{Foundations}
\chapter{Introduction}
\label{ch:1}
\sessionhead{1}{Thursday 10 September 2026 \;$\cdot$\; reading: Oates, Chapter 1}

This chapter accompanies the first session. It does three things. First, it sets out how the course runs and how it is assessed: twelve Thursdays, one essay, and questions for discussion at the end of every session. Second, it answers the question that gives the course its direction --- what research is, and what it is not --- by tracing the definition through the history that produced it. Third, it introduces the reading material: one textbook, and one published study that serves as a map of the whole course. The chapter ends, like every chapter, with the questions for discussion and the sources used.


\section{What this course is for}

IKT448 is a methods course with one product: an essay that plans or evaluates a piece of research in cyber security. Everything in the twelve sessions serves that product, and everything in this script follows the order in which the essay is built. The course runs on Zoom on twelve Thursdays, each with a lecture in three blocks and a final block of questions for discussion. Participation in those discussions counts for 20\,\% of the grade, a small group project for another 20\,\%, and the essay for 60\,\%, of which 40\,\% is the written report and 20\,\% its presentation and discussion in the oral exam.

The course description sets out six learning outcomes. Two of them frame the whole course: an understanding of the structure of scientific investigations, their breadth, depth and limitations (LO5), and the ability to choose a suitable method for a topic in cyber security, such as a master's thesis (LO6). The remaining four --- an overview of qualitative and quantitative methods with a cyber security emphasis (LO1), an understanding of formal methods and models and their validity (LO2), the ability to assess research-ethical aspects with an emphasis on privacy and social responsibility (LO3), and insight into how cyber security fits into an overall ICT setting (LO4) --- are served by particular sessions. The introduction serves LO5 and LO6; each later chapter states which outcomes it serves.

A word on where you stand. A bachelor's degree certifies that a graduate knows the foundations and methods of a field and can apply them with justification. A master's degree adds independence and specialisation: methods applied to new problems, and a thesis that must argue its own choices. A doctorate adds the creation of new knowledge. The steps are the same pipeline, scaled rather than replaced, and this course supplies the method side of the master's thesis: the strategies, the methods, and the justification the thesis must contain.

\section{The essay in one paragraph}

The essay has two possible layouts. Layout A plans work to be done: a proposed or possible master's topic is developed into an introduction and motivation, a problem and research questions, relevant work, a justified choice of methods with strengths, weaknesses and bias, ethics and privacy, and a conclusion. Layout B evaluates work already done: a published paper or master's thesis is summarised and its methods, research questions, limitations, bias, validity and ethics are assessed. Both layouts are about ten pages, may be written alone or in a group of two or three, in English or Norwegian, and must state which AI tools were used and for what. Nobody is expected to answer their research questions in the essay --- only to justify how they would.

The design of the course rests on one mechanism. The questions at the end of every session are building blocks of layout A. A student who answers them for his or her own topic, week by week, has most of the essay drafted before the deadline, and has the method chapter of the master's thesis in hand.

\section{What research is}

\subsection{Three claims and one question}

Three claims circulate in every security organisation: that 95\,\% of breaches are caused by human error; that agentic AI will handle incident response autonomously; and that code review finds most defects before release. All three sound plausible, all three are repeated with confidence, and all three prompt the same question, which this course trains you to ask as a matter of habit: \emph{how do they know?}

The question has a structure. Information reaches us through a supply chain that begins with original studies in peer-reviewed journals and conferences, passes through syntheses such as systematic reviews and textbooks, continues through the professional layer of standards, industry and vendor reports, and ends in media and social feeds. Reach and speed increase along the chain; quality control and detail decrease. A qualified finding becomes a headline, and a headline becomes a slogan. Most claims reach us from the right end of the chain, and the task is to trace them back to the left.

\begin{center}\small
\begin{center}
    \resizebox{0.96\textwidth}{!}{%
    \begin{tikzpicture}
      \node[kbox, text width=2.6cm, minimum height=1.5cm] (a) at (0,0)
        {\textbf{Original studies}\\ \scriptsize peer-reviewed journals \& conferences};
      \node[kbox, text width=2.6cm, minimum height=1.5cm] (b) at (3.2,0)
        {\textbf{Syntheses}\\ \scriptsize systematic reviews, textbooks};
      \node[gbox, text width=2.6cm, minimum height=1.5cm] (c) at (6.4,0)
        {\textbf{Professional layer}\\ \scriptsize standards, industry \& vendor reports};
      \node[gbox, text width=2.6cm, minimum height=1.5cm] (d) at (9.6,0)
        {\textbf{Media}\\ \scriptsize tech journalism, blogs, keynotes};
      \node[gbox, text width=2.6cm, minimum height=1.5cm] (e) at (12.8,0)
        {\textbf{Feeds}\\ \scriptsize social media, forums, reposts};
      \draw[karrow] (a) -- (b); \draw[karrow] (b) -- (c);
      \draw[garrow] (c) -- (d); \draw[garrow] (d) -- (e);
      \draw[garrow] (0,-1.35) -- (12.8,-1.35);
      \node[font=\scriptsize, text=uiagrey, anchor=west] at (0,-1.7)
        {reach and speed increase $\rightarrow$};
      \draw[karrow] (12.8,-2.15) -- (0,-2.15);
      \node[font=\scriptsize, text=uiared, anchor=west] at (0,-2.5)
        {$\leftarrow$ quality control and detail increase};
    \end{tikzpicture}}
  \end{center}
  
  Each step to the right compresses and simplifies: a qualified finding
  becomes a headline, and a headline becomes a slogan. Most claims reach us
  from the \textbf{right end} of the chain; the task is to trace them back to
  the left.
\end{center}


Four questions do the tracing: (1) is there a primary source, as opposed to a retelling; (2) who produced it, with what independence, funding and incentive to reach this result; (3) how do they know --- with which method, sample and comparison; and (4) what was actually claimed, as opposed to what the repost says. Applied to the first claim, the trail runs from a LinkedIn post to a keynote slide to a vendor report and ends there, without a method section; the claim is unverifiable as stated. Applied to the second, the trail reaches a peer-reviewed review (Sapkota, Roumeliotis \& Karkee, 2025) whose method is visible --- a literature search and a conceptual taxonomy --- but in which incident response appears as an illustrative scenario, not as an evaluation. The same four questions therefore lead to different conclusions, and the difference lies not in the confidence of the claim but in the traceability of the evidence.

\subsection{A working definition}

\begin{block}{Research (working definition for this course)}
Research refers to the creation of new knowledge through a systematic process of collecting and analysing data, whose methods and findings are communicated so that others can evaluate and criticise them.
\end{block}

Three conditions follow, and removing any one of them means that what remains is no longer research. First, an original contribution: the study adds something not previously known --- at master's level, new to the field, or a known method applied to a new problem, with justification. Secondly, a systematic process: planned methods rather than anecdotes. Thirdly, openness to criticism: documentation that allows others to inspect the work. Oates, Griffiths and McLean (2022, Chapter 1) list the same three marks and contrast them with everyday thinking. For this audience the contrast is concrete: a penetration-test report satisfies the second and third conditions but rarely the first, because it produces knowledge about one client rather than knowledge beyond it.

\subsection{Where the definition comes from}

The definition is not arbitrary. Each of its parts was established against resistance over roughly 2,400 years: evidence over authority, method over intuition, and published, criticisable work over private conviction. The timeline runs from Aristotle's observation and logic through Ibn al-Haytham's systematic experiments, Bacon's empiricism, Descartes' doubt as method, the first scientific journals of 1665 and Newton's union of theory and experiment, to Popper's falsifiability in 1934, Kuhn's paradigms in 1962 and today's open science and replication movements.

\begin{center}\small
\begin{center}
    \resizebox{0.98\textwidth}{!}{%
    \begin{tikzpicture}
      \draw[very thick, uiagrey] (0,0) -- (15.6,0);
      \foreach \x in {0.6,2.4,4.2,6.0,7.8,9.6,11.4,13.2,15.0}{\fill[uiared] (\x,0) circle (0.09);}
      \node[gbox, text width=2.3cm, minimum height=1.25cm, font=\scriptsize] at (0.6,1.15)
        {\textbf{$\sim$350 BCE} Aristotle\\ observation + logic};
      \node[gbox, text width=2.3cm, minimum height=1.25cm, font=\scriptsize] at (4.2,1.15)
        {\textbf{1620} Bacon\\ empiricism: knowledge from evidence};
      \node[gbox, text width=2.3cm, minimum height=1.25cm, font=\scriptsize] at (7.8,1.15)
        {\textbf{1665} first journals\\ peer review begins};
      \node[gbox, text width=2.3cm, minimum height=1.25cm, font=\scriptsize] at (11.4,1.15)
        {\textbf{1934} Popper\\ falsifiability: claims must be testable};
      \node[gbox, text width=2.3cm, minimum height=1.25cm, font=\scriptsize] at (15.0,1.15)
        {\textbf{today}\\ open science, replication};
      \node[kbox, text width=2.3cm, minimum height=1.25cm, font=\scriptsize] at (2.4,-1.15)
        {\textbf{$\sim$1021} Ibn al-Haytham\\ systematic experiments};
      \node[kbox, text width=2.3cm, minimum height=1.25cm, font=\scriptsize] at (6.0,-1.15)
        {\textbf{1637} Descartes\\ doubt as method};
      \node[kbox, text width=2.3cm, minimum height=1.25cm, font=\scriptsize] at (9.6,-1.15)
        {\textbf{1687} Newton\\ theory + experiment united};
      \node[kbox, text width=2.3cm, minimum height=1.25cm, font=\scriptsize] at (13.2,-1.15)
        {\textbf{1962} Kuhn\\ paradigms: science moves in worldviews};
      \foreach \x in {0.6,4.2,7.8,11.4,15.0}{\draw[uiagrey] (\x,0.09)--(\x,0.5);}
      \foreach \x in {2.4,6.0,9.6,13.2}{\draw[uiagrey] (\x,-0.09)--(\x,-0.5);}
    \end{tikzpicture}}
  \end{center}
  
  The definition on the previous slide is not arbitrary. Each of its parts
  was established against resistance: evidence over authority, method over
  intuition, and published, \textbf{criticisable} work over private
  conviction.
\end{center}


Every rule in this course --- justify, document, cite, state limitations --- is a condensed lesson from this history. Empiricism reappears as your data (Sessions 8--9), method and doubt as the research process (Session 3), journals and peer review as the literature review (Session 5), falsifiability as questions and hypotheses (Session 4), paradigms as approaches (Session 2), and the replication crisis as protocols and limitations (Sessions 5 and 11). The history also produced the field this course belongs to: information systems, which grew out of both the technical and the social sciences.

\subsection{Where this course sits}

Laudon and Laudon (2022, section 1-3) describe the study of information systems as a multidisciplinary field with two traditions. The technical approach emphasises mathematically based models and the physical technology and formal capabilities of systems; its disciplines are computer science, management science and operations research. The behavioural approach is concerned with the issues that arise in the development and long-term use of systems --- strategic integration, design, implementation, utilisation and management; its disciplines are sociology, psychology and economics, and its focus is on attitudes, policy and behaviour rather than technical solutions. The sociotechnical view combines the two: optimal performance is achieved by jointly optimising the social and the technical system.

Information systems research therefore studies technology in organisations and draws on both traditions. Cyber security sits on the same seam, with technical controls on one side and human behaviour and organisational practice on the other. Most master's theses in this field need methods from both, which is why Sessions 8 and 9 cover quantitative and qualitative data.

\subsection{Who decides what counts as knowledge}

For most of history the decisive question was not whether a claim was true but whether it was permitted. Galileo was forced to recant in 1633; the Catholic Church acknowledged the error in 1992. The Enlightenment made openness to criticism a virtue: claims earn trust by withstanding scrutiny, not by being licensed. Today the filter has moved a second time. Publication is unrestricted, and evaluation has become the scarce resource: millions of papers appear each year, no individual can verify more than a small fraction of them, and trust therefore attaches to a system rather than to persons. Disinformation and organised opposition to evidence-based science operate in the third position of this sequence --- content ranked by engagement, not by evidence.

\begin{block}{Peer review}
Peer review refers to the assessment of a manuscript, before publication, by independent experts in the same field, usually anonymously: is the method sound, does the evidence support the claims, and is the contribution new?
\end{block}

Peer review guarantees a minimum standard --- the method was examined by independent experts before anyone else saw it --- but it does not guarantee truth. Errors, and occasionally fraud, pass review, which is why replication and reading the method yourself still matter. In this course you take both roles: as an author, whose essay is defended in the oral exam, and as a reviewer, in the small project and the discussions.

\subsection{Why trust science, if it is never complete}

The distinguishing mark of science is not being right but changing in response to evidence. The round-earth theory was refined by every new measurement, from Eratosthenes' circumference to Newton's predicted flattening, the oblate spheroid confirmed by expeditions, the geoid measured by satellites, and the GPS that works because of it. The flat-earth claim has stood still since antiquity. Science revises its claims when evidence contradicts them, makes risky and checkable predictions, and states its own limits; pseudo-science is immune to evidence, explains everything after the fact, and claims certainty.

Three consequences follow. First, in science a theory denotes the highest status a claim can attain --- extensively evidenced --- and not ``just a guess''. Secondly, science is never complete, and neither will your research be; it must be honest, documented, and one step forward. Thirdly, the question to ask of any claim is not whether it is certain but whether it changes under evidence, and whether you can check how.

\subsection{Claims, and when they are testable}

\begin{block}{Claim (thesis)}
A claim refers to a statement that asserts something is the case and that can, in principle, be true or false.
\end{block}

A claim is testable when three conditions hold: its terms are precise, it predicts something observable, and evidence could in principle refute it. ``Agentic AI shortens incident response time'' is testable once response time and a baseline are defined; ``our architecture is elegant'' is not testable as stated, because elegance predicts nothing observable; ``zero trust is better'' is not testable yet, and becomes so once ``better'' is operationalised as fewer successful lateral movements or lower incident cost. A claim that fails one of the three conditions is not wrong; it is unfinished. Session 4 turns this into craft: a research question is a testable claim in the form of a question.

\subsection{What research is not}

Four familiar statements mark the boundary. ``I built a tool and it works'' has no question, no systematic evaluation and no new knowledge. ``I searched the web and summarised'' collects opinions rather than generating evidence. ``Our penetration test found twelve vulnerabilities'' is valuable practice without a generalisable knowledge claim. ``Everyone on the team agrees'' offers consensus, which is not evidence. However, each of these can become research: build the tool and evaluate it systematically, or turn the penetration test into a study of why those vulnerabilities recur. Oates calls this the step from ``just developing a system'' to design and creation research (Session 6), and many cyber security master's theses lie on exactly this line.

\section{Relevance}

\begin{block}{Evidence-based practice}
Evidence-based practice refers to making professional decisions on the basis of the best available evidence --- from research and from systematically collected data --- rather than on habit, authority or marketing.
\end{block}

Security is already an evidence-based field: threat-intelligence reports, vulnerability research, incident statistics and post-incident reviews are its everyday material. It performs well when it runs on evidence and poorly when it runs on fear. This course teaches you to produce such evidence and to judge it, and there are three situations in which you will need both. First, your master's thesis: the essay written in this course can become its method chapter. Secondly, reading claims critically --- vendor white papers, benchmark blogs and ``studies show'' headlines --- with the question \emph{how do they know?} as a habit. Thirdly, arguing at work: ``we should migrate to X'' is persuasive when you can design the small study that would settle it.

\section{Reading material}

The course uses one book. Oates, Griffiths and McLean's \emph{Researching Information Systems and Computing} (2nd edition, 2022) is written for our field --- every example comes from information systems and computing --- and is organised around the 6Ps of research (Session 3). Every strategy and every method has its own chapter, and every chapter ends with evaluation criteria, which are the lens through which your essay and small project are graded. Chapter references in this script follow the second edition, the one available through the UiA library; a third edition (2025/26) is almost identical and mainly adds material on AI in research. Read the assigned chapter before the session: the first slides of every session ask questions the reading answers, and the cost of skipping grows every week, because each session builds on terms the reading introduced.

One published study serves as a map of the course. Sapkota, Roumeliotis and Karkee (2025) address a problem --- the terms ``AI agent'' and ``agentic AI'' are used interchangeably, which leads to conceptual ambiguity and misaligned system design --- with a method, a conceptual review based on a literature search across twelve platforms and a comparative analysis along dimensions such as autonomy, coordination, memory and interaction. Their finding is a taxonomy: AI agents are single, task-specific, tool-augmented systems, whereas agentic AI denotes systems of several specialised agents with orchestration, shared memory and goal decomposition. The implication is a shared vocabulary for design and evaluation, with causality, coordination failure and security as open challenges. The pattern problem $\rightarrow$ method $\rightarrow$ finding $\rightarrow$ implication recurs in every paper you will read this semester, and in your own essay.

The components of that study map onto the sessions of this course: purpose and problem (Session 3), research question (Session 4), literature (Session 5), strategy (Sessions 6--7), data and analysis (Sessions 8--9), ethics (Session 10) and write-up (Session 11). Every session is one row of this table, applied to your own essay and master's thesis. The same project can be run twice --- as a development project that ships a password-manager onboarding flow, and as a research project that ships the same flow and investigates where users abort, and why. The research version is more useful to the client, not less, because it explains why the numbers move; and it is what a master's thesis requires.


\section{Worked examples and tables}

The examples of the session are reproduced here in the form used on the slides, so that they can be reread without the deck.

\subsection{Tracing two claims to their source}
The four questions of Section 1.3 applied to two of the three claims. The same procedure yields different conclusions, and the difference lies in the traceability of the evidence, not in the confidence of the claim.
\begin{center}\small
Four questions to ask of any claim:
  \begin{enumerate}\setlength{\itemsep}{0.15em}
    \item \textbf{Primary source?} --- locate the original document, not a retelling
    \item \textbf{Who produced it?} --- independence, funding, and the incentive to reach this result
    \item \textbf{How do they know?} --- method, sample, comparison
    \item \textbf{What was actually claimed?} --- as opposed to what the repost says
  \end{enumerate}
  
  \footnotesize
  \renewcommand{\arraystretch}{1.3}
  \begin{tabular}{@{}p{6.3cm}p{6.5cm}@{}}
    \toprule
    \textbf{``95\,\% of breaches are human error''} & \textbf{``Agentic AI will handle incident response''}\\
    \midrule
    LinkedIn post $\rightarrow$ keynote slide $\rightarrow$ vendor report $\rightarrow$ \textbf{no method section}: the trail ends & Vendor blog $\rightarrow$ peer-reviewed review (Sapkota et al., 2025) $\rightarrow$ method visible: literature search and conceptual taxonomy; incident response appears as an \textbf{illustrative scenario}, not as an evaluation\\
    Assessment: unverifiable as stated & Assessment: a documented design proposal; the empirical evidence is still to be produced\\
    \bottomrule
  \end{tabular}
  

  {\color{uiagrey}\small The same four questions lead to different
  conclusions. The difference lies not in the confidence of the claim but in
  the \textbf{traceability} of the evidence.}
\end{center}


\subsection{Science and pseudo-science}
The demarcation in practical clothes. The table is Popper's argument reduced to three rows; the three consequences below it are what the argument means for a master's thesis.
\begin{center}\small
\small
  \renewcommand{\arraystretch}{1.25}
  \begin{tabular}{@{}p{6.2cm}p{6.4cm}@{}}
    \toprule
    \textbf{Science} & \textbf{Pseudo-science}\\
    \midrule
    Revises its claims when evidence contradicts them & Immune to evidence --- always the same answer\\
    Makes risky, checkable predictions (GPS works \emph{because} the earth's shape was refined for 2{,}000 years) & Explains everything \emph{after} the fact (the flat earth never moved under evidence)\\
    States its own limits and uncertainty & Claims certainty\\
    \bottomrule
  \end{tabular}
  

  Three consequences:
  \begin{enumerate}\setlength{\itemsep}{0.15em}
    \item In science, \textbf{theory} denotes the \emph{highest} status a
          claim can attain --- extensively evidenced --- not ``just a guess''.
    \item Science is \textbf{never complete}; neither will your research be,
          nor mine. It must be honest, documented, and one step forward:
          \textbf{research progresses over time}.
    \item The question to ask of any claim is therefore not ``is it certain?''
          but \textbf{``does it change under evidence --- and can I check
          how?''}
  \end{enumerate}
\end{center}


\subsection{Testable and untestable claims}
Three claims from a security workplace, tested against the three conditions of precision, observability and refutability.
\begin{center}\small
\small
  \renewcommand{\arraystretch}{1.5}
  \begin{tabular}{@{}p{6.4cm}p{6.5cm}@{}}
    \toprule
    ``Agentic AI shortens incident response time'' & \textbf{Testable} --- once response time and a baseline are defined. Sapkota et al.\ (2025) propose the scenario; nobody has measured it yet\\
    ``Our architecture is elegant'' & \textbf{Not testable} as stated --- ``elegant'' predicts nothing observable\\
    ``Zero trust is better'' & \textbf{Not testable \emph{yet}} --- it becomes testable once ``better'' is operationalised: fewer successful lateral movements? lower incident cost?\\
    \bottomrule
  \end{tabular}
  

  {\color{uiagrey}Session 4 turns this into craft: a research question is a
  testable claim in the form of a question.}
\end{center}


\subsection{What research is not}
\begin{center}\small
\small
  \renewcommand{\arraystretch}{1.35}
  \begin{tabular}{@{}p{4.4cm}p{9.2cm}@{}}
    \toprule
    \textbf{This alone\dots} & \textbf{\dots is not research, because\dots}\\
    \midrule
    ``I built a tool and it works'' & No question, no systematic evaluation, no new knowledge\\
    ``I searched the web and summarised'' & Collecting opinions is not generating evidence\\
    ``Our penetration test found 12 vulnerabilities'' & Valuable practice --- but no generalisable knowledge claim\\
    ``Everyone on the team agrees'' & Consensus is not evidence\\
    \bottomrule
  \end{tabular}
  
  \begin{block}{However}
    Each of these can \textbf{become} research: build the tool \emph{and}
    evaluate it systematically; turn the penetration test into a study of
    \emph{why} those vulnerabilities recur. Oates calls this the step from
    ``just developing a system'' to \textbf{design and creation research} ---
    session 6. Many cyber security master's theses lie on exactly this line.
  \end{block}
\end{center}


\subsection{The components of any study}
The study of Sapkota et al.\ (2025) decomposed into the components that the sessions of this course treat one by one.
\begin{center}\small
\small
  \renewcommand{\arraystretch}{1.3}
  \begin{tabular}{@{}p{3.1cm}p{6.6cm}p{3.2cm}@{}}
    \toprule
    \textbf{Component} & \textbf{In Sapkota et al.\ (2025)} & \textbf{In this course}\\
    \midrule
    Purpose and problem & Two terms conflated; a taxonomy is needed & Session 3\\
    Research question   & How do AI agents and agentic AI differ in architecture, autonomy and application? & Session 4\\
    Literature          & Search across twelve platforms, Boolean queries & Session 5\\
    Strategy            & Conceptual review; comparative taxonomy & Sessions 6--7\\
    Data and analysis   & Comparison tables across dimensions; application cases & Sessions 8--9\\
    Ethics              & AI writing assistance declared; funding disclosed; open access & Session 10\\
    Write-up            & The \emph{Information Fusion} article itself & Session 11\\
    \bottomrule
  \end{tabular}
  

  Every session of this course is one row of this table --- applied to
  \textbf{your own} essay and master's thesis.
\end{center}


\subsection{The same project, twice}
\begin{center}\small
\small
  \renewcommand{\arraystretch}{1.35}
  \begin{tabular}{@{}p{6.3cm}p{6.5cm}@{}}
    \toprule
    \textbf{As a development project} & \textbf{As a research project}\\
    \midrule
    ``We build a password manager onboarding flow for the client.'' & ``We build the flow --- \emph{and investigate} where users abort, and why.''\\
    Success = it ships and works & Success = it ships \emph{and} we can say something new and evidenced about onboarding drop-off\\
    Report: what we built and how & Report: what we built, how we studied it, what we found, what it means beyond this client\\
    \bottomrule
  \end{tabular}
  

  Same code, same client, same semester. The right column is what this
  course adds, and what a master's thesis requires. The research version is
  \emph{more} useful to the client, not less: it explains \emph{why} the
  numbers move.
\end{center}



\section*{Questions for discussion}
\begin{enumerate}\setlength{\itemsep}{0.7em}
    \item \textbf{The three claims} from the start --- human error, agentic
          AI, code review: which one would you accept, and what would it take
          to check it?
          \answer{The human-error figure is commonly traced to a 2014 IBM
          report without a published method, and ``involve'' is not
          ``cause''. The agentic-AI claim rests on a peer-reviewed review whose
          incident-response case is an illustrative scenario, not an
          evaluation. The code-review claim becomes testable once ``most
          defects'' and ``before release'' are defined.}
    \item \textbf{Your thesis} --- what could it investigate, and how would
          anyone know the answer?
          \answer{On the first day, ``not yet testable'' is the usual state.
          The next steps are the same for every topic: state what is unknown,
          who needs to know it, and what would count as an answer. Session 4
          turns this into a research question.}
    \item \textbf{Access} --- does unrestricted access to information
          democratise knowledge, or bury it?
          \answer{Both. Access widens the supply of claims, not their
          evaluation. Evaluation --- asking how the authors know --- is the
          skill this course trains.}
  \end{enumerate}

\section*{Key points}
\begin{itemize}
  \item Research = new knowledge, systematic, open to criticism
  \item Evidence-based practice: \emph{how do they know?}
  \item Course literature: Oates, Griffiths \& McLean (2022)
  \item Assessment: participation 20, small project 20, essay 60
\end{itemize}
\medskip\noindent\textbf{Six terms from this session:} research $\cdot$ evidence-based practice $\cdot$ original contribution $\cdot$ rigour $\cdot$ relevance $\cdot$ essay layout A (the research proposal in all but name)

\section*{Sources}
\begin{itemize}\small
  \item Oates, B.\,J., Griffiths, M., \& McLean, R. (2022). \emph{Researching information systems and computing} (2nd ed.). SAGE. --- Chapter 1.
  \item Sapkota, R., Roumeliotis, K.\,I., \& Karkee, M. (2025). AI agents vs.\ agentic AI: A conceptual taxonomy, applications and challenges. \emph{Information Fusion}, \emph{126}, 103599. \url{https://doi.org/10.1016/j.inffus.2025.103599}
  \item Oates, B.\,J., Griffiths, M., \& McLean, R. (2022). \emph{Researching information systems and computing} (2nd ed.). SAGE. --- Chapters 19 (Philosophical paradigms: positivism) and 20 (Alternative philosophical paradigms).
  \item Laudon, K.\,C., \& Laudon, J.\,P. (2022). \emph{Management information systems: Managing the digital firm} (17th ed., Global ed.). Pearson. --- Chapter 1, section 1-3.
\end{itemize}

\chapter{Approaches to Research}
\label{ch:2}
\sessionhead{2}{Thursday 17 September 2026 \;$\cdot$\; reading: Oates, Chapters 19--20}

Every study rests on assumptions about what can be known and how. This chapter makes those assumptions visible. It distinguishes the quantitative, qualitative and mixed approaches, presents the worldviews behind them, and orders all of the decisions a study requires with the research onion. Two guidelines for mixed methods research and four current studies of ChatGPT adoption serve as examples. The chapter matters for Part III of the essay: the approach you choose must be justified, and this chapter supplies the vocabulary for the justification.


\section{Warm-up}

A headline such as ``fitness apps make people healthier'' can be approached in at least four ways, and each carries assumptions about what evidence is. Sending a questionnaire to hundreds of users counts, compares and generalises. Interviewing people who quit asks what happened and why. Looking at the usage logs and seeing when people drop off measures behaviour. Doing a bit of all three combines them. These are not just different methods; they are different assumptions about what knowledge is, and that is what the word ``approach'' means in research.

Last session settled that research is the creation of new knowledge through a systematic process, open to criticism. This session asks how that process is set up, and the answer depends on the question. Part III of the essay --- the choice of methods --- has to justify the approach, and this chapter supplies the vocabulary.

\section{Three approaches}

\begin{block}{Quantitative research}
Quantitative research refers to an approach for testing objective theories by examining the relationships among measurable variables, yielding numerical data that are analysed with statistics.
\end{block}
\begin{block}{Qualitative research}
Qualitative research refers to an approach for exploring and understanding the meaning that individuals or groups ascribe to a social or human problem, with data collected in the participants' setting and analysed inductively.
\end{block}
\begin{block}{Mixed methods research}
Mixed methods research refers to an approach that collects both quantitative and qualitative data and integrates the two, so that the combined understanding exceeds what either strand provides alone.
\end{block}

Creswell's point is that these are not two boxes but the ends of a continuum: a study sits somewhere on it, and it tends towards one end more than the other. The core ideas differ. A quantitative study works deductively, from theory to hypothesis to measurement to test, and aims at generalisation and replication; its typical questions are which factors, how many, and whether one thing affects another. A qualitative study works inductively, from data in a real setting towards concepts and explanation, and aims at depth and meaning; its typical questions are how people do something and what it means to them. A mixed-methods study combines the two deliberately, for instance in an explanatory sequential design that runs a survey first and then interviews a smaller sample to explain the pattern.

Four recent studies illustrate the ends of the continuum. Biloš and Budimir (2024) surveyed 694 Generation Z respondents in Croatia with an extended UTAUT2 model, analysed the data with factor analysis and hierarchical regression, and found that performance expectancy, social influence, hedonic motivation, habit and personal innovativeness predict the intention to use ChatGPT, whereas effort expectancy, facilitating conditions and price value do not. Huy et al.\ (2024) surveyed 671 ChatGPT users in Vietnam with UTAUT2 integrated with task--technology fit, used covariance-based structural equation modelling, and found that most factors affect use while effort expectancy, social influence and trust do not, and that use drives continuance and recommendation. Both are deductive and postpositivist; the same theory family, two samples, and partly different results --- which is what replication looks like in practice.

In contrast, Mayer, Baygi and Buwalda (2025) studied one consultancy with semi-structured interviews, first with entry-level professionals and later with senior consultants and managers, using job crafting as a theoretical lens. They found that entry-level professionals use generative AI to craft tasks, relationships and the meaning of their work, and that a new form appears --- signal crafting, the use of AI to shape how others perceive the value of one's work. The study is inductive and set in a real context, with meaning as its object; the theory is a lens rather than a hypothesis, and the study extends it rather than tests it.

Mixed methods have their own guidelines in our field. Venkatesh, Brown and Bala (2013) ask three questions: whether a mixed methods approach is appropriate for the question at all; how meta-inferences --- the conclusions that integrate both strands --- are developed; and how their quality is validated. They list seven purposes for combining the strands: complementarity, completeness, developmental sequencing, expansion, corroboration or confirmation, compensation, and diversity. Venkatesh, Brown and Sullivan (2016) turn the guidelines into fourteen design properties with a decision tree, from the research questions and the paradigm stance to the mixing strategy, the timing of the strands and the validation of inferences. For your essay the consequence is simple: a mixed-methods proposal states its purpose in these terms, the sequence of the strands, which strand has priority, and where the two are integrated. Without these four decisions it is two studies, not one.

\section{Worldviews}

Every study rests on assumptions about what exists and how it can be known, whether or not the author states them. Creswell calls these assumptions worldviews; Oates (Chapters 19--20) calls them philosophical paradigms.

\begin{block}{Worldview (paradigm)}
A worldview refers to a basic set of beliefs that guide action: an ontology, or what exists, and an epistemology, or how we can know it. The assumptions are usually implicit; the essay makes them explicit.
\end{block}

Four worldviews carry most of the weight. Postpositivism assumes an objective reality that can be measured imperfectly, favours deductive testing, and produces claims such as ``which factors predict abandonment of fitness apps''. Constructivism assumes multiple realities constructed through experience, favours inductive understanding, and produces claims about what quitting an app means to the person who quits. The transformative worldview starts from politics and power, researches with marginalised groups, and aims at change. Pragmatism starts from the problem, uses whatever works, and underlies most mixed-methods research. Creswell's framework connects the worldview to a research design and to methods: the three components together constitute the approach.

\subsection{The research onion}

Saunders, Lewis and Thornhill (2019) order the decisions a study requires as layers of an onion, read from the outside in: philosophy, approach to theory development, methodological choice, strategy, time horizon, and techniques and procedures.

\begin{center}\small
\begin{center}
\begin{tikzpicture}[every node/.style={align=center, text=ink, inner sep=1pt}, scale=0.86, transform shape]
\path[use as bounding box] (-4.85,-3.45) rectangle (8.60,2.95);
\fill[uiared!30!white, draw=black!55, line width=0.4pt] (0.00,0) ellipse (4.72 and 2.85);
\fill[uiared!24!white, draw=black!55, line width=0.4pt] (-0.24,0) ellipse (4.15 and 2.54);
\fill[uiared!18!white, draw=black!55, line width=0.4pt] (-0.59,0) ellipse (3.46 and 2.17);
\fill[uiared!12!white, draw=black!55, line width=0.4pt] (-0.96,0) ellipse (2.74 and 1.80);
\fill[uiared!6!white] (-1.81,0) ellipse (1.61 and 1.04);
\draw[uiared!85!black, line width=0.6pt, dash pattern=on 1.5pt off 1.3pt] (-1.81,0) ellipse (1.61 and 1.04);
\fill[white, draw=black!55, line width=0.4pt] (-2.03,0) ellipse (1.25 and 0.44);
\node[font=\onionopt] at (-0.55,2.62) {Positivism};
\node[font=\onionopt] at (2.84,2.00) {Critical\\realism};
\node[font=\onionopt] at (4.10,-0.15) {Interpretivism};
\node[font=\onionopt] at (3.10,-1.85) {Post-\\modernism};
\node[font=\onionopt] at (-0.55,-2.62) {Pragmatism};
\node[font=\onionopt] at (-0.54,2.20) {Deduction};
\node[font=\onionopt] at (3.03,1.00) {Abduction};
\node[font=\onionopt] at (-0.54,-2.20) {Induction};
\node[font=\onionopt] at (-0.89,1.87) {Mono-method\\quantitative};
\node[font=\onionopt] at (1.49,1.35) {Mono-method\\qualitative};
\node[font=\onionopt] at (2.20,0.55) {Multi-method\\quantitative};
\node[font=\onionopt] at (2.20,-0.55) {Multi-method\\qualitative};
\node[font=\onionopt] at (1.49,-1.35) {Mixed-method\\simple};
\node[font=\onionopt] at (-0.89,-1.87) {Mixed method\\complex};
\node[font=\onionopt] at (-1.89,1.42) {Experiment};
\node[font=\onionopt] at (-0.03,1.42) {Survey};
\node[font=\onionopt] at (0.20,0.90) {Archival research};
\node[font=\onionopt] at (0.72,0.35) {Case study};
\node[font=\onionopt] at (0.72,-0.35) {Ethnography};
\node[font=\onionopt] at (0.20,-0.90) {Action research};
\node[font=\onionopt] at (-1.92,-1.42) {Narrative\\inquiry};
\node[font=\onionopt] at (-0.08,-1.42) {Grounded\\theory};
\node[font=\onionopt] at (-1.81,0.62) {Cross-sectional};
\node[font=\onionopt] at (-1.81,-0.62) {Longitudinal};
\node[font=\onionsub] at (-2.03,0.00) {Data collection\\and data analysis};
\ldr{( 1.20, 2.62) -- (5.05, 2.62)}
\node[font=\onionlay, anchor=west] at (5.15, 2.62) {Philosophy};
\ldr{( 1.30, 2.20) -- (5.05, 2.20)}
\node[font=\onionlay, anchor=west] at (5.15, 2.14) {Approach to theory\\development};
\ldr{( 0.90, 1.78) -- (5.05, 1.78) -- (5.05, 1.30)}
\node[font=\onionlay, anchor=west] at (5.15, 1.30) {Methodological\\choice};
\ldr{( 1.30, 0.02) -- (5.05, 0.02)}
\node[font=\onionlay, anchor=west] at (5.15, 0.02) {Strategy (ies)};
\ldr{(-0.70,-0.62) -- (3.30,-2.55) -- (5.05,-2.55)}
\node[font=\onionlay, anchor=west] at (5.15,-2.55) {Time horizon};
\ldr{(-2.70,-0.36) -- (-2.70,-3.10) -- (3.55,-3.10)}
\node[font=\onionlay, anchor=west] at (3.65,-3.10) {Techniques and\\procedures};
\draw[-{Stealth[length=2.4mm]}, uiared, line width=0.9pt] (7.95,2.62) -- (7.95,-3.10);
\node[font=\fontsize{6pt}{7pt}\selectfont, text=uiared, rotate=-90] at (8.35,-0.24) {read from the outside in};
\end{tikzpicture}
  \end{center}
\end{center}


The first layer is the philosophy. Positivism assumes one observer-independent reality and leads to hypothesis testing with large samples; critical realism assumes real structures that are observable only indirectly and searches for generative mechanisms; interpretivism assumes socially constructed realities and reconstructs meaning through in-depth interviews; postmodernism treats reality as an effect of language and power; pragmatism judges reality by problem solving and lets the method follow the question. Oates covers positivism, interpretivism and critical research; Saunders adds the other three. The second layer is the approach to theory development --- deduction from theory to test, typical of UTAUT and TAM studies; induction from data to theory; abduction, moving between the two from a surprising observation. The third is the methodological choice: mono-method with a single type of data, multi-method with several types within one paradigm, and mixed-method combining quantitative and qualitative data, in a simple sequential form or a complex repeated one. The fourth layer holds the strategies --- experiment, survey, archival research, case study, ethnography, action research, grounded theory and narrative inquiry; the fifth the time horizon, cross-sectional or longitudinal; and the sixth the techniques of collection and analysis.

The purpose of the onion is methodological coherence: the outer layers constrain the inner ones, so that an interpretivist philosophy combined with a large-scale survey needs explaining. The dependency does not run through all six layers, however. Layers 4 and 5 are independent: a case study can be cross-sectional or longitudinal, and so can a survey. Overall, the onion makes decisions explicit and open to criticism, which is what separates research from everyday reasoning. It is not the method that makes a study scientific but the traceable justification of the choice --- and that justification is Part III of your essay, layer by layer.

\section{One question, three studies}

The same question can be studied three ways, and the exercise shows how the question itself changes with the approach. Take ``why do users abandon fitness apps?''. Study A, quantitative and postpositivist, refines it to which factors predict abandonment, surveys 500 users, tests hypotheses about habit, price value and social influence, and reports odds ratios. Study B, qualitative and constructivist, refines it to how users experience abandoning an app, interviews twelve people who quit, and reports themes. Study C, mixed and pragmatist, runs the survey first and then interviews a sample of those who quit to explain the strongest predictor, in an explanatory sequential design. The refined question, the data, the analysis and the kind of knowledge produced differ in each case, and the essay must say which of the three it chooses and why.

Behind the surveys in this chapter stands one research programme. Venkatesh (2022) is a conceptual paper without data of its own: it identifies properties of AI tools, of the humans who work with them and of the operations settings that could hinder adoption, and uses UTAUT to propose research directions --- individual, technology and environmental characteristics, and interventions. A research agenda is a legitimate product of research, and reading it first shows where a study sits and where the gaps are. The ChatGPT studies above, and my own study of ChatGPT adoption among Big Four consultants, take their framework from this line of work.

\section{How to choose}

Three criteria decide, in this order. First, the research problem: new or unexplored concepts call for a qualitative approach, known factors to test or compare for a quantitative one, and both for mixed methods. Secondly, the audience: who must be convinced --- a clinic that trusts outcomes, a design team that trusts stories, or a regulator that wants both. Thirdly, your own experience and resources: statistics skills, interview skills, access to participants, and time. Recognising these criteria in a client's wording --- ``does it work, and how fast?'' is postpositivist, ``how do users experience it?'' is constructivist --- is this chapter applied to your own thesis.


\section{Worked examples and tables}

\subsection{The vocabulary --- one map, two sources}
Creswell and Oates use different words for the same decisions. The map translates between them, and the essay uses one vocabulary consistently.
\begin{center}\small
\small Creswell and Oates use \textbf{different words for similar layers}.
  This map holds for the whole course:
  \begin{center}
    \begin{tikzpicture}[node distance=0.32cm and 0.8cm]
      \node[kbox, text width=11.6cm] (l1)
        {\textbf{Research approach} \;{\scriptsize(Creswell)}\quad
         \emph{the overall plan: quantitative, qualitative or mixed methods}};
      \node[kbox, text width=11.6cm, below=of l1] (l2)
        {\textbf{Research design} \;{\scriptsize(Creswell)} \;=\;
         \textbf{research strategy} \;{\scriptsize(Oates)}\quad
         \emph{the type of study: survey, experiment, case study, ethnography, \dots}};
      \node[kbox, text width=11.6cm, below=of l2] (l3)
        {\textbf{Research methods}\quad
         \emph{the concrete techniques: questionnaire, interview, observation,
         statistical or thematic analysis}};
      \draw[karrow] (l1) -- node[right=2pt,font=\scriptsize\itshape,kgrey]{from broad\dots} (l2);
      \draw[karrow] (l2) -- node[right=2pt,font=\scriptsize\itshape,kgrey]{\dots to specific} (l3);
    \end{tikzpicture}
  \end{center}
  \begin{alertblock}{\small Common essay mistake}
    \small ``Our method is a case study.'' --- No: the case study is your
    \textbf{strategy/design}; interviews and document analysis are your \textbf{methods}.
  \end{alertblock}
\end{center}


\subsection{Not two boxes --- a continuum}
\begin{center}\small
\begin{center}
    \begin{tikzpicture}[scale=1.0]
      % axis
      \draw[very thick, kgrey, {Stealth}-{Stealth}] (-5.6,0) -- (5.6,0);
      % three anchors
      \node[kbox, minimum width=3.2cm] (quant) at (-4.2,1.1) {\textbf{Quantitative}\\numbers, variables,\\testing theories};
      \node[gbox, minimum width=3.2cm] (mixed) at (0,1.1) {\textbf{Mixed methods}\\integrating both\\kinds of data};
      \node[kbox, minimum width=3.2cm] (qual)  at (4.2,1.1) {\textbf{Qualitative}\\words, meanings,\\exploring settings};
      \draw[garrow] (quant) -- (-4.2,0.15);
      \draw[garrow] (mixed) -- (0,0.15);
      \draw[garrow] (qual)  -- (4.2,0.15);
      % labels below
      \node[font=\scriptsize\itshape, kgrey] at (-4.2,-0.45) {closed-ended data};
      \node[font=\scriptsize\itshape, kgrey] at (0,-0.45)    {both};
      \node[font=\scriptsize\itshape, kgrey] at (4.2,-0.45)  {open-ended data};
    \end{tikzpicture}
  \end{center}
  
  Creswell's point: qualitative and quantitative are the \textbf{ends of a
  continuum}. A study sits somewhere on it --- it tends toward one end more
  than the other. Mixed methods lives in the middle.
  

  {\color{kgrey}\small Source: Creswell \& Creswell (2018), Ch.\ 1}
\end{center}


\subsection{Deductive and inductive reasoning}
The engine underneath the approaches: deduction runs from theory to test, induction from data to theory, and the same study can use both in sequence.
\begin{center}\small
\begin{center}
    \resizebox{0.92\textwidth}{!}{%
    \begin{tikzpicture}[node distance=0.38cm and 1.4cm]
      % deductive column
      \node[kbox, text width=4.9cm] (t1)
        {\textbf{1.\ Theory}\\{\scriptsize\color{kgrey}e.g.\ self-determination theory:
          external rewards undermine intrinsic motivation}};
      \node[kbox, text width=4.9cm, below=of t1] (h1)
        {\textbf{2.\ Hypothesis}\\{\scriptsize\color{kgrey}Hypothesis 1 (H1): reward-driven users
          abandon the app more often}};
      \node[kbox, text width=4.9cm, below=of h1] (d1)
        {\textbf{3.\ Data collection}\\{\scriptsize\color{kgrey}questionnaire, 500 users,
          validated scales}};
      \node[kbox, text width=4.9cm, below=of d1] (c1)
        {\textbf{4.\ Confirm / reject}\\{\scriptsize\color{kgrey}statistics support H1 $\rightarrow$
          theory survives this test}};
      \draw[karrow] (t1) -- (h1);
      \draw[karrow] (h1) -- (d1);
      \draw[karrow] (d1) -- (c1);
      \node[above=0.22cm of t1, font=\bfseries] {Deductive = theory-\emph{testing}};
      % inductive column
      \node[gbox, text width=4.9cm, right=1.6cm of c1] (d2)
        {\textbf{1.\ Data collection}\\{\scriptsize\color{kgrey}12 open interviews with
          former users}};
      \node[gbox, text width=4.9cm, above=of d2] (p2)
        {\textbf{2.\ Patterns (codes)}\\{\scriptsize\color{kgrey}``felt judged'', ``lost my
          streak'', ``relief after deleting''}};
      \node[gbox, text width=4.9cm, above=of p2] (th2)
        {\textbf{3.\ Themes}\\{\scriptsize\color{kgrey}codes condense into
          ``quantified guilt''}};
      \node[gbox, text width=4.9cm, above=of th2] (g2)
        {\textbf{4.\ Theory / model}\\{\scriptsize\color{kgrey}a model of how self-tracking
          turns into pressure}};
      \draw[garrow] (d2) -- (p2);
      \draw[garrow] (p2) -- (th2);
      \draw[garrow] (th2) -- (g2);
      \node[above=0.22cm of g2, font=\bfseries] {Inductive = theory-\emph{building}};
    \end{tikzpicture}}
  \end{center}
  
  {\scriptsize Deduction runs \textbf{top-down} from theory (typical for
  quantitative research); induction runs \textbf{bottom-up} from data (typical
  for qualitative research). In practice research \textbf{cycles}: induction
  proposes a theory, deduction then tests it.}
\end{center}


\subsection{Two current quantitative studies}
\begin{center}\small
\footnotesize
  \renewcommand{\arraystretch}{1.2}
  \begin{tabular}{@{}p{2.2cm}p{5.2cm}p{5.8cm}@{}}
    \toprule
    & \textbf{Bilo\v{s} \& Budimir (2024)} & \textbf{Huy et al.\ (2024)}\\
    \midrule
    Sample & 694 Generation Z respondents, Croatia; online questionnaire, seven-point Likert & 671 ChatGPT users, Vietnam; questionnaire, face-to-face and online\\
    Theory & Extended UTAUT2 & UTAUT2 integrated with task--technology fit; trust and curiosity added\\
    Analysis & Confirmatory and exploratory factor analysis; hierarchical regression & Reliability, EFA/CFA, then covariance-based SEM (AMOS): use $\rightarrow$ continuance intention $\rightarrow$ word of mouth; curiosity as moderator\\
    Result & Performance expectancy, social influence, hedonic motivation, habit and personal innovativeness predict intention; effort expectancy, facilitating conditions and price value do not (65\,\% of variance explained) & Most UTAUT2 and TTF factors affect use; effort expectancy, social influence and trust do not; use drives continuance and recommendation\\
    \bottomrule
  \end{tabular}
  

  \small Both are deductive and postpositivist: theory $\rightarrow$
  hypotheses $\rightarrow$ measurement $\rightarrow$ test. Same theory family,
  two samples, partly different results --- replication in practice.
\end{center}


\subsection{The seven purposes of mixed methods research}
\begin{center}\small
\small Mixed methods research combines quantitative and qualitative methods
  in the same inquiry. The guidelines address three questions: (1) is a mixed
  methods approach \textbf{appropriate}; (2) how are \textbf{meta-inferences}
  --- the conclusions that integrate both strands --- developed; (3) how is
  their quality \textbf{validated}.
  
  \footnotesize
  \renewcommand{\arraystretch}{1.15}
  \begin{tabular}{@{}p{3.2cm}p{9.8cm}@{}}
    \toprule
    \textbf{Purpose} & \textbf{Why combine the two strands}\\
    \midrule
    Complementarity & To gain complementary views of the same phenomenon or relationship\\
    Completeness & To obtain a complete picture of a phenomenon\\
    Developmental & Questions or hypotheses for one strand emerge from the inferences of the previous one\\
    Expansion & To explain or expand the understanding obtained in a previous strand\\
    Corroboration / confirmation & To assess the credibility of inferences obtained from one strand\\
    Compensation & To compensate for the weaknesses of one approach with the other\\
    Diversity & To obtain divergent views of the same phenomenon\\
    \bottomrule
  \end{tabular}
  

  {\color{uiagrey}\scriptsize Purposes adapted from Venkatesh et al.\ (2013), Table 1. An essay that proposes mixed methods names its purpose in these terms.}
\end{center}


\subsection{Designing a mixed-methods study}
\begin{center}\small
\small The 2016 extension turns the guidelines into design decisions: 14
  properties of a mixed-methods study, each with a question to answer and a
  decision tree that connects them. The ones that decide most:
  
  \begin{columns}[T]
    \column{0.5\textwidth}
    \textbf{Foundations}
    \begin{itemize}\setlength{\itemsep}{0.15em}
      \item \textbf{Research questions} --- questions, aims or hypotheses; independent or dependent
      \item \textbf{Purpose} --- one of the seven on the previous slide
      \item \textbf{Paradigm stance} --- single or multiple; pragmatism, critical realism, dialectical
    \end{itemize}
    \column{0.48\textwidth}
    \textbf{Design strategies}
    \begin{itemize}\setlength{\itemsep}{0.15em}
      \item \textbf{Exploratory or confirmatory} --- develop or test theory
      \item \textbf{Strands and phases} --- one study or several
      \item \textbf{Mixing and timing} --- fully or partially mixed; concurrent or sequential
      \item \textbf{Inference quality} --- how the meta-inferences are validated
    \end{itemize}
  \end{columns}
  
  {\color{uiagrey}\small For your essay: a mixed-methods proposal states the
  purpose, the sequence of strands, which strand has priority, and where the
  two are integrated. Without these four decisions it is two studies, not one.}
\end{center}


\subsection{Three mini-studies to classify}
Three short descriptions to classify along the way --- quantitative, qualitative or mixed, and why. The answers are given below the descriptions.
\begin{center}\small
\small
  \begin{enumerate}\setlength{\itemsep}{0.7em}
    \item A researcher analyses \textbf{six months of telemetry logs} from
          50{,}000 app installs and models which onboarding events predict churn.
    \item A researcher spends \textbf{three weeks in a security operations
          centre (SOC)} observing how analysts handle the flood of automated
          alerts, then interviews eight of them.
    \item A researcher runs a \textbf{questionnaire} (n = 412) testing whether
          perceived usefulness predicts students' use of artificial intelligence (AI) chatbots, then
          \textbf{interviews ten students} to interpret an unexpected result.
  \end{enumerate}
  
  {\color{kgrey}\small Answers: (1) quantitative --- measured behaviour,
  deductive modelling; (2) qualitative --- meanings in a real setting,
  inductive; (3) mixed methods --- explanatory sequential. In the discussion at the end you
  will make such calls yourselves, with justification.}
\end{center}


\subsection{Creswell's framework and the four worldviews}
\begin{center}\small
\begin{center}
    \begin{tikzpicture}[node distance=0.9cm and 1.1cm]
      \node[kbox, minimum width=3.0cm, minimum height=1.5cm] (wv)  {\textbf{Philosophical}\\\textbf{worldview}\\ \scriptsize postpositivist, constructivist,\\ \scriptsize transformative, pragmatist};
      \node[kbox, minimum width=3.0cm, minimum height=1.5cm, right=of wv] (de) {\textbf{Research}\\\textbf{design}\\ \scriptsize survey, experiment,\\ \scriptsize case study, ethnography, \dots};
      \node[kbox, minimum width=3.0cm, minimum height=1.5cm, right=of de] (me) {\textbf{Research}\\\textbf{methods}\\ \scriptsize questions, data collection,\\ \scriptsize analysis, interpretation};
      \node[gbox, minimum width=7.6cm, minimum height=1.1cm, below=1.1cm of de]
        (ap) {\textbf{Research approach}\\ \scriptsize qualitative $\cdot$ quantitative $\cdot$ mixed methods};
      \draw[karrow] (wv) -- (de);
      \draw[karrow] (de) -- (me);
      \draw[garrow] (wv.south) to[out=-70,in=160] (ap.west);
      \draw[garrow] (de.south) -- (ap.north);
      \draw[garrow] (me.south) to[out=-110,in=20] (ap.east);
    \end{tikzpicture}
  \end{center}
  
  Approach $=$ worldview $+$ design $+$ methods, made consistent.
  \hfill {\color{kgrey}\small Adapted from Creswell \& Creswell (2018), Fig.\ 1.1}
\end{center}

\begin{center}\small
\small
  \renewcommand{\arraystretch}{1.35}
  \begin{tabular}{@{}p{2.6cm}p{4.6cm}p{5.6cm}@{}}
    \toprule
    \textbf{Worldview} & \textbf{Knowledge is\dots} & \textbf{Typical IT/health example}\\
    \midrule
    \textcolor{kristiania}{\textbf{Postpositivism}} & objective; found by measuring and testing causes & A/B test: which onboarding flow reduces app churn?\\
    \textcolor{kristiania}{\textbf{Constructivism}} & constructed; multiple meanings in social settings & Interviews: what does ``secure enough'' \emph{mean} to small-business sysadmins?\\
    \textcolor{kristiania}{\textbf{Transformative}} & tied to power and change; research should empower & Co-designing an accessible web frontend \emph{with} visually impaired users\\
    \textcolor{kristiania}{\textbf{Pragmatism}} & whatever works to solve the problem & Mixed-methods evaluation of a password manager\\
    \bottomrule
  \end{tabular}
  

  {\color{kgrey}\small Based on Creswell \& Creswell (2018), Ch.\ 1}
\end{center}


\subsection{The layers of the onion}
\begin{center}\small
\centering
  \footnotesize
  \begin{tabular}{@{}L{2.4cm}L{3.9cm}L{4.6cm}@{}}
    \toprule
    \textbf{Position} & \textbf{Ontology} & \textbf{Typical consequence} \\
    \midrule
    Positivism & One observer-independent reality & Hypothesis testing, large samples \\
    Critical realism & Real structures, observable only indirectly & Search for generative mechanisms \\
    Interpretivism & Socially constructed realities & In-depth interviews, reconstruction of meaning \\
    Postmodernism & Reality as an effect of language and power & Deconstruction of dominant narratives \\
    Pragmatism & Reality judged by problem solving & Method follows the research question \\
    \bottomrule
  \end{tabular}
  \par
  \raggedright
  \src{Every subsequent layer depends on this choice. Oates (Ch.\ 19--20) covers positivism, interpretivism and critical research; Saunders adds critical realism, postmodernism and pragmatism.}
  \onionloc{1}{0}{0}{0}{0}{0}{Layer 1}
\end{center}

\begin{center}\small
\begin{minipage}{0.85\textwidth}
  {\color{uiared}\bfseries Approach to theory development}
  \begin{itemize}\footnotesize\setlength{\itemsep}{0.25em}
    \item \textbf{Deduction} --- theory first, then test (typical of UTAUT and TAM studies)
    \item \textbf{Induction} --- data first, theory as the outcome
    \item \textbf{Abduction} --- moving between the two; the starting point is a surprising observation
  \end{itemize}
  
  {\color{uiared}\bfseries Methodological choice}
  \begin{itemize}\footnotesize\setlength{\itemsep}{0.25em}
    \item \textbf{Mono-method} --- a single type of data, quantitative or qualitative
    \item \textbf{Multi-method} --- several types of data within one paradigm
    \item \textbf{Mixed-method} --- quantitative and qualitative combined
      \begin{itemize}\scriptsize
        \item \emph{simple}: one combination, usually sequential
        \item \emph{complex}: repeated combination across several phases
      \end{itemize}
  \end{itemize}
  \end{minipage}
  \onionloc{0}{1}{1}{0}{0}{0}{Layers 2--3}
\end{center}

\begin{center}\small
\begin{columns}[T,onlytextwidth]
    \begin{column}{0.42\textwidth}
      {\color{uiared}\bfseries Research strategies}
      \begin{itemize}\footnotesize\setlength{\itemsep}{0.15em}
        \item Experiment
        \item Survey
        \item Archival research
        \item Case study
        \item Ethnography
        \item Action research
        \item Grounded theory
        \item Narrative inquiry
      \end{itemize}
    \end{column}
    \begin{column}{0.44\textwidth}
      {\color{uiared}\bfseries Time horizon}
      \begin{itemize}\footnotesize\setlength{\itemsep}{0.15em}
        \item Cross-sectional --- a snapshot
        \item Longitudinal --- change over time
      \end{itemize}
      
      {\color{uiared}\bfseries Techniques and procedures}
      \begin{itemize}\footnotesize\setlength{\itemsep}{0.15em}
        \item Collection: questionnaire, interview, observation, secondary data
        \item Analysis: descriptive and inferential statistics, SEM, thematic analysis, coding
      \end{itemize}
    \end{column}
    \begin{column}{0.13\textwidth}\ \end{column}
  \end{columns}
  
  \src{Sessions 6--7 take the strategies one by one; sessions 8--9 the techniques.}
  \onionloc{0}{0}{0}{1}{1}{1}{Layers 4--6}
\end{center}


\subsection{One question, three studies --- side by side}
The three studies of Section 2.4 in one table: the refined question, the data, the analysis and the kind of knowledge each produces.
\begin{center}\small
\small
  \begin{columns}[T]
    \column{0.55\textwidth}
    \textbf{Refined question:} \emph{Which factors predict abandonment of
    fitness apps within 90 days?}
    \begin{itemize}\setlength{\itemsep}{0.3em}
      \item \textbf{Theory:} self-determination theory $\rightarrow$ hypotheses
            (hypothesis 1 (H1): external rewards $\rightarrow$ higher churn, \dots)
      \item \textbf{Data:} online questionnaire, n = 500 current/former users;
            validated scales
      \item \textbf{Analysis:} logistic regression
    \end{itemize}
    \column{0.42\textwidth}
    \begin{block}{What you get}
      Generalisable pattern: e.g.\ goal-setting halves the odds of churn.
    \end{block}
    \begin{alertblock}{What you miss}
      \emph{Why} goal-setting works, and everything you didn't
      think to ask about.
    \end{alertblock}
  \end{columns}
\end{center}

\begin{center}\small
\small
  \begin{columns}[T]
    \column{0.55\textwidth}
    \textbf{Refined question:} \emph{How do former users describe the process
    of abandoning their fitness app?}
    \begin{itemize}\setlength{\itemsep}{0.3em}
      \item \textbf{Data:} 12 semi-structured interviews, 45--60 min,
            participants recruited until themes repeat (saturation)
      \item \textbf{Analysis:} thematic analysis, inductive coding
      \item \textbf{Result:} themes --- streak anxiety, life disruptions,
            ``quantified guilt'', quiet relief after quitting
    \end{itemize}
    \column{0.42\textwidth}
    \begin{block}{What you get}
      Deep understanding of the \emph{experience}; concepts nobody
      had put in a questionnaire yet.
    \end{block}
    \begin{alertblock}{What you miss}
      How common each theme is; you cannot generalise to ``users in general''.
    \end{alertblock}
  \end{columns}
\end{center}

\begin{center}\small
\small
  \textbf{Refined question:} \emph{Which factors predict abandonment --- and how
  do users explain the strongest ones?}
  \begin{center}
    \begin{tikzpicture}[node distance=0.5cm and 1.0cm]
      \node[kbox, minimum width=3.6cm, minimum height=1.3cm] (q1)
        {\textbf{QUANT}\\survey, n = 500\\regression};
      \node[gbox, minimum width=3.6cm, minimum height=1.3cm, right=of q1] (q2)
        {\textbf{qual}\\10 follow-up interviews\\on surprising results};
      \node[kbox, minimum width=3.6cm, minimum height=1.3cm, right=of q2] (int)
        {\textbf{Integration}\\joint interpretation\\\& recommendations};
      \draw[karrow] (q1) -- node[above,font=\scriptsize]{informs} (q2);
      \draw[karrow] (q2) -- (int);
    \end{tikzpicture}
  \end{center}
  
  This is an \textbf{explanatory sequential design}: quantitative first,
  qualitative second to \emph{explain} the numbers. Note the deliberate order
  and integration point --- that is what makes it a design, not a coincidence.
  \emph{Notation:} CAPITAL letters (QUANT) mark the \textbf{dominant} component,
  lower-case (qual) the supporting one --- a standard convention in mixed
  methods research.
\end{center}

\begin{center}\small
\small
  \renewcommand{\arraystretch}{1.3}
  \begin{tabular}{@{}p{2.4cm}p{3.4cm}p{3.4cm}p{3.5cm}@{}}
    \toprule
     & \textbf{Study A (quant)} & \textbf{Study B (qual)} & \textbf{Study C (mixed)}\\
    \midrule
    Worldview   & postpositivist & constructivist & pragmatist\\
    Question    & \emph{which factors predict\dots} & \emph{how do users experience\dots} & both, in sequence\\
    Sample      & 500, random & 12, purposive & 500 + 10\\
    Claim type  & generalisable pattern & rich understanding & pattern + explanation\\
    Weakness    & no ``why'' & no ``how many'' & cost, complexity\\
    \bottomrule
  \end{tabular}
  

  \begin{alertblock}{The essay question behind every essay}
    Not ``which approach is best?'' --- but
    \textbf{``which approach fits \emph{your} question, and can you argue it?''}
  \end{alertblock}
\end{center}


\subsection{A study from my own research}
\begin{center}\small
\begin{block}{ChatGPT adoption in Big Four consulting firms (Amling, ongoing)}
    \textbf{Question:} why do consultants adopt --- or quietly avoid ---
    ChatGPT at work?\\[0.3em]
    \textbf{Design:} mixed methods. A survey based on the UTAUT2 technology-%
    acceptance model, extended with \emph{confidence}, analysed with
    structural equation modelling (quantitative) --- combined with qualitative
    interviews, with data from Germany and Norway.
  \end{block}
  
  Why show you this? Because the pattern you just saw --- numbers for the
  \emph{what}, interviews for the \emph{why} --- is not a textbook idealisation.
  It is how technology-adoption research is actually done, including in your
  lecturer's ongoing work.
\end{center}


\subsection{How to choose}
\begin{center}\small
\begin{enumerate}\setlength{\itemsep}{0.6em}
    \item \textbf{The research problem.}\\
          New, unexplored concept $\rightarrow$ qualitative.
          Known factors to test or compare $\rightarrow$ quantitative.
          Numbers that need explaining $\rightarrow$ mixed.
    \item \textbf{Your audience.}\\
          Who must be convinced --- a clinic that thinks in outcomes,
          or a design team that thinks in user stories?
    \item \textbf{Your own experience and resources.}\\
          Statistics skills? Interview skills? Access to participants? Time?
  \end{enumerate}
  
  {\color{kgrey}\small Based on Creswell \& Creswell (2018), Ch.\ 1}
\end{center}


\section*{Terms from this session}
\begin{center}\small
\footnotesize
  \renewcommand{\arraystretch}{1.12}
  \begin{tabular}{@{}p{3.2cm}p{10.3cm}@{}}
    \toprule
    \textbf{Research approach} & overall plan of a study: quantitative, qualitative or mixed methods\\
    \textbf{Research design / strategy} & the type of study (survey, experiment, case study, \dots); Creswell: \emph{design}, Oates: \emph{strategy}\\
    \textbf{Research method} & concrete technique for collecting or analysing data (interview, questionnaire, \dots)\\
    \textbf{Worldview / paradigm} & set of beliefs about reality and knowledge guiding the research\\
    \textbf{Ontology} & assumptions about what reality is\\
    \textbf{Epistemology} & assumptions about how we can know reality\\
    \textbf{Postpositivism} & one imperfectly knowable reality; test and refine theories by measurement\\
    \textbf{Constructivism} & multiple constructed realities; understand meanings in context\\
    \textbf{Deductive / inductive} & theory $\rightarrow$ data (testing) vs.\ data $\rightarrow$ theory (exploring)\\
    \textbf{Triangulation} & examining one phenomenon from several angles (methods, data, sources) to corroborate findings\\
    \textbf{Explanatory sequential} & mixed design: quantitative first, qualitative second to explain the numbers\\
    \bottomrule
  \end{tabular}
  

  {\color{kgrey}\scriptsize Keep this slide next to you while writing Part III of the essay.}
\end{center}



\section*{Questions from the reading}

  \begin{enumerate}\setlength{\itemsep}{0.4em}
    \item On which two assumptions does \textbf{positivism} rest, according to Oates?
          \answer{That the world exists independently of us and is ordered and regular (ontology), and that we can discover it through observation and measurement --- the scientific method (epistemology).}
    \item What does \textbf{interpretivism} set out to produce instead of universal laws?
          \answer{An understanding of the social context: how people make sense of their world, with multiple subjective realities and the researcher as part of the situation studied.}
    \item What distinguishes \textbf{critical research} from interpretive research?
          \answer{It goes beyond understanding: it identifies and challenges power relations, conflicts and contradictions, with the aim of emancipation and change.}
    \item Which strategies does Oates link to which paradigm --- and which strategy fits any of them?
          \answer{Experiments and surveys lean positivist; case studies, action research and ethnography lean interpretive or critical; design and creation can be carried out under any paradigm.}
  \end{enumerate}

\section*{Questions for discussion}

  \begin{enumerate}\setlength{\itemsep}{0.25em}
    \item \textbf{Your topic} --- write your question twice: once in a quantitative form (\emph{which factors, how many, does X affect Y}) and once in a qualitative form (\emph{how do, what does X mean to}). Which version is more interesting, and which is more feasible?
          \answer{Most topics allow both. The quantitative form needs measurable variables and enough cases; the qualitative form needs access to people and time for analysis. Decide with today's three criteria: fit to the question, fit to your resources, and the audience of the result.}
    \item \textbf{The telemetry study} --- a company analyses usage logs of 40{,}000 users and then interviews twelve of them. Which approach, and which worldview?
          \answer{Mixed methods, explanatory sequential: the quantitative phase finds the pattern, the qualitative phase explains it. The worldview is pragmatist: the question decides the method.}
    \item \textbf{Your client's question} --- ``does it work, and how fast?'' versus ``how do users experience it?'' Which worldview does each carry?
          \answer{The first is postpositivist: an objective outcome, measured and compared. The second is constructivist: meaning as experienced by users, reconstructed from their accounts. Recognising this in a client's wording is today's content applied to your own thesis.}
  \end{enumerate}

\section*{Key points}
\begin{itemize}
  \item Three approaches: quantitative, qualitative, mixed --- a continuum, not a binary
  \item Every study rests on a worldview --- and the research onion orders all six decisions, from philosophy to technique
  \item Mixed methods: name the purpose (Venkatesh et al., 2013) and the design decisions (2016)
  \item Two current ChatGPT studies, one research agenda --- and the justification your essay needs (Part III)
\end{itemize}

\section*{Sources}
\begin{itemize}\small
  \item Bilo\v{s}, A., \& Budimir, B. (2024). Understanding the adoption dynamics of ChatGPT among Generation Z: Insights from a modified UTAUT2 model. \emph{Journal of Theoretical and Applied Electronic Commerce Research}, \emph{19}(2), 863--879. \url{https://doi.org/10.3390/jtaer19020045}
  \item Huy, L.\,V., Nguyen, H.\,T.\,T., Vo-Thanh, T., Thinh, N.\,H.\,T., \& Tran, T.\,T.\,D. (2024). Generative AI, why, how, and outcomes: A user adoption study. \emph{AIS Transactions on Human-Computer Interaction}, \emph{16}(1), 1--27. \url{https://doi.org/10.17705/1thci.00198}
  \item Presthus, W., \& Munkvold, B.\,E. (2016). How to frame your contribution to knowledge? \emph{NOKOBIT}, Bergen.
  \item Venkatesh, V. (2022). Adoption and use of AI tools: A research agenda grounded in UTAUT. \emph{Annals of Operations Research}, \emph{308}(1), 641--652. \url{https://doi.org/10.1007/s10479-020-03918-9}
  \item Venkatesh, V., Brown, S.\,A., \& Bala, H. (2013). Bridging the qualitative--quantitative divide: Guidelines for conducting mixed methods research in information systems. \emph{MIS Quarterly}, \emph{37}(1), 21--54. \url{https://doi.org/10.25300/MISQ/2013/37.1.02}
  \item Venkatesh, V., Brown, S.\,A., \& Sullivan, Y.\,W. (2016). Guidelines for conducting mixed-methods research: An extension and illustration. \emph{Journal of the Association for Information Systems}, \emph{17}(7), 435--494. \url{https://doi.org/10.17705/1jais.00433}
  \item Mayer, A.-S., Baygi, R.\,M., \& Buwalda, R. (2025). Generation AI: Job crafting by entry-level professionals in the age of generative AI. \emph{Business \& Information Systems Engineering}, \emph{67}(5), 595--613. \url{https://doi.org/10.1007/s12599-025-00959-x}
  \item Saunders, M., Lewis, P., \& Thornhill, A. (2019). \emph{Research methods for business students} (8th ed.). Pearson.
  \item Oates, B.\,J., Griffiths, M., \& McLean, R. (2022). \emph{Researching information systems and computing} (2nd ed.). SAGE. --- Chapters 19 (Philosophical paradigms: positivism) and 20 (Alternative philosophical paradigms).
  \item Laudon, K.\,C., \& Laudon, J.\,P. (2022). \emph{Management information systems: Managing the digital firm} (17th ed., Global ed.). Pearson. --- Chapter 1, section 1-3 (technical and behavioural approaches).
\end{itemize}

\chapter{Purpose and Products of Research}
\label{ch:3}
\sessionhead{3}{Thursday 24 September 2026 \;$\cdot$\; reading: Oates, Chapters 1--3}

Why do people do research, and what does research produce? This chapter answers both questions with Oates' 6Ps --- purpose, products, process, participants, paradigm and presentation --- and walks one published study through all six. It then turns to the products: the forms a contribution to knowledge can take, from a taxonomy to a confirmed model, and the difference between descriptive and prescriptive knowledge. The last section contrasts research with the other system that produces the knowledge organisations actually use, standards and frameworks, and draws five rules for the method part of the essay.


\section{The 6Ps}

Oates organises the whole of research around six aspects that every study has, whether or not the author has thought about them: purpose, products, process, participants, paradigm and presentation. The purpose asks why the study is done; the products ask what comes out of it; the process describes the sequence of activities; the participants are the people involved and the way they are treated; the paradigm is the set of assumptions behind the study; and the presentation is the way in which the study is reported. The six are also the skeleton of the essay: purpose and questions belong to Part I, the process and the paradigm to Part III, the participants to Part IV, and the presentation is the essay itself, about ten pages and referenced.

A published study makes the six concrete. Huy et al.\ (2024) asked which factors drive ChatGPT use and whether use leads to continued use and recommendation; the gap was that UTAUT2 and task--technology fit had not been combined for generative AI and that curiosity was untested. Their products were an integrated model, a six-item scale for ChatGPT use, and evidence that most UTAUT2 and TTF factors matter while effort expectancy, social influence and trust do not. The process was deductive: literature, then model and hypotheses, then a questionnaire with five-point Likert items, back-translated and piloted, then exploratory and confirmatory factor analysis and covariance-based structural equation modelling. The participants were 671 ChatGPT users in Vietnam, a convenience sample recruited face to face in four cities and online, supplemented by twenty interviews and five AI professors who helped to develop the use scale. The paradigm was positivist --- constructs measured, hypotheses tested, one model for all respondents --- and the presentation is a journal article with introduction, theory and hypotheses, methodology, results and discussion. Note the twenty interviews inside a quantitative study: a developmental mixed-methods element, used to build the instrument rather than to answer the question.

\section{Purpose}

\begin{block}{Purpose}
The purpose of a study refers to its reason for doing research: the topic of interest, why it is important or useful to study, and the specific research question or objectives that follow.
\end{block}

Oates' reasons for doing research translate readily into the security field. To add to the body of knowledge: how do phishing tactics evolve over time? To solve a problem: reduce false positives in a security operations centre. To test or evaluate a method: does an API gateway actually reduce attack surface? To find out ``what happens if'': what if we push passkeys to all customers? To develop a model or tool: build a threat model for a hospital network. To contribute to personal development: it counts, but write the essay's motivation on one of the others. Most cyber security proposals pursue the second or third reason; the essay's purpose statement is the second half of the funnel that turns a topic into a problem and a problem into a purpose.

Research for practice has a long history in our field. McAfee and Brynjolfsson (2008) argued in the \emph{Harvard Business Review} that competition in IT-intensive industries had accelerated since the mid-1990s --- winners pull away faster and leaders change more often --- because enterprise IT lets a firm replicate a better process across the whole organisation at once. The claim rested on the authors' own analyses of US industry data, published separately in peer-reviewed form and summarised for managers in the article. The same mechanism --- a technology that propagates good practice fast --- is the argument made for AI today, which makes the 2008 article a template for reading 2026 claims. An HBR article sits in the professional layer of the supply chain; it is useful precisely because it rests on research the reader can trace, and the trace should be checked before it is cited.

\subsection{From topic to problem to purpose}

The funnel narrows in three steps. A topic is a broad area of interest, such as ``password security''. A problem is something specific that is wrong, unknown or contested within the topic --- ``users reuse passwords despite training''. A purpose statement says what the study will do about the problem and to what end --- ``to identify the situational factors that predict reuse, so that training can target them''. A purpose statement that needs two ``and''s is still too wide; it names one problem, one owner and one intended use of the result, and it is Part I of the essay in miniature. Five drafts of the same statement, from ``to look into cyber security in companies'' to ``to identify which password-reuse triggers a security awareness programme should address first, using interviews with 20 employees of one Norwegian SME'', show the movement from a topic to a purpose with a population, a method and a use.

\section{Products}

\begin{block}{Products}
The products of research refer to its outcomes, especially the contribution to knowledge: an answer to the research questions, but also unexpected findings.
\end{block}

Oates lists the products research can produce: a new or improved product, a new theory, a re-interpretation of an existing theory, an in-depth analysis of a situation, an exploration of a topic, a critical analysis, a new research tool or technique, and a new or improved model. In IT terms these are an intrusion-detection component, a theory of why administrators ignore alerts, a critique of a maturity model, a security-awareness questionnaire, and so on. The contribution is what the world knows or has after the study that it did not know or have before.

Hevner (2021) adds a distinction that clarifies most master's theses. Science is both investigation (the process) and knowledge (the outcome), and the knowledge a study adds falls into one of two bases. Descriptive knowledge --- the ``what'' --- comprises phenomena and the regularities among them: observations, classifications, measurements, laws, patterns and theories. Prescriptive knowledge --- the ``how'' --- comprises human-built artefacts: constructs, models, methods and instantiations, together with design theories that state principles of form and function. A study shows its rigour and relevance by how it consumes existing knowledge and what it adds to these two bases. For your essay the consequence is one sentence: say which kind of knowledge your thesis will add --- a finding about how things are, or an artefact and the knowledge of how to build it. Most cyber security theses produce the second, and should say so.

Presthus and Munkvold (2016), writing at NOKOBIT for master's students and junior researchers, offer a taxonomy of the forms a contribution can take. Contributions to practice include lessons learned, experience reports, guidelines and roadmaps, heuristics, critical success factors and patterns. Theoretical contributions rise from building blocks upwards: concept, construct, rich insight and case description; framework or taxonomy, method and proposition; generative mechanism, hypothesis and model; and mid-range theory, design theory and grand theory. Four ways to contribute to an existing theory cover what most master's theses actually do: confirmation or replication, extension, contradiction and elimination. The authors consider this sufficient, provided the contribution is stated precisely, and give three pieces of advice: balance your own idea against what is already known; be critical of what you read and ask whether the field needs one more study on the topic; and decide the intended contribution before the study starts.

\subsection{What is expected at master's level}

The essay guideline is explicit on three points. You are not expected to answer your research questions in the essay --- you are expected to justify how you would. Your topic should be relevant, well understood and well delimited: context matters, and so does feasibility. A good essay states its contribution to practice and, where applicable, to academia, and argues why the study is worth doing. Nobody expects a new attack class; ``we evaluate passkey onboarding in a real organisation, with a justified method'' is an excellent master's-level topic.

\section{Two systems of knowledge production}

A large share of the knowledge applied in security organisations does not come from research at all. Research makes descriptive and explanatory claims, asks what is the case and why, tests them through methodological transparency, replicability and peer review, and produces theories, models and effect sizes. Practice and standardisation make normative and prescriptive claims, ask what should be done, test them through expert consensus, committee votes and use in the field, and produce standards, frameworks and guidelines. Both produce knowledge, but they answer different questions and justify themselves in different ways.

The practice-based sources are familiar. Formal standardisation --- ISO/IEC 27001, ISO 31000 --- draws its legitimacy from the consensus of national standards bodies; government frameworks such as the NIST Cybersecurity Framework, NSM's grunnprinsipper and BSI IT-Grundschutz from a state mandate; professional bodies' frameworks such as COBIT from the profession and its certification; and consulting knowledge --- maturity models and frameworks --- from market adoption and references. These sources are normatively binding without having been tested empirically, and students cite them as if they were literature.

The ISO process shows why. A proposal becomes a working draft, a committee draft, an enquiry draft, a final draft and a published standard; the committee stage iterates until consensus, a failed ballot returns the draft, and a systematic review every five years confirms, revises or withdraws the document. The process does loop, but the loops are about consensus, not evidence. A standard holds because the procedure was followed; nothing in it tests whether the measures work.

\begin{center}\small
\begin{center}
  \begin{tikzpicture}[
    stage/.style={draw=uiared, fill=uiared!8, rounded corners=2pt,
                  text width=1.72cm, align=center, font=\scriptsize,
                  minimum height=0.88cm, text=ink},
    fwd/.style={-{Stealth[length=2.2mm]}, uiared, thick},
    back/.style={-{Stealth[length=1.8mm]}, black!60, line width=0.55pt},
    lbl/.style={font=\fontsize{6pt}{7pt}\selectfont, text=black!65, align=center}]
    \node[stage] (a) at (-5.20,0) {Proposal\\(NWIP)};
    \node[stage] (b) at (-3.12,0) {Working draft\\(WD)};
    \node[stage] (c) at (-1.04,0) {Committee draft\\(CD)};
    \node[stage] (d) at ( 1.04,0) {Enquiry\\(DIS)};
    \node[stage] (e) at ( 3.12,0) {Approval\\(FDIS)};
    \node[stage] (f) at ( 5.20,0) {Publication\\(IS)};
    \foreach \x/\y in {a/b,b/c,c/d,d/e,e/f} \draw[fwd] (\x) -- (\y);
    \draw[back] ([xshift=-5mm]c.north) to[out=105,in=75,looseness=3.4] ([xshift=5mm]c.north);
    \node[lbl, anchor=east] at (-1.85,1.05) {iterates until consensus};
    \draw[back] (d.north) to[out=105,in=50,looseness=0.75] (c.north east);
    \draw[back] (e.north) to[out=105,in=42,looseness=0.60] ([xshift=2mm]c.north east);
    \node[lbl] at (3.05,1.36) {returned if the ballot fails};
    \node[draw=uiared, fill=white, rounded corners=2pt, text=ink,
          text width=2.35cm, align=center, minimum height=0.62cm,
          font=\fontsize{6pt}{7pt}\selectfont] (r) at (5.20,-1.20) {Systematic review\\(every five years)};
    \draw[back] (f.south) -- (r.north);
    \draw[back] (r.west) -- (-5.20,-1.20) -- (a.south);
    \node[lbl, fill=white, inner sep=1.5pt] at (-0.40,-1.20) {confirmed, revised or withdrawn};
  \end{tikzpicture}
  \end{center}
  
  \begin{itemize}\footnotesize\setlength{\itemsep}{0.2em}
    \item The process does loop --- but the loops are about \textbf{consensus}, not about evidence.
    \item A standard holds because the procedure was followed. Nothing in it tests whether the measures work.
  \end{itemize}
\end{center}


Seven dimensions summarise the comparison. Research makes descriptive claims, validates them by evidence, replication and peer review, corrects errors by refutation, moves slowly, aims at generalisation, addresses the research community, and fails when a hypothesis is rejected. Standards make prescriptive claims, validate them by consensus, voting and uptake, correct errors through the document's revision cycle, move fast and market-driven, aim at applicability, address the organisation, the auditor and the regulator, and fail when the standard is not adopted. The error-correction row is the most interesting: a standard gets revised, but never refuted.

The two systems are connected in four ways. First, practice supplies the phenomena: RPA, agile, DevOps, zero trust, cloud governance and agentic AI all appeared in practice first, and research responds with a lag of two to four years against product cycles of months; some of the constructs research works with are taken from vendor material, and must be examined before operationalisation, or marketing vocabulary becomes a measurement. Secondly, research tests after the fact: does the prescribed measure work, under which conditions, with what effect size and side effects, and is the effect causal or self-selection? Standards assert effectiveness implicitly, through the authority of the procedure, and empirical evidence of effectiveness is neither required nor supplied. Thirdly, standards are themselves an object of research: Meyer and Rowan (1977) explain that structures are adopted for legitimacy and become decoupled from lived practice, and DiMaggio and Powell (1983) explain the spread of standards through coercive, mimetic and normative isomorphism --- which is why standards spread without evidence of effectiveness. Lastly, research feeds back into practice through Mode 2 knowledge production (Gibbons et al., 1994), design science research (Hevner et al., 2004; Peffers et al., 2007), engaged scholarship (Van de Ven \& Johnson, 2006) and evidence-based management (Rousseau, 2006).

Five rules follow for the method part of the essay. Establish where the constructs come from --- theory or a vendor paper. Separate the normative from the descriptive: a standard is data, not a literature base. Treat standards as citable but not as evidence: they show what is required, not what works. Use the onion as a checklist and justify each layer explicitly, especially the step from philosophy to method. Lastly, proximity to practice is access to data, not an argument: relevance does not substitute for methodological rigour.

\section{The research process}

Oates' model of the research process runs from experiences and motivation through a literature review to research questions, a conceptual framework, one or more strategies, data generation methods, data analysis and the presentation of the results. The model is a map rather than a schedule: real projects loop back, and the literature review in particular is revisited when the question changes. The sessions of this course follow the model in order, and the essay's four parts do the same.


\section{Worked examples and tables}

\subsection{The 6Ps as the essay's skeleton}
\begin{center}\small
\small
  \renewcommand{\arraystretch}{1.35}
  \begin{tabular}{@{}p{2.6cm}p{10.6cm}@{}}
    \toprule
    \textbf{P} & \textbf{Where it appears in your essay}\\
    \midrule
    Purpose       & Part I: problem, aim, objective, research question\\
    Products      & Part I \& IV: expected contribution to practice/academia\\
    Process       & Part II \& III: literature plan, strategy, methods, data plan\\
    Participants  & Part IV: ethics, SIKT approval, who provides data\\
    Paradigm      & Part III: the (implicit) worldview behind your choices\\
    Presentation  & the essay document itself: 10 pages, referenced\\
    \bottomrule
  \end{tabular}
  

  If you can fill this table for your own project, your essay is
  essentially outlined.
\end{center}


\subsection{Oates' reasons, translated to IT}
\begin{center}\small
\small
  \renewcommand{\arraystretch}{1.3}
  \begin{tabular}{@{}p{4.3cm}p{9.0cm}@{}}
    \toprule
    \textbf{Reason} & \textbf{IT / security example}\\
    \midrule
    Add to the body of knowledge & how do phishing tactics evolve over time?\\
    Solve a problem              & reduce alert fatigue in a security team\\
    Test or evaluate something   & does our new API gateway actually scale?\\
    Find out ``what happens if'' & what if we drop passwords for passkeys?\\
    Personal development         & learn threat modelling by studying it properly\\
    Qualification (you!)         & pass this course; write your master's thesis\\
    \bottomrule
  \end{tabular}
  

  Most real projects mix several reasons --- but \textbf{one} should dominate,
  because it decides what a ``good answer'' looks like.
  \hfill {\color{kgrey}\small Based on Oates et al.\ (2022), Ch.\ 2}
\end{center}


\subsection{The funnel}
\begin{center}\small
\begin{center}
    \resizebox{0.92\textwidth}{!}{%
    \begin{tikzpicture}[node distance=0.32cm]
      \node[gbox, text width=11.8cm] (t)
        {\textbf{Topic} (broad area of interest)\\
         {\scriptsize\color{kgrey}``password security''}};
      \node[kbox, text width=9.6cm, below=of t] (pr)
        {\textbf{Problem} (something specific that is wrong, unknown or needed)\\
         {\scriptsize\color{kgrey}``employees reuse passwords across systems ---
          despite mandatory awareness training''}};
      \node[kbox, text width=7.4cm, below=of pr] (pu)
        {\textbf{Purpose} (what this study will find out, for whom)\\
         {\scriptsize\color{kgrey}``to identify which factors drive password
          reuse among employees of small businesses, so that training can
          target them''}};
      \node[gbox, text width=5.2cm, below=of pu] (rq)
        {\textbf{Research question}\\ {\scriptsize\color{kgrey}$\rightarrow$ next session!}};
      \draw[karrow] (t) -- (pr);
      \draw[karrow] (pr) -- (pu);
      \draw[garrow] (pu) -- (rq);
    \end{tikzpicture}}
  \end{center}
  
  Each step \textbf{narrows}: fewer words about more specific things.
\end{center}


\subsection{Five purpose statements, rated}
The same purpose statement in five drafts, from a topic to a statement with a population, a method and a use. The pattern of the good ones is stable: one problem, one owner, one intended use of the result.
\begin{center}\small
\footnotesize
  \renewcommand{\arraystretch}{1.3}
  \begin{tabular}{@{}p{9.2cm}p{4.0cm}@{}}
    \toprule
    \textbf{Draft} & \textbf{Assessment}\\
    \midrule
    ``to look into cyber security in companies'' & no problem, no owner, no so-that\\
    ``to make the client's website better'' & a goal, not a purpose --- better \emph{how}, for \emph{whom}?\\
    ``to study why developers ignore security warnings so that tooling can be redesigned'' & problem + so-that: workable\\
    ``to prove that TypeScript reduces bugs'' & ``prove'' presumes the answer --- investigate, don't prove\\
    ``to identify onboarding steps where password-manager users abort, so that the client can prioritise fixes'' & specific, owned, actionable: strong\\
    \bottomrule
  \end{tabular}
  

  The pattern of the good ones: \emph{to \textbf{[verb]} [specific thing] so
  that [someone] can [act]}.
\end{center}


\subsection{Types of contribution, with IT examples}
\begin{center}\small
\small
  \renewcommand{\arraystretch}{1.28}
  \begin{tabular}{@{}p{4.0cm}p{9.3cm}@{}}
    \toprule
    \textbf{Contribution type} & \textbf{Example}\\
    \midrule
    New or improved \textbf{product} & an accessible component library, evaluated\\
    New or improved \textbf{theory/model} & a model of why developers ignore security warnings\\
    New or improved \textbf{method} & a lighter threat-modelling method for small teams\\
    New \textbf{evidence} & measured effect of security warnings on click-through\\
    Re-interpretation / \textbf{synthesis} & systematic review of passkey adoption studies\\
    \textbf{Tool/technique transfer} & applying A/B testing practice to security UX\\
    \bottomrule
  \end{tabular}
  

  {\color{kgrey}\small Based on Oates et al.\ (2022), Ch.\ 2; for framing your
  contribution, see Presthus \& Munkvold (2016) --- short and written for juniors.}
\end{center}


\subsection{Two kinds of knowledge}
\begin{center}\small
\small
  Science is both \textbf{investigation} (the process) and \textbf{knowledge}
  (the outcome); rigorous IS research has to deliver on both. The knowledge a
  study adds falls into one of two bases:
  
  \begin{columns}[T]
    \column{0.48\textwidth}
    \begin{block}{Descriptive knowledge ($\Omega$) --- \emph{what}}
      Phenomena and the regularities among them: observations,
      classifications, measurements, laws, patterns, theories. Added by
      discovering something that existed but was not yet understood.
    \end{block}
    \column{0.48\textwidth}
    \begin{alertblock}{Prescriptive knowledge ($\Lambda$) --- \emph{how}}
      Knowledge of human-built artefacts: \textbf{constructs}, \textbf{models},
      \textbf{methods}, \textbf{instantiations} --- and design theories that
      state principles of form and function.
    \end{alertblock}
  \end{columns}
  
  A study shows its rigour and relevance by how it \emph{consumes} existing
  knowledge and what it \emph{adds} to these two bases. For your essay: say
  which kind your thesis will add --- a finding about how things are, or an
  artefact and the knowledge of how to build it. Most cyber security theses
  produce the second, and should say so.
  

  {\color{uiagrey}\scriptsize Hevner (2021), building on Gregor \& Hevner (2013) and Hevner et al.\ (2004).}
\end{center}


\subsection{Forms of contribution in IS}
\begin{center}\small
\footnotesize
  \begin{columns}[T]
    \column{0.47\textwidth}
    \textbf{Contributions to practice}
    \begin{itemize}\setlength{\itemsep}{0.1em}
      \item Lessons learned; experience report
      \item Guidelines or roadmap; heuristics
      \item Critical success factors; patterns
    \end{itemize}
    
    \textbf{Theoretical contributions}, from building blocks upwards
    \begin{itemize}\setlength{\itemsep}{0.1em}
      \item Concept, construct, rich insight, case description
      \item Framework or taxonomy; method; proposition
      \item Generative mechanism; hypothesis; model
      \item Mid-range theory; design theory; grand theory
    \end{itemize}
    \column{0.5\textwidth}
    \textbf{Four ways to contribute to an existing theory}
    \begin{itemize}\setlength{\itemsep}{0.1em}
      \item \textbf{Confirmation} or replication --- it still holds, or holds in a new setting
      \item \textbf{Extension} --- an added construct or relationship
      \item \textbf{Contradiction} --- evidence against part of it
      \item \textbf{Elimination} --- a part is obsolete in this setting
    \end{itemize}
    
    \begin{block}{Their advice}
      (1) Balance your own idea against what is already known; (2) be
      critical of what you read --- does the field need one more study on
      this?; (3) decide the intended contribution \emph{before} the study
      starts.
    \end{block}
  \end{columns}
  
  {\color{uiagrey}\scriptsize Most master's theses confirm, extend or contradict an existing model --- sufficient, provided it is stated precisely. Munkvold: Department of Information Systems, UiA.}
\end{center}


\subsection{Practice-based knowledge production}
\begin{center}\small
\centering
  \footnotesize
  \begin{tabular}{@{}L{3.3cm}L{4.2cm}L{4.6cm}@{}}
    \toprule
    \textbf{Type} & \textbf{Examples} & \textbf{Source of legitimacy} \\
    \midrule
    Formal standardisation & ISO/IEC 27001, ISO 31000 & Consensus of national standards bodies \\
    Government frameworks & NIST CSF, NSM grunnprinsipper, BSI IT-Grundschutz & State mandate \\
    Professional bodies & COBIT (ISACA), IIA standards & Profession and certification \\
    Consulting knowledge & Maturity models, frameworks & Market adoption, references \\
    \bottomrule
  \end{tabular}
  \par
  \raggedright
  \src{These sources are normatively binding without having been tested empirically.}
\end{center}

\begin{center}\small
\centering
  \footnotesize
  \begin{tabular}{@{}L{3.1cm}L{4.4cm}L{4.4cm}@{}}
    \toprule
    \textbf{Dimension} & \textbf{Research} & \textbf{Standards and practice} \\
    \midrule
    Type of claim & Descriptive, explanatory & Prescriptive, normative \\
    Validation & Evidence, replication, peer review & Consensus, voting, uptake \\
    Error correction & Refutation, revision of theory & Revision cycle of the document \\
    Pace & Slow & Fast, market-driven \\
    Reach & Generalisation is the aim & Applicability is the aim \\
    Audience & The research community & Organisation, auditor, regulator \\
    Failure means & The hypothesis is rejected & The standard is not adopted \\
    \bottomrule
  \end{tabular}
\end{center}


\subsection{Huy et al.\ (2024) through the 6Ps}
\begin{center}\small
\small
  \renewcommand{\arraystretch}{1.12}
  \begin{tabular}{@{}p{2.4cm}p{10.9cm}@{}}
    \toprule
    Purpose       & Which factors drive ChatGPT use, and does use lead to continued use and recommendation? Gap: UTAUT2 and task--technology fit not yet combined for generative AI; curiosity untested\\
    Products      & An integrated model; a six-item scale for ChatGPT use; evidence that most UTAUT2 and TTF factors matter while effort expectancy, social influence and trust do not; implications for providers\\
    Process       & Deductive: literature $\rightarrow$ model and hypotheses $\rightarrow$ questionnaire (five-point Likert, back-translated, piloted) $\rightarrow$ EFA/CFA, then covariance-based SEM\\
    Participants  & 671 ChatGPT users in Vietnam, convenience sample, face-to-face in four cities and online; 20 interviews and five AI professors to develop the use scale\\
    Paradigm      & Positivist: constructs measured, hypotheses tested, one model for all respondents\\
    Presentation  & Journal article: introduction, theory and hypotheses, methodology, results, discussion with contributions and limitations\\
    \bottomrule
  \end{tabular}
  

  Note the 20 interviews inside a quantitative study: a developmental
  mixed-methods element (session 2), used to build the instrument, not to
  answer the question.
\end{center}


\subsection{A model of the research process}
\begin{center}\small
\begin{center}
    \resizebox{\textwidth}{!}{%
    \begin{tikzpicture}[node distance=0.35cm and 0.5cm]
      \node[gbox, text width=3.1cm] (exp) {\textbf{Experiences \&\\motivation}};
      \node[kbox, text width=3.1cm, right=of exp] (lit) {\textbf{Literature\\review}};
      \node[kbox, text width=3.1cm, right=of lit] (rq) {\textbf{Research\\question(s)}};
      \node[kbox, text width=3.1cm, right=of rq] (str) {\textbf{Strategy}};
      \node[kbox, text width=3.1cm, below=0.9cm of str] (dgm) {\textbf{Data generation\\method(s)}};
      \node[kbox, text width=3.1cm, left=of dgm] (ana) {\textbf{Data\\analysis}};
      \node[gbox, text width=3.1cm, left=of ana] (ans) {\textbf{Answers \&\\contribution}};
      \draw[karrow] (exp) -- (lit);
      \draw[karrow] (lit) -- (rq);
      \draw[karrow] (rq) -- (str);
      \draw[karrow] (str) -- (dgm);
      \draw[karrow] (dgm) -- (ana);
      \draw[karrow] (ana) -- (ans);
      \draw[garrow, dashed] (ans.north) to[out=90,in=-140] node[above,font=\scriptsize,kgrey,pos=0.4]{new questions emerge} (rq.south west);
    \end{tikzpicture}}
  \end{center}
  
  \small Not a straight line in practice --- the dashed arrow is where real
  projects loop. Your \textbf{essay} describes one planned pass through
  this model. \hfill {\color{kgrey}\small Adapted from Oates et al.\ (2022), Ch.\ 3}
\end{center}


\section*{Terms from this session}
\begin{center}\small
\footnotesize
  \renewcommand{\arraystretch}{1.2}
  \begin{tabular}{@{}p{3.4cm}p{10.1cm}@{}}
    \toprule
    \textbf{6Ps} & Purpose, Products, Process, Participants, Paradigm, Presentation --- the aspects every study must consider\\
    \textbf{Purpose} & why the study is done: topic, importance, question(s), objectives\\
    \textbf{Products} & the outcomes, especially the contribution to knowledge\\
    \textbf{Topic vs.\ problem} & broad area of interest vs.\ a specific wrong/unknown/needed thing with an owner and a cost\\
    \textbf{Purpose statement} & one sentence: \emph{to (find out X) so that (Y benefits)}\\
    \textbf{Output vs.\ contribution} & the thing produced vs.\ what the world knows or can do now\\
    \textbf{Research process} & literature $\rightarrow$ questions $\rightarrow$ strategy $\rightarrow$ data generation $\rightarrow$ analysis $\rightarrow$ answers\\
    \textbf{Master's expectation} & not original research --- methodology known, applied well, justified\\
    \bottomrule
  \end{tabular}
\end{center}



\section*{Questions from the reading}

  \begin{enumerate}\setlength{\itemsep}{0.4em}
    \item What are the \textbf{6Ps} --- and which two are the focus of Chapter 2?
          \answer{Purpose, Products, Process, Participants, Paradigm, Presentation --- Ch.\ 2 focuses on \textbf{purpose} and \textbf{products}.}
    \item Oates lists several \textbf{reasons for doing research}: which one applies most to you this semester?
          \answer{Personal --- but the book's list includes: solving a problem, curiosity, qualification, contributing to knowledge. For you, the honest first answer is the qualification --- and that is fine.}
    \item What is the difference between an \textbf{output} and a \textbf{contribution to knowledge}?
          \answer{Outputs are everything the project produces; the \textbf{contribution} is the \emph{new knowledge} added --- new or improved evidence, methodology, analysis, concepts/theories, or products.}
    \item In Oates' model of the research process (Fig.\ 3.1): which component follows directly after the literature review?
          \answer{The \textbf{research question(s)} --- in Fig.\ 3.1 the literature review feeds into (and refines) the questions before strategy and methods follow.}
  \end{enumerate}

\section*{Questions for discussion}

  \begin{enumerate}\setlength{\itemsep}{0.25em}
    \item \textbf{Your topic through the funnel} --- topic $\rightarrow$ problem (who has it, what does it cost?) $\rightarrow$ purpose statement (\emph{to \dots\ so that \dots}). Can you say it in one sentence?
          \answer{If the sentence needs two ``and''s, the topic is still too wide. The purpose statement names one problem, one owner and one intended use of the result --- that is Part I of the essay in miniature.}
    \item \textbf{The hardest P} --- in the phishing-warning study, which of the 6Ps is hardest to fill in from the description, and why?
          \answer{Usually products and paradigm: the study reports what it did, but rarely what kind of knowledge it claims to add or which worldview its design assumes. Both must be made explicit in your essay, which is why they have their own sessions.}
    \item \textbf{Your contribution} --- one sentence for practice and, if you can, one for academia. Who would act differently because of your result?
          \answer{A contribution to practice names a decision that becomes better informed; a contribution to academia names a gap in the literature that becomes smaller. At master's level one of the two, clearly stated, is sufficient.}
  \end{enumerate}

\section*{Key points}
\begin{itemize}
  \item The 6Ps: purpose, products, process, participants, paradigm, presentation --- Huy et al.\ (2024) as the worked example
  \item From topic to problem to purpose --- the funnel
  \item Contribution: descriptive or prescriptive knowledge (Hevner); the forms a junior researcher can target (Presthus \& Munkvold)
  \item Two systems of knowledge production: research tests, standards prescribe --- cite standards, do not treat them as evidence
\end{itemize}

\section*{Sources}
\begin{itemize}\small
  \item Hevner, A.\,R. (2021). The duality of science: Knowledge in information systems research. \emph{Journal of Information Technology}, \emph{36}(1), 72--76. \url{https://doi.org/10.1177/0268396220945714}
  \item Huy, L.\,V., Nguyen, H.\,T.\,T., Vo-Thanh, T., Thinh, N.\,H.\,T., \& Tran, T.\,T.\,D. (2024). Generative AI, why, how, and outcomes: A user adoption study. \emph{AIS Transactions on Human-Computer Interaction}, \emph{16}(1), 1--27. \url{https://doi.org/10.17705/1thci.00198}
  \item Presthus, W., \& Munkvold, B.\,E. (2016). How to frame your contribution to knowledge? A guide for junior researchers in information systems. \emph{NOKOBIT}, \emph{24}(1). Paper presented at NOKOBIT 2016, Bergen.
  \item McAfee, A., \& Brynjolfsson, E. (2008). Investing in the IT that makes a competitive difference. \emph{Harvard Business Review}, \emph{86}(7/8), 98--107.
  \item DiMaggio, P.\,J., \& Powell, W.\,W. (1983). The iron cage revisited: Institutional isomorphism and collective rationality in organizational fields. \emph{American Sociological Review}, \emph{48}(2), 147--160.
  \item Gibbons, M., Limoges, C., Nowotny, H., Schwartzman, S., Scott, P., \& Trow, M. (1994). \emph{The new production of knowledge}. SAGE.
  \item Hevner, A.\,R., March, S.\,T., Park, J., \& Ram, S. (2004). Design science in information systems research. \emph{MIS Quarterly}, \emph{28}(1), 75--105.
  \item Meyer, J.\,W., \& Rowan, B. (1977). Institutionalized organizations: Formal structure as myth and ceremony. \emph{American Journal of Sociology}, \emph{83}(2), 340--363.
  \item Peffers, K., Tuunanen, T., Rothenberger, M.\,A., \& Chatterjee, S. (2007). A design science research methodology for information systems research. \emph{Journal of Management Information Systems}, \emph{24}(3), 45--77.
  \item Rousseau, D.\,M. (2006). Is there such a thing as evidence-based management? \emph{Academy of Management Review}, \emph{31}(2), 256--269.
  \item Saunders, M., Lewis, P., \& Thornhill, A. (2019). \emph{Research methods for business students} (8th ed.). Pearson.
  \item Van de Ven, A.\,H., \& Johnson, P.\,E. (2006). Knowledge for theory and practice. \emph{Academy of Management Review}, \emph{31}(4), 802--821.
  \item Oates, B.\,J., Griffiths, M., \& McLean, R. (2022). \emph{Researching information systems and computing} (2nd ed.). SAGE. --- Chapters 1--3 (Introduction; The purpose and products of research; Overview of the research process).
  \item Laudon, K.\,C., \& Laudon, J.\,P. (2022). \emph{Management information systems: Managing the digital firm} (17th ed., Global ed.). Pearson. --- Chapter 8 (Securing information systems) and Chapter 11 (Managing knowledge and artificial intelligence) as maps of possible topics.
\end{itemize}

\part{Design}
\chapter{Formulating Research Questions}
\label{ch:4}
\sessionhead{4}{Thursday 1 October 2026 \;$\cdot$\; reading: Oates, Chapter 2}

A research question is the pivot of the essay: it decides the approach, the strategy and the feasibility of the study. This chapter separates aim, objective, question and hypothesis, presents the shapes a question can take, and shows the five ways in which questions fail --- and how each failure is repaired. Real research questions from the papers read in this course serve as the reference; the chapter ends with the question--strategy fit that the following chapters develop.


\section{Anatomy}

Four terms are often used as if they were one, and the essay must keep them apart.

\begin{block}{Aim, objective, research question, hypothesis}
The aim refers to the overall intent of the study in one sentence. Objectives refer to the concrete, checkable steps towards the aim. A research question refers to the question the study answers. A hypothesis refers to a testable predicted answer to a quantitative question.
\end{block}

An example makes the layers visible. The aim is to understand why employees reuse passwords; the objectives are to survey 150 employees, to interview twelve of them, and to compare the two; the research question is which factors are associated with password reuse among SME employees; and the hypothesis is that perceived account value is negatively associated with reuse. A qualitative study has no hypothesis; a quantitative study can have several.

\begin{block}{Research question}
A research question refers to a question that is answerable through systematic investigation --- by collecting and analysing data --- within the resources of the study.
\end{block}

Five criteria decide whether a draft qualifies. It must be anchored, following directly from the problem and purpose; focused, one thing at a time and not double-barrelled; researchable, so that data could actually answer it rather than only opinion, prophecy or morality; feasible with the resources of a three-month study; and open, in the sense that it is not trivially yes or no, and invites analysis.

Real research questions from the papers read in this course show the range. Presthus and Munkvold (2016) ask what forms of knowledge contribution can be targeted in information systems research --- a descriptive question answered by a taxonomy. Enholm et al.\ (2022) ask three review questions: what enables or inhibits AI use in organisations, what types of AI use exist, and through which mechanisms AI value is realised --- one topic, three facets. Woodruff et al.\ (2024) ask how knowledge workers expect generative AI to affect their industries --- exploratory and qualitative, answered by workshops. Boodraj et al.\ (2026) ask which skills and competencies are most in demand for cyber security project managers --- descriptive, answered from 161,000 job postings. Huy et al.\ (2024) state no question at all but a model with hypotheses; in quantitative work the hypothesis often replaces the question. Two observations follow: review and qualitative papers state their questions verbatim, and every question names its population and its phenomenon. Test your own draft against the same two markers.

\section{Question types}

Qualitative questions follow a script. They begin with ``how'' or ``what'', name a central concept, use exploratory verbs --- explore, understand, describe, discover --- and specify the participants and the setting: ``how do security analysts make sense of alert fatigue in a Norwegian SOC?''. Quantitative questions take one of three shapes. Descriptive questions ask how frequent or large something is: ``how widespread is password reuse among SME employees?''. Comparative questions ask whether groups differ: ``do developers with security training produce fewer vulnerabilities than those without?''. Relationship questions ask whether variables are associated: ``is perceived account value associated with password reuse?''. The relationship question often comes with a hypothesis, a predicted, testable answer derived from theory.

Two templates help when the question will not settle. PICO --- population, intervention, comparison, outcome --- turns ``does security training help?'' into ``among SME employees (P), does quarterly phishing simulation (I), compared with annual e-learning (C), reduce click rates on phishing e-mails (O) within six months?''. GQM --- goal, question, metric --- was developed for software engineering and works for any evaluation: the goal is stated, the questions that would show whether it is met are derived, and each question receives a measurable metric.

\section{From bad to good}

Questions fail in five recognisable ways. Too broad: ``how can cyber security be improved?'' is a career, not a study. Trivial yes or no: ``is phishing dangerous?'' is answered before it is asked. Not researchable: ``should companies be punished for breaches?'' is a value judgement. Double-barrelled: ``how do users choose and remember passwords?'' is two studies. Lastly, assumption baked in: ``why does TypeScript make developers happier?'' assumes the effect before anyone has shown it. Every one of these is repaired the same way: define the terms, name the population, and turn the assumption into a sub-question. ``Is zero trust better than perimeter security?'' becomes ``how does the adoption of zero-trust access controls affect the number of successful lateral movements in mid-sized organisations?'' --- and the variables, the population and the comparison become visible.

A study usually has one main question that steers it and at most two sub-questions that decompose it. If the sub-questions could each be a study of their own, the main question is a topic, not a question. Design and creation questions carry an additional part: the artefact and its evaluation. ``Can we build a tool that makes code review faster?'' becomes ``does an automated pre-review of pull requests reduce human review time without reducing defect detection, in a team of eight developers over three months?'' --- the evaluation criteria, the population and the period are now part of the question.

\section{Fit check}

The question decides more than the essay's first page. A ``how widespread'' question implies a survey; a ``how do people make sense'' question implies interviews or ethnography; a ``does X cause Y'' question implies an experiment; a ``does the artefact work'' question implies design and creation with an evaluation. The next two sessions take these strategies one by one, and the question written in this session is tested against each of them.


\section{Worked examples and tables}

\subsection{The four layers}
\begin{center}\small
\small
  \renewcommand{\arraystretch}{1.4}
  \begin{tabular}{@{}p{2.6cm}p{5.1cm}p{5.3cm}@{}}
    \toprule
    \textbf{Term} & \textbf{What it is} & \textbf{Example (password reuse)}\\
    \midrule
    \textbf{Aim} & the overall intent, one sentence & understand why employees reuse passwords\\
    \textbf{Objective(s)} & concrete, checkable steps toward the aim & survey 150 employees; interview 8 of them\\
    \textbf{Research question} & the question the study will answer & which factors are associated with password reuse among SMB employees?\\
    \textbf{Hypothesis} & a testable predicted answer (quantitative only) & H1: perceived account value is negatively associated with reuse\\
    \bottomrule
  \end{tabular}
  

  Aim = destination, objectives = milestones, question = what you ask the
  world, hypothesis = your predicted answer.
  \hfill {\color{kgrey}\small Based on Creswell \& Creswell (2018), Ch.\ 7}
\end{center}


\subsection{Question types}
\begin{center}\small
\begin{columns}[T]
    \column{0.55\textwidth}
    Qualitative studies ask \textbf{one central question} plus a few
    sub-questions. Typical shape:
    
    \begin{itemize}\setlength{\itemsep}{0.4em}
      \item begins with \textbf{how} or \textbf{what} (not \emph{why} ---
            too causal for an exploratory study)
      \item uses exploratory verbs: \emph{explore, understand, describe,
            discover}
      \item names the \textbf{participants} and the \textbf{setting}
    \end{itemize}
    \column{0.42\textwidth}
    \begin{block}{Example}
      \emph{What does ``being secure'' mean to system administrators in
      Norwegian small businesses?}\\[0.3em]
      {\small Sub: How do they describe trade-offs between security and
      day-to-day operations?}
    \end{block}
  \end{columns}
  
  {\color{kgrey}\small Based on Creswell \& Creswell (2018), Ch.\ 7}
\end{center}

\begin{center}\small
\small
  \renewcommand{\arraystretch}{1.4}
  \begin{tabular}{@{}p{2.9cm}p{4.6cm}p{5.6cm}@{}}
    \toprule
    \textbf{Shape} & \textbf{Template} & \textbf{Example}\\
    \midrule
    \textbf{Descriptive} & how frequent/large is X in group G? & how widespread is password reuse among SMB employees?\\
    \textbf{Comparative} & does group A differ from group B on X? & do developers with security training reuse passwords less than those without?\\
    \textbf{Relationship} & is X associated with / does X affect Y? & does perceived account value affect password reuse?\\
    \bottomrule
  \end{tabular}
  

  Relationship questions often come with \textbf{hypotheses} --- predicted,
  testable answers derived from theory. Variables must be
  \textbf{measurable}, and the question names them explicitly.
  \hfill {\color{kgrey}\small Based on Creswell \& Creswell (2018), Ch.\ 7}
\end{center}


\subsection{PICO and GQM}
Two templates for questions that will not settle: PICO for evaluations with a comparison, GQM for evaluations of a goal through questions and metrics.
\begin{center}\small
\small Four slots turn a vague comparison idea into a complete question:
  
  \footnotesize
  \renewcommand{\arraystretch}{1.3}
  \begin{tabular}{@{}p{2.8cm}p{3.3cm}p{3.3cm}p{3.3cm}@{}}
    \toprule
    & \textbf{Cyber} & \textbf{Frontend} & \textbf{Backend}\\
    \midrule
    \textbf{P}opulation & SMB employees & junior developers & small on-call teams\\
    \textbf{I}ntervention & quarterly phishing simulations & accessible component library & pair programming\\
    \textbf{C}omparison & annual e-learning & plain HTML/CSS & solo + review\\
    \textbf{O}utcome & click-through rate & task time, WCAG violations & defect density\\
    \bottomrule
  \end{tabular}
  

  Read any column top to bottom and the question assembles itself:
  \emph{``Among [P], does [I], compared to [C], improve [O]?''} --- if a slot
  stays empty, the question is not yet a comparison, and PICO has told you
  exactly what is missing.
  \hfill {\color{kgrey}\scriptsize Schardt et al.\ (2007)}
\end{center}

\begin{center}\small
\small \textbf{Goal} (object, purpose, viewpoint, context) $\rightarrow$
  \textbf{Questions} $\rightarrow$ \textbf{Metrics}:
  \begin{center}
    \resizebox{0.92\textwidth}{!}{%
    \begin{tikzpicture}
      \node[kbox, text width=6.4cm, minimum height=1.0cm] (g) at (0,0)
        {\textbf{Goal:} improve the password-manager \emph{onboarding}
         (object) for \emph{new users} (viewpoint) at the client (context)};
      \node[gbox, text width=4.4cm, minimum height=0.85cm] (q1) at (-3.0,-1.7)
        {\textbf{Q1:} where do users abort?};
      \node[gbox, text width=4.4cm, minimum height=0.85cm] (q2) at (3.0,-1.7)
        {\textbf{Q2:} how long does setup take?};
      \node[kbox, text width=4.4cm, minimum height=0.85cm] (m1) at (-3.0,-3.2)
        {\textbf{M1:} drop-off rate per step};
      \node[kbox, text width=4.4cm, minimum height=0.85cm] (m2) at (3.0,-3.2)
        {\textbf{M2:} median time to first stored login};
      \draw[karrow] (g) -- (q1); \draw[karrow] (g) -- (q2);
      \draw[karrow] (q1) -- (m1); \draw[karrow] (q2) -- (m2);
    \end{tikzpicture}}
  \end{center}
  \footnotesize Second example, one line: \emph{Goal:} CI pipeline
  reliability (developer viewpoint) $\rightarrow$ \emph{Q:} how often do
  builds fail for non-code reasons? $\rightarrow$ \emph{M:} share of failed
  builds with infrastructure causes. The discipline: \textbf{no metric
  without a question, no question without a goal}.
  \hfill {\color{kgrey}\scriptsize Basili (1994); Molléri \& Petersen (2024)}
\end{center}


\subsection{The five failure modes, and their repairs}
\begin{center}\small
\small
  \renewcommand{\arraystretch}{1.35}
  \begin{tabular}{@{}p{3.3cm}p{9.9cm}@{}}
    \toprule
    \textbf{Failure} & \textbf{Symptom}\\
    \midrule
    Too broad & ``How can cyber security be improved?'' --- a career, not a study\\
    Trivial yes/no & ``Is phishing dangerous?'' --- answerable without data\\
    Not researchable & ``Should companies be punished for breaches?'' --- a value judgement\\
    Double-barreled & ``How do users choose passwords and why do they reuse them?'' --- two studies\\
    Assumption baked in & ``Why does TypeScript make teams faster?'' --- presumes the effect exists\\
    \bottomrule
  \end{tabular}
  

  Every one of these appears in essay drafts every year. Learn to see them
  in \textbf{your own} questions.
\end{center}

\begin{center}\small
\begin{center}
    \begin{tikzpicture}[node distance=0.4cm]
      \node[gbox, text width=11.4cm] (v1)
        {\textbf{Draft:} ``How can companies protect themselves against phishing?''\\
         {\scriptsize\color{kgrey}too broad; ``protect'' unmeasurable; which companies?}};
      \node[kbox, text width=11.4cm, below=of v1] (v2)
        {\textbf{Better:} ``Does anti-phishing training reduce click rates?''\\
         {\scriptsize\color{kgrey}focused --- but which training, whom, over what period?}};
      \node[kbox, text width=11.4cm, below=of v2] (v3)
        {\textbf{Good:} ``To what extent does quarterly simulation-based
         anti-phishing training reduce click-through rates on phishing e-mails
         among employees of Norwegian SMBs over six months?''\\
         {\scriptsize\color{kgrey}population, intervention, outcome, timeframe --- PICO-complete}};
      \draw[karrow] (v1) -- (v2);
      \draw[karrow] (v2) -- (v3);
    \end{tikzpicture}
  \end{center}
\end{center}

\begin{center}\small
\small
  \begin{center}
    \begin{tikzpicture}[node distance=0.35cm]
      \node[gbox, text width=11.6cm] (v1)
        {\textbf{Draft:} ``Can we build a tool that makes code reviews better?''\\
         {\scriptsize\color{kgrey}``better'' unmeasured; ``can we build'' is almost always yes}};
      \node[kbox, text width=11.6cm, below=of v1] (v2)
        {\textbf{Good:} ``To what extent does an automated review-checklist
         bot reduce review turnaround time and missed-defect rate in a small
         backend team, compared to the team's current process?''\\
         {\scriptsize\color{kgrey}artefact named, two outcomes, comparison baseline, population}};
      \draw[karrow] (v1) -- (v2);
    \end{tikzpicture}
  \end{center}
  
  For design \& creation questions, the repair always adds the
  \textbf{evaluation}: outcomes and a baseline. ``Can it be built'' is a
  development milestone; ``how well does it work, compared to what'' is the
  research question.
\end{center}


\subsection{Hierarchies and hypotheses}
\begin{center}\small
\begin{center}
    \begin{tikzpicture}[node distance=0.45cm and 0.5cm]
      \node[kbox, text width=10.6cm] (mq)
        {\textbf{Main RQ:} which factors are associated with password reuse
         among SMB employees?};
      \node[gbox, text width=5.1cm, below left=0.5cm and -4.6cm of mq] (s1)
        {\textbf{SQ1:} how widespread is reuse in the studied SMBs? {\scriptsize(descriptive)}};
      \node[gbox, text width=5.1cm, below right=0.5cm and -4.6cm of mq] (s2)
        {\textbf{SQ2:} how do employees explain their reuse decisions? {\scriptsize(qualitative)}};
      \draw[karrow] (mq) -- (s1);
      \draw[karrow] (mq) -- (s2);
    \end{tikzpicture}
  \end{center}
  
  Rules of thumb:
  \begin{itemize}\setlength{\itemsep}{0.3em}
    \item Sub-questions \textbf{decompose} the main question --- together they
          answer it; none reaches beyond it.
    \item Two to three sub-questions are plenty at master's level.
    \item Notice: SQ1 + SQ2 above quietly sketch a \textbf{mixed-methods}
          design --- the question hierarchy and the strategy must fit
          (sessions 6--7).
  \end{itemize}
\end{center}

\begin{center}\small
\footnotesize
  \begin{columns}[T]
    \column{0.48\textwidth}
    \textbf{Design \& creation project}\\[0.2em]
    \textbf{Main:} how well does a review-checklist bot reduce turnaround in
    a small team?
    \begin{itemize}\setlength{\itemsep}{0.2em}
      \item SQ1: which checklist items catch the most defects? {\scriptsize(log analysis)}
      \item SQ2: how do reviewers experience the bot? {\scriptsize(short interviews)}
    \end{itemize}
    \column{0.48\textwidth}
    \textbf{Case study project}\\[0.2em]
    \textbf{Main:} how does one SMB rebuild operational trust after a
    ransomware incident?
    \begin{itemize}\setlength{\itemsep}{0.2em}
      \item SQ1: which routines changed, and when? {\scriptsize(documents, timeline)}
      \item SQ2: how do staff explain what made the changes stick? {\scriptsize(interviews)}
    \end{itemize}
  \end{columns}
  
  Same rule both times: sub-questions \textbf{decompose} the main question
  --- and each quietly names its data source, which is exactly what
  Parts III--IV will need.
\end{center}

\begin{center}\small
\small
  \renewcommand{\arraystretch}{1.4}
  \begin{tabular}{@{}p{7.0cm}p{6.0cm}@{}}
    \toprule
    \textbf{Hypothesis} & \textbf{Assessment}\\
    \midrule
    H1: employees with security training reuse fewer passwords than those without & testable: two groups, one measurable outcome\\
    H2: the new warning design is better & not testable --- ``better'' has no measure or direction\\
    H3: perceived account value is negatively associated with reuse & testable: named variables, stated direction, theory-derived\\
    \bottomrule
  \end{tabular}
  

  A hypothesis needs: \textbf{named variables}, a \textbf{direction}, and a
  route to \textbf{measurement}. And remember: hypotheses belong to
  quantitative questions --- a qualitative study states expectations, not
  hypotheses.
\end{center}


\subsection{The question decides more than you think}
\begin{center}\small
\small
  \renewcommand{\arraystretch}{1.35}
  \begin{tabular}{@{}p{4.6cm}p{4.0cm}p{4.4cm}@{}}
    \toprule
    \textbf{If your question says\dots} & \textbf{\dots it implies} & \textbf{\dots and demands}\\
    \midrule
    \emph{how many, to what extent, does X affect Y} & quantitative & measurable variables, enough participants\\
    \emph{how do people experience, what does X mean} & qualitative & access to participants, interview time\\
    \emph{what predicts X --- and how do people explain it} & mixed methods & both of the above (cost!)\\
    \emph{can an artefact that does X be built and how well does it work} & design \& creation & build skills \emph{and} an evaluation plan\\
    \bottomrule
  \end{tabular}
  

  Feasibility is a \textbf{quality criterion}, not an afterthought: a
  beautiful question your group cannot answer by June is a bad question
  \emph{for this essay}.
\end{center}


\section*{Terms from this session}
\begin{center}\small
\footnotesize
  \renewcommand{\arraystretch}{1.2}
  \begin{tabular}{@{}p{3.4cm}p{10.1cm}@{}}
    \toprule
    \textbf{Aim} & the overall intent of the study, one sentence\\
    \textbf{Objective} & a concrete, checkable step toward the aim\\
    \textbf{Research question} & the question the study answers through systematic investigation\\
    \textbf{Hypothesis} & a testable predicted answer, derived from theory (quantitative)\\
    \textbf{Sub-question} & a component question that helps decompose and answer the main question\\
    \textbf{PICO} & Population, Intervention, Comparison, Outcome --- a template for comparative questions\\
    \textbf{GQM} & Goal, Question, Metric --- a measurement-driven template from software engineering\\
    \textbf{Double-barreled} & a question that secretly asks two things at once --- split it\\
    \bottomrule
  \end{tabular}
\end{center}



\section*{Questions from the reading}

  \begin{enumerate}\setlength{\itemsep}{0.4em}
    \item Oates lists several \textbf{reasons} for doing research. Name three that could apply to a master's thesis.
          \answer{For example: to add to the body of knowledge, to solve a problem, to find evidence to inform practice, to test or refine a theory, to develop a system and learn from it --- and, honestly, to obtain a degree.}
    \item Which \textbf{products} can research produce, according to Oates?
          \answer{A new or improved product, a new theory, a re-interpretation of an existing theory, an in-depth analysis of a situation, an exploration of a topic, a critical analysis, a new research tool or technique, a new or improved model.}
    \item Where do \textbf{research topics} come from?
          \answer{Your own interests and experience, the literature (gaps and ``further research'' sections), practice and work problems, supervisors and current debates in the field.}
    \item Which \textbf{criteria} does Oates propose for judging a candidate topic?
          \answer{Interest to you, feasibility with the time, skills and access you have, a clear purpose, and a contribution that others would recognise --- plus the ethical and practical constraints of the setting.}
  \end{enumerate}

\section*{Questions for discussion}

  \begin{enumerate}\setlength{\itemsep}{0.25em}
    \item \textbf{Your main question} --- draft it in one sentence, then run the five-point checklist: precise, feasible, researchable, single-barrelled, no assumption baked in. Which criterion fails first?
          \answer{Most first drafts fail on precision (``effective'', ``better'') or on a baked-in assumption (``why does X cause Y?'' before anyone has shown that it does). Repair by defining the terms and by turning the assumption into a sub-question.}
    \item \textbf{One question or several?} --- your topic tempts you to ask three things. What does the hierarchy look like?
          \answer{One main question that steers the study; at most two sub-questions that decompose it. If the sub-questions could each be a study of their own, the main question is a topic, not a question.}
    \item \textbf{Repair a neighbour's question} --- ``Is zero trust better than perimeter security?'' Diagnose and rewrite.
          \answer{Trivial yes/no and undefined terms. One repair: ``How does the adoption of zero-trust access controls affect the number of successful lateral movements in mid-sized organisations?'' --- now the variables, the population and the comparison are visible.}
  \end{enumerate}

\section*{Key points}
\begin{itemize}
  \item Aim, objective, question, hypothesis --- four layers that must not blur
  \item Question types: qualitative scripts, quantitative shapes, PICO and GQM
  \item The five failure modes --- and how to repair a question
  \item The question decides the approach, the strategy and the feasibility
\end{itemize}

\section*{Sources}
\begin{itemize}\small
  \item Enholm, I.\,M., Papagiannidis, E., Mikalef, P., \& Krogstie, J. (2022). Artificial intelligence and business value: A literature review. \emph{Information Systems Frontiers}, \emph{24}, 1709--1734. \url{https://doi.org/10.1007/s10796-021-10186-w}
  \item Huy, L.\,V., Nguyen, H.\,T.\,T., Vo-Thanh, T., Thinh, N.\,H.\,T., \& Tran, T.\,T.\,D. (2024). Generative AI, why, how, and outcomes: A user adoption study. \emph{AIS Transactions on Human-Computer Interaction}, \emph{16}(1), 1--27. \url{https://doi.org/10.17705/1thci.00198}
  \item Presthus, W., \& Munkvold, B.\,E. (2016). How to frame your contribution to knowledge? A guide for junior researchers in information systems. \emph{NOKOBIT}, \emph{24}(1). Paper presented at NOKOBIT 2016, Bergen.
  \item Woodruff, A., Shelby, R., Kelley, P.\,G., Rousso-Schindler, S., Smith-Loud, J., \& Wilcox, L. (2024). How knowledge workers think generative AI will (not) transform their industries. In \emph{Proceedings of the CHI Conference on Human Factors in Computing Systems (CHI '24)}. ACM. \url{https://doi.org/10.1145/3613904.3642700}
  \item Boodraj, M., Fairbanks, J., Davies, C., \& Akcakaya, O. (2026). Beyond the firewall: The rise of the cybersecurity project manager. In \emph{AMCIS 2026 Proceedings} (Emergent Research Forum paper). AIS. \url{https://aisel.aisnet.org/amcis2026/sig_itpm/sig_itpm/2}
  \item Oates, B.\,J., Griffiths, M., \& McLean, R. (2022). \emph{Researching information systems and computing} (2nd ed.). SAGE. --- Chapter 2 (The purpose and products of research).
\end{itemize}

\chapter{Reviewing the Literature}
\label{ch:5}
\sessionhead{5}{Thursday 8 October 2026 \;$\cdot$\; reading: Oates, Chapter 6}

A literature review does four jobs: it justifies the problem, avoids reinventing what exists, grounds the concepts, and informs the method. This chapter turns the review into a documented, repeatable procedure --- from the research question to a keyword table, a Boolean string, explicit screening criteria and a PRISMA flow --- and then into a synthesis rather than a summary. One published systematic review and my own review serve as worked examples. The chapter corresponds to Part II of the essay.


\section{Why review}

\begin{block}{Literature review}
A literature review refers to a systematic examination of the existing knowledge relevant to a research question, undertaken both for the process of finding and evaluating it and for the written account that results.
\end{block}

The review does four jobs for the essay. First, it justifies the problem: it shows the gap, or the conflict, that the study addresses. Secondly, it avoids reinventing what exists --- someone may already have answered your question, or shown that the obvious design fails. Thirdly, it grounds the concepts: the definitions and theories the study uses. Lastly, it informs the method: how others did similar studies, and what they learned about doing them. Part II of the essay asks for the keywords and search strings, the databases and dates, the inclusion and exclusion criteria, the initial results, and a table of the included studies --- everything on that list is a step in this chapter.

vom Brocke et al.\ (2009) describe the review as a five-phase framework: definition of the review scope, conceptualisation of the topic, the literature search, the analysis and synthesis, and the research agenda. Their argument is that the search must be documented in enough detail to be repeated, which is what turns a review from a reading list into evidence. The chapter follows the phases.

\section{Where to look}

Sources differ in quality, and the quality hierarchy is not the same as the ease of finding. Peer-reviewed journals --- for our field the AIS basket of eight, \emph{IEEE Transactions on Information Forensics and Security}, \emph{Computers \& Security} --- and peer-reviewed conference proceedings --- ICIS, ECIS, USENIX Security, CCS --- form the core; books and textbooks add synthesis; grey literature such as NIST and ENISA reports, vendor white papers and industry surveys adds practice and data but must be labelled as what it is; blogs and vendor pages add currency and marketing. The databases of our field are the AIS eLibrary, IEEE Xplore, the ACM Digital Library, Scopus and Web of Science, with Google Scholar as a wide net and a poor filter. Every database has its own syntax for Boolean operators, wildcards and field limits, and a search log records which syntax was used where.

\section{Systematic search}

The search proceeds from the research question to a keyword table. Each concept in the question receives a column of synonyms and related terms; ``password reuse'' becomes ``password reuse'', ``credential reuse'' and ``password management''; ``employees'' becomes ``employees'', ``staff'', ``end users'' and ``knowledge workers''. The concept with the most synonyms is usually the technical one, with tool names, acronyms and vendor terms, and it needs the longest OR block. Boolean operators build the string: OR joins synonyms within a concept, AND joins the concepts, NOT excludes a term and is used sparingly, and quotation marks and wildcards fix phrases and stems. A first string is tested in one database, the hit count is logged, and the string is adjusted --- which is the normal course of a search, not a failure.

Screening follows explicit criteria stated before the hits are read: inclusion criteria such as peer-reviewed publications in English from 2015 onwards that study organisational settings, and exclusion criteria such as purely technical measures without a human or organisational dimension. A criterion invented after reading the hits is not a criterion but a preference. Kitchenham and Charters (2007) set out the procedure for software engineering as a protocol written first; Kitchenham et al.\ (2009) reviewed how the procedure had been applied. The PRISMA 2020 statement (Page et al., 2021) supplies the reporting standard: a flow diagram that records how many records were identified, screened, assessed for eligibility and included, and why the others were excluded. The flow diagram is a template, and every review fills it in.



Enholm, Papagiannidis, Mikalef and Krogstie (2022) show the same procedure in a published review. They wrote a protocol first, following the Cochrane handbook; they asked three questions --- what enables or inhibits AI use in organisations, what types of AI use exist, and through which mechanisms AI value is realised; they combined two keyword sets with wildcards and searched Google Scholar and nine databases, including Scopus, IEEE Xplore and the AIS library, within a stated period in September 2020; they applied explicit inclusion and exclusion criteria; two authors assessed the quality of each paper independently for rigour, credibility and relevance; and they extracted the 43 included papers into a concept matrix and synthesised them in three parts and a research agenda. Everything Part II of your essay asks for is visible in one paper.

\section{Read and synthesise}

Reading for a review is different from reading for interest. A paper is read in the right order --- abstract, conclusion, figures and tables, method, and only then the full text --- and read for extraction: what is the question, the method, the sample, the finding, and the limitation. The extraction goes into a concept matrix, which Webster and Watson (2002) introduced as the instrument of synthesis: one row per paper, one column per concept, and in each cell what the paper says about the concept. Watson and Webster (2020) extended the idea for the age of graph databases, in which the matrix becomes a network of concepts and relationships; Session 11 returns to both.

\begin{block}{Synthesis}
Synthesis refers to constructing a new, integrated statement out of several sources --- a claim that no single source makes on its own, organised by concept rather than by paper.
\end{block}

The matrix is read by column, and the column produces a paragraph: what the papers agree on, where they diverge, what explains the divergence --- method, setting, definition --- and what the divergence means for your own design. Three papers that disagree therefore appear in the essay not as three paragraphs, one per paper, but as one argument. Four moves make the argument: establish what is known, contrast where sources conflict, locate the gap, and point at the implication for the study. A review written paper by paper is a summary; a review written concept by concept is a synthesis, and the difference is visible in the first paragraph.

Two habits keep the review under control. The reference manager --- Zotero, free, with a browser button and plug-ins for Word and LaTeX --- is set up on the first day of the search, because reconstructing sources in November has ruined better essays than any examiner. And the search log, the keyword table and the concept matrix live in one spreadsheet, so that Part II can be written from the log rather than from memory.


\section{Worked examples and tables}

\subsection{The five phases, walked through}
\begin{center}\small
\footnotesize
  \renewcommand{\arraystretch}{1.12}
  \begin{tabular}{@{}p{3.3cm}p{4.6cm}p{5.2cm}@{}}
    \toprule
    \textbf{Phase} & \textbf{What you decide / do} & \textbf{Example}\\
    \midrule
    1.\ Review scope & focus, goal, coverage, audience & empirical studies on workplace password behaviour, 2015--2026\\
    2.\ Conceptualisation & working definitions of the key terms; the concepts behind the question & what counts as ``reuse''? candidate theories (security fatigue)? $\rightarrow$ your keyword table\\
    3.\ Search process & journals $\rightarrow$ databases $\rightarrow$ keywords $\rightarrow$ backward/forward $\rightarrow$ evaluation & the search log and criteria from today's steps 1--3\\
    4.\ Analysis \& synthesis & extract by \textbf{concept}, not by paper & the concept matrix (later this session)\\
    5.\ Research agenda & what remains open --- and therefore your study & the gap paragraph ending in your question\\
    \bottomrule
  \end{tabular}
  {\color{kgrey}\scriptsize vom Brocke et al.\ (2009). Phase 5 is the bridge
  from Part II to Part I: the review \emph{ends} where your question
  \emph{begins}.}
\end{center}


\subsection{Source types and databases}
\begin{center}\small
\small
  \renewcommand{\arraystretch}{1.3}
  \begin{tabular}{@{}p{4.3cm}p{5.6cm}p{3.2cm}@{}}
    \toprule
    \textbf{Source type} & \textbf{Examples} & \textbf{Use for}\\
    \midrule
    Peer-reviewed journals & \emph{IEEE TSE, Computers \& Security} & core evidence\\
    Conference proceedings & ICSE, CHI, USENIX Security, CCS & core evidence (big in computing!)\\
    Books \& textbooks & Oates; specialist monographs & concepts, methods\\
    Grey literature & standards (NIST, OWASP), industry reports (DORA, Verizon DBIR) & practice context --- flag it as grey\\
    Blogs, vendor pages & engineering blogs, marketing & leads only --- verify elsewhere\\
    \bottomrule
  \end{tabular}
  

  Computing is unusual: \textbf{top conferences are peer-reviewed and count
  as much as journals}. But a vendor whitepaper is not evidence --- session 1's
  ``95\,\% human error'' claim lives exactly there.
\end{center}

\begin{center}\small
\footnotesize
  \renewcommand{\arraystretch}{1.22}
  \begin{tabular}{@{}p{4.2cm}p{9.0cm}@{}}
    \toprule
    \textbf{Database} & \textbf{What it is good for}\\
    \midrule
    \textbf{Oria} & the library's front door --- searches most of the below\\
    \textbf{AIS eLibrary} & the information systems field's own library, incl.\ the ``Senior Scholars'' basket of top journals\\
    \textbf{ACM Digital Library} & computing conferences and journals\\
    \textbf{IEEE Xplore} & engineering and computer science\\
    \textbf{EBSCOhost} & broad multidisciplinary coverage incl.\ business and IS\\
    \textbf{Scopus / Web of Science} & broad, curated, good citation chasing\\
    \textbf{Google Scholar} & broadest recall --- but no quality filter; use for snowballing\\
    \textbf{SAGE Research Methods} & not papers --- \emph{methods}: look up any method you meet in this course\\
    \bottomrule
  \end{tabular}
  

  Strategy: run your string in \textbf{two or three} databases, not ten ---
  and \textbf{snowball}: follow the reference lists (backward) and citations
  (forward) of your best hits. {\color{kgrey}(AIS eLibrary, EBSCOhost and the
  ACM DL are the three used in the real review later in this session.)}
\end{center}


\subsection{From question to string to criteria}
The three steps of the systematic search on the running example: the keyword table, the Boolean string with its test log, and the screening criteria stated before reading.
\begin{center}\small
Take the question apart into \textbf{concepts}, then collect
  \textbf{synonyms} per concept:
  
  \small
  \renewcommand{\arraystretch}{1.35}
  \begin{tabular}{@{}p{3.6cm}p{9.6cm}@{}}
    \toprule
    \textbf{Concept} & \textbf{Synonyms \& variants}\\
    \midrule
    password reuse & password reuse, credential reuse, password sharing, password habits\\
    employees & employee*, staff, worker*, end user*\\
    small business & SMB, SME, ``small and medium-sized enterprise*''\\
    \bottomrule
  \end{tabular}
  

  Tricks of the trade: \texttt{*} truncation (\texttt{employee*} finds
  \emph{employees}), \texttt{"quotation marks"} for exact phrases, and
  check the \textbf{keywords} of your first good hit --- authors hand you
  their vocabulary for free.
\end{center}

\begin{center}\small
\begin{center}
    \begin{tikzpicture}[node distance=0.35cm]
      \node[kbox, text width=12.2cm] (s)
        {\texttt{("password reuse" OR "credential reuse") AND (employee* OR staff)}\\
         \texttt{AND (SMB OR SME OR "small business")}};
    \end{tikzpicture}
  \end{center}
  
  \begin{itemize}\setlength{\itemsep}{0.4em}
    \item \textbf{OR} joins synonyms \emph{within} a concept
          (widens the net)
    \item \textbf{AND} joins the concepts \emph{together}
          (narrows the net)
    \item \textbf{NOT} excludes --- use sparingly; it silently removes good
          hits too
  \end{itemize}
  
  \begin{alertblock}{Iterate and log}
    First string returns 4{,}000 hits? Add a concept. Returns 3? Drop one.
    \textbf{Write every attempt into the search log} --- that log is
    literally a deliverable of Part II.
  \end{alertblock}
\end{center}

\begin{center}\small
Decide \textbf{before} screening what counts:
  
  \small
  \begin{columns}[T]
    \column{0.48\textwidth}
    \textbf{Inclusion (example)}
    \begin{itemize}\setlength{\itemsep}{0.25em}
      \item peer-reviewed, 2015 or later
      \item empirical data on password behaviour
      \item workplace context
      \item English or Scandinavian language
    \end{itemize}
    \column{0.48\textwidth}
    \textbf{Exclusion (example)}
    \begin{itemize}\setlength{\itemsep}{0.25em}
      \item purely technical cracking papers
      \item consumer-only context
      \item opinion pieces, posters
      \item no full text available
    \end{itemize}
  \end{columns}
  
  Screen in two passes: first \textbf{title + abstract} (fast), then
  \textbf{full text} for the survivors. Every decision is countable ---
  which is exactly what the PRISMA figure shows.
\end{center}


\subsection{Kitchenham's process and PRISMA 2020}
\begin{center}\small
\begin{center}
    \resizebox{0.96\textwidth}{!}{%
    \begin{tikzpicture}
      \node[kbox, text width=4.1cm, minimum height=2.5cm, anchor=west] (pl) at (0,0)
        {\textbf{1.\ Planning}\\[2pt]
         {\scriptsize identify the need for a review\\
          formulate the review question(s)\\
          write the \textbf{review protocol}:\\
          sources, search string, criteria,\\
          quality-assessment plan}};
      \node[kbox, text width=4.1cm, minimum height=2.5cm, anchor=west] (co) at (4.75,0)
        {\textbf{2.\ Conducting}\\[2pt]
         {\scriptsize execute the search\\
          select studies (criteria)\\
          \textbf{quality assessment}\\
          data extraction\\
          synthesis}};
      \node[kbox, text width=4.1cm, minimum height=2.5cm, anchor=west] (re) at (9.5,0)
        {\textbf{3.\ Reporting}\\[2pt]
         {\scriptsize document every step\\
          (this is where PRISMA\\
          enters) so the review is\\
          \textbf{replicable} and its limits\\
          are visible}};
      \draw[karrow] (pl) -- (co);
      \draw[karrow] (co) -- (re);
    \end{tikzpicture}}
  \end{center}
  
  The core idea: the protocol is written \textbf{before} the search --- so
  the criteria cannot quietly drift to fit the papers you happen to like.
  \hfill {\color{kgrey}\scriptsize Kitchenham et al.\ (2009)}
\end{center}

\begin{center}\small
\footnotesize
  \begin{columns}[T]
    \column{0.47\textwidth}
    \textbf{The protocol (example contents)}
    \begin{itemize}\setlength{\itemsep}{0.25em}
      \item review question(s)
      \item sources: AIS eLibrary, EBSCOhost, ACM DL
      \item search string (documented verbatim)
      \item inclusion / exclusion criteria
      \item quality-assessment criteria and scoring
      \item data-extraction fields (the matrix columns)
    \end{itemize}
    \column{0.5\textwidth}
    \textbf{Quality assessment (from the real review)}
    \begin{itemize}\setlength{\itemsep}{0.25em}
      \item five criteria: QA1 topic relevance, QA2 research context, QA3 methodological detail, QA4 data collection, QA5 analysis approach
      \item scoring per criterion: 1 / 0.5 / 0
      \item thresholds: high $\geq 4$; medium 2.5 to $<4$; low $<2.5$ (excluded)
      \item result: 32 high, 5 medium, 0 excluded
    \end{itemize}
  \end{columns}
  
  \begin{alertblock}{\small For your essay}
    \small A proportionate protocol --- question, sources, string, criteria
    --- is precisely the ``plan'' that Part II asks you to develop and report.
  \end{alertblock}
  
  {\color{kgrey}\scriptsize Kitchenham et al.\ (2009); QA example: Amling et
  al.\ (under review)}
\end{center}

\begin{center}\small
\begin{block}{PRISMA 2020 (Page et al., 2021)}
    A \textbf{reporting guideline} for systematic reviews: a
    \textbf{27-item checklist} covering everything from the search strategy
    and selection process to synthesis and limitations --- \emph{plus} the
    well-known \textbf{flow diagram}.
  \end{block}
  
  \begin{itemize}\setlength{\itemsep}{0.4em}
    \item The flow diagram is only the \textbf{visible summary}; PRISMA
          thinking means \textbf{every decision is documented and countable}.
    \item The full diagram has \textbf{two identification streams}:
          via \emph{databases} --- and via \emph{other methods}
          (forward/backward citation search, expert contacts).
    \item It also tracks \emph{reports not retrieved} --- honesty about what
          you could not get.
  \end{itemize}
  
  For your essay: the diagram plus documented criteria are the visible
  artefacts --- the next slide shows what the complete diagram looks like in
  a real, current review.
\end{center}


\subsection{A published systematic review, step by step}
\begin{center}\small
\footnotesize
  \renewcommand{\arraystretch}{1.15}
  \begin{tabular}{@{}p{2.4cm}p{10.9cm}@{}}
    \toprule
    Protocol & Written first, following the Cochrane handbook: questions, search strategy, criteria, quality assessment, synthesis method\\
    Questions & (1) What enables or inhibits AI use in organisations? (2) What are the types of AI use? (3) Through which mechanisms is AI value realised?\\
    Search & Two keyword sets (AI technologies; organisational perspective) combined with wildcards; Google Scholar plus Scopus, Business Source Complete, Emerald, Taylor \& Francis, Springer, Web of Knowledge, ABI/Inform, IEEE Xplore, AIS library; 14--30 September 2020; separate search in the AIS basket of eight\\
    Screening & Explicit inclusion and exclusion criteria (English; peer-reviewed; organisational focus)\\
    Quality & Two authors independently: rigour, credibility, relevance\\
    Result & 43 papers extracted into a \textbf{concept matrix}; synthesis in three parts (enablers and inhibitors, typologies of use, first- and second-order effects) and a research agenda\\
    \bottomrule
  \end{tabular}
  

  \small Everything Part II of your essay asks for is visible in one paper:
  protocol, strings, databases, dates, criteria, counts, matrix. Their Figure 1
  is the PRISMA-style flow of the selection.
\end{center}


\subsection{Reading and the concept matrix}
\begin{center}\small
Nobody reads papers front to back. The efficient order:
  
  \begin{enumerate}\setlength{\itemsep}{0.4em}
    \item \textbf{Abstract} --- is this relevant at all? (30 seconds)
    \item \textbf{Conclusion} --- what do they claim to have found? (2 min)
    \item \textbf{Figures and tables} --- the evidence at a glance (2 min)
    \item \textbf{Method} --- \emph{how} do they know? Only now decide whether
          to trust the conclusion (10 min)
    \item \textbf{Full read} --- only for the handful that survive steps 1--4
  \end{enumerate}
  
  While reading, extract into your \textbf{literature table} --- or, one
  step better, a \textbf{concept matrix} (Webster \& Watson, 2002): concepts
  as columns, papers as rows. It forces synthesis by concept instead of a
  paper-by-paper shopping list, and it goes straight into Part II.
\end{center}

\begin{center}\small
Concepts as \textbf{columns} (from your conceptualisation, phase 2),
  papers as \textbf{rows}; a cell marks that the paper addresses the concept:
  
  \footnotesize
  \renewcommand{\arraystretch}{1.35}
  \begin{tabular}{@{}p{3.0cm}cccc@{}}
    \toprule
    & \textbf{Concept A} & \textbf{Concept B} & \textbf{Concept C} & \textbf{Concept D}\\
    \midrule
    Paper 1 & $\times$ & $\times$ &          & \\
    Paper 2 &          & $\times$ &          & $\times$\\
    Paper 3 & $\times$ &          & $\times$ & \\
    Paper 4 & $\times$ & $\times$ & $\times$ & $\times$\\
    \bottomrule
  \end{tabular}
  

  Two reading directions: a \textbf{row} shows what one paper covers; a
  \textbf{column} shows the state of evidence for one concept --- dense
  columns mean converging work, sparse columns mean thin evidence, empty
  columns mean a gap. Cells can carry more than a mark: a finding direction,
  a method tag, a context.
  \hfill {\color{kgrey}\scriptsize Webster \& Watson (2002)}
\end{center}

\begin{center}\small
The same structure filled in, with the P/N/B cell coding from the real
  review:
  
  \footnotesize
  \renewcommand{\arraystretch}{1.3}
  \begin{tabular}{@{}p{3.0cm}cccc@{}}
    \toprule
    & \textbf{Security fatigue} & \textbf{Convenience} & \textbf{Risk perception} & \textbf{Training effect}\\
    \midrule
    Smith (2021), survey        & N & P &   & \\
    Lee (2022), experiment      &   & P &   & P\\
    Hansen (2023), interviews   & N &   & B & \\
    Berg (2024), mixed          & N & B & N & P\\
    \bottomrule
  \end{tabular}
  

  Cells: \textbf{P} = reaction/effect reported as positive, \textbf{N} =
  negative, \textbf{B} = both. Reading the columns: fatigue is consistently
  negative (converging evidence), convenience is mostly positive with one
  ``both'' (a tension to explain), risk perception is thin and mixed (a gap
  candidate). The matrix hands you the argument moves.
  \hfill {\color{kgrey}\scriptsize Webster \& Watson (2002); Amling et al.\ (under review)}
\end{center}


\subsection{From synthesis to argument}
\begin{center}\small
Every good review paragraph makes the same four moves:
  
  \small
  \begin{enumerate}\setlength{\itemsep}{0.45em}
    \item \textbf{Establish} --- what converging sources agree on:\\
          {\color{kgrey}\emph{``Survey studies consistently report that security
          fatigue drives password reuse (Smith, 2021; Lee, 2022).''}}
    \item \textbf{Complicate} --- the conflict, exception or boundary:\\
          {\color{kgrey}\emph{``However, the only field study found no such
          effect once account value was controlled (Hansen, 2023).''}}
    \item \textbf{Locate the gap} --- what nobody has examined:\\
          {\color{kgrey}\emph{``No study has examined this in small businesses,
          where IT support is minimal.''}}
    \item \textbf{Point at your study} --- the \emph{therefore}:\\
          {\color{kgrey}\emph{``We therefore ask which factors are associated
          with password reuse among SMB employees.''}}
  \end{enumerate}
\end{center}

\begin{center}\small
\small
  \begin{itemize}\setlength{\itemsep}{0.35em}
    \item \textbf{One block = one concept = one claim.} Plan the review as
          blocks taken from your matrix columns; if you cannot state a
          block's claim in one sentence, split it.
    \item \textbf{Every claim carries evidence.} A claim without citations is
          an opinion; citations without a claim are a shopping list. The
          argument is the \emph{combination}: claim + sources + your
          reasoning why they support it.
    \item \textbf{Weigh, don't just count.} Three surveys and one field study
          disagree? Say \emph{why} the designs might explain the difference ---
          this is where session 6's knowledge is applied.
    \item \textbf{The matrix feeds the moves.} A full column = move 1
          (establish); mixed P/N in a column = move 2 (complicate); an empty
          column or row = move 3 (the gap).
  \end{itemize}
  \begin{alertblock}{The test, again}
    \small Every paragraph ends facing \textbf{your} study. If a paragraph
    could be deleted without weakening the case for your research question,
    delete it.
  \end{alertblock}
\end{center}

\begin{center}\small
\small
  \begin{columns}[T]
    \column{0.48\textwidth}
    \textbf{Summary (weak)}
    \begin{itemize}\setlength{\itemsep}{0.3em}
      \item ``Smith (2021) found A.''
      \item ``Lee (2022) found B.''
      \item ``Hansen (2023) found C.''
    \end{itemize}
    {\scriptsize\color{kgrey}A shopping list. The reader must do the
    thinking.}
    \column{0.48\textwidth}
    \textbf{Synthesis (strong)}
    \begin{itemize}\setlength{\itemsep}{0.3em}
      \item ``Survey studies consistently find A (Smith, 2021; Lee, 2022) ---
            but the single field study contradicts this (Hansen, 2023),
            suggesting the effect may not survive real workplaces.
            \textbf{This gap motivates our question.}''
    \end{itemize}
  \end{columns}
  
  \begin{alertblock}{The test}
    Every paragraph of your review should end pointing at \textbf{your}
    study: what is known, what conflicts, what is missing --- and therefore
    why your question matters.
  \end{alertblock}
\end{center}


\subsection{The toolchain for Part II}
\begin{center}\small
\footnotesize
  \renewcommand{\arraystretch}{1.35}
  \begin{tabular}{@{}p{2.8cm}p{3.4cm}p{6.6cm}@{}}
    \toprule
    \textbf{Job} & \textbf{Tool} & \textbf{Why / how}\\
    \midrule
    Collect \& cite & \textbf{Zotero} (free) & browser button saves each hit with full metadata; cite-while-you-write plugins for Word and LaTeX --- start it \emph{today}, not in November\\
    Search log \& matrix & \textbf{Excel} / Sheets & one sheet per artefact: the log (string, database, date, hits) and the concept matrix live naturally in a spreadsheet\\
    Writing & \textbf{Word} or \textbf{LaTeX} & pick one \emph{as a group}, now; Word + Zotero plugin is the low-friction default; LaTeX if the group already knows it --- happy to share templates\\
    Later & Nettskjema, JASP & data collection (S8) and analysis (S8--9) --- same rule: decide early, learn once\\
    \bottomrule
  \end{tabular}
  

  One rule beats all tool debates: \textbf{everyone in the group uses the
  same tools from session 5 on}. Merging two reference managers in submission
  week has ruined better essays than any examiner.
\end{center}


\section*{Terms from this session}
\begin{center}\small
\footnotesize
  \renewcommand{\arraystretch}{1.2}
  \begin{tabular}{@{}p{3.6cm}p{9.9cm}@{}}
    \toprule
    \textbf{Literature review} & systematic examination of existing knowledge relevant to the question --- process and product\\
    \textbf{Peer review} & pre-publication quality control by independent experts\\
    \textbf{Grey literature} & non-peer-reviewed sources (standards, industry reports) --- context, flagged as such\\
    \textbf{Boolean operators} & AND (narrow), OR (widen), NOT (exclude, carefully)\\
    \textbf{Search log} & documented record of every string, database, date and hit count\\
    \textbf{Inclusion/exclusion criteria} & pre-defined rules deciding which hits enter the review\\
    \textbf{Snowballing} & following reference lists (backward) and citations (forward) of good hits\\
    \textbf{PRISMA 2020} & reporting guideline for systematic reviews: 27-item checklist + flow diagram with two identification streams\\
    \textbf{Concept matrix} & synthesis table with concepts as columns and papers as rows (Webster \& Watson)\\
    \bottomrule
  \end{tabular}
\end{center}



\section*{Questions from the reading}

  \begin{enumerate}\setlength{\itemsep}{0.4em}
    \item Which \textbf{purposes} does a literature review serve, according to Oates?
          \answer{Two main ones: \textbf{find} a research idea and its material --- and \textbf{gather evidence} that the topic is worthwhile, the foundations are known, and the work does not merely repeat others.}
    \item What is the difference between \textbf{peer-reviewed} and \textbf{grey} literature --- and why does it matter for your essay?
          \answer{Peer-reviewed sources passed independent expert quality control and are rated higher; grey literature (reports, standards, vendor material) gives practice context --- use it, but flag it.}
    \item What must be documented so that your \textbf{search can be repeated} by others?
          \answer{Everything needed to \textbf{repeat} the search: databases, strings, dates, criteria, hit counts --- the search log.}
    \item What does Oates say about citing sources you have \textbf{not read yourself}?
          \answer{Cite only what you have read; summarising or paraphrasing still requires citing the original --- otherwise it is plagiarism, however unintended.}
  \end{enumerate}

\section*{Questions for discussion}

  \begin{enumerate}\setlength{\itemsep}{0.25em}
    \item \textbf{Your search string} --- decompose your research question into concepts and synonyms. Which concept has the most synonyms, and which database would you test first?
          \answer{The concept with the most synonyms is usually the technical one (tool names, acronyms, vendor terms); it needs the longest OR-block. Test in one database first --- Scopus or IEEE Xplore for our field --- log the hit count, then adjust.}
    \item \textbf{Inclusion and exclusion} --- write three criteria of each kind for your review. Which criterion would exclude the most papers?
          \answer{Typically the publication type (no grey literature) and the setting (only studies in organisations of a certain size or sector). State them before screening; a criterion invented after reading the hits is not a criterion but a preference.}
    \item \textbf{Synthesis, not summary} --- three papers on your topic disagree. How does that appear in Part II of your essay?
          \answer{Not as three paragraphs, one per paper, but as one argument: what they agree on, where they diverge, what explains the divergence (method, setting, definition), and what that means for your own design. The concept matrix makes the pattern visible.}
  \end{enumerate}

\section*{Key points}
\begin{itemize}
  \item The literature review does four jobs --- and follows a five-phase framework
  \item Source types differ in quality; our field has its own databases
  \item From question to keyword table to Boolean string to explicit screening --- documented, repeatable (PRISMA)
  \item Synthesis, not summary: extract into a concept matrix, then write in blocks
\end{itemize}

\section*{Sources}
\begin{itemize}\small
  \item Amling, T., et al.\ (under review). When Robotic Process Automation takes over --- do employees cheer or jeer? A systematic literature review.
  \item Enholm, I.\,M., Papagiannidis, E., Mikalef, P., \& Krogstie, J. (2022). Artificial intelligence and business value: A literature review. \emph{Information Systems Frontiers}, \emph{24}, 1709--1734. \url{https://doi.org/10.1007/s10796-021-10186-w}
  \item Kitchenham, B., \& Charters, S. (2007). \emph{Guidelines for performing systematic literature reviews in software engineering} (EBSE Technical Report EBSE-2007-01). Keele University and Durham University.
  \item Kitchenham, B., Pearl Brereton, O., Budgen, D., Turner, M., Bailey, J., \& Linkman, S. (2009). Systematic literature reviews in software engineering --- A systematic literature review. \emph{Information and Software Technology}, \emph{51}(1), 7--15.
  \item Page, M.\,J., McKenzie, J.\,E., Bossuyt, P.\,M., Boutron, I., Hoffmann, T.\,C., Mulrow, C.\,D., et al.\ (2021). The PRISMA 2020 statement: An updated guideline for reporting systematic reviews. \emph{BMJ}, \emph{372}, n71. \url{https://doi.org/10.1136/bmj.n71}
  \item vom Brocke, J., Simons, A., Niehaves, B., Riemer, K., Plattfaut, R., \& Cleven, A. (2009). Reconstructing the giant: On the importance of rigour in documenting the literature search process. In \emph{Proceedings of the 17th European Conference on Information Systems (ECIS 2009)}.
  \item Watson, R.\,T., \& Webster, J. (2020). Analysing the past to prepare for the future: Writing a literature review a roadmap for release 2.0. \emph{Journal of Decision Systems}, \emph{29}(3), 129--147. \url{https://doi.org/10.1080/12460125.2020.1798591}
  \item Webster, J., \& Watson, R.\,T. (2002). Analyzing the past to prepare for the future: Writing a literature review. \emph{MIS Quarterly}, \emph{26}(2), xiii--xxiii.
  \item Oates, B.\,J., Griffiths, M., \& McLean, R. (2022). \emph{Researching information systems and computing} (2nd ed.). SAGE. --- Chapter 6 (Reviewing the literature).
\end{itemize}

\chapter{Research Strategies 1: Survey, Design and Creation, Experiment}
\label{ch:6}
\sessionhead{6}{Thursday 15 October 2026 \;$\cdot$\; reading: Oates, Chapters 7--9}

This chapter presents the first three research strategies: the survey, which trades depth for breadth; design and creation, in which knowledge is demonstrated through an artefact and its evaluation; and the experiment, which buys causal claims with control. For each strategy the chapter gives the definition, a worked example and the criteria for judging it. The chapter closes with the choice: fit to the question, fit to the resources, the main risk --- and the strategies used in the papers read this semester.


\section{The menu}

Oates distinguishes six research strategies --- survey, design and creation, experiment, case study, action research and ethnography --- and each is defined by the kind of knowledge it produces, the data it generates and the way its data are analysed. A strategy is the overall approach to answering the question; the methods within it are the techniques of generating and analysing data (Sessions 8 and 9). This session takes the first three; Session 7 takes the other three and the decision between them.

\section{Survey}

\begin{block}{Survey}
A survey refers to the collection of the same kinds of data from a large group of people or events in a standardised, systematic way, in order to find patterns that can be generalised to a larger population than the group actually studied.
\end{block}

Four concepts make or break a survey. The population is everyone about whom the study wants to claim something; the sampling frame is the actual list from which the sample is drawn; the probability sampling technique --- random, stratified, cluster --- decides whether the results can be generalised statistically, and a convenience sample decides that they cannot; and the response rate decides how much of the intended sample is actually described. Typical questions are how widespread, which factors are associated with an outcome, and whether groups differ. The data are usually generated by questionnaires, sometimes by interviews, observations or documents, and analysed with statistics.

The four survey studies read this semester show the range of the strategy. Biloš and Budimir (2024) reached 694 respondents through an online questionnaire with seven-point Likert items; Huy et al.\ (2024) 671 respondents face to face and online; Grassini, Aasen and Møgelvang (2024) 104 Norwegian students, with a stop rule fixed in advance; and Jo and Park (2024) 351 workers recruited through a polling company's panel. All four adapted validated scales, and all four report a measurement model before the structural model. Their sampling decisions differ --- a large convenience sample, a small convenience sample with a pre-registered stop rule, a paid panel --- and each is defensible if it is stated and its limits discussed. None generalises to ``Croatian consumers'' or ``Norwegian students'' as populations.

\section{Design and creation}

\begin{block}{Design and creation}
Design and creation refers to a strategy in which new IT artefacts are developed --- constructs, models, methods or instantiations --- and the knowledge gained through building and evaluating them is the contribution.
\end{block}

The strategy is the natural home of most cyber security theses, and also their most common weakness. Building a system is development; building it, asking a question about it, evaluating it systematically and claiming something that holds beyond the prototype is research. Oates describes a five-step cycle --- awareness of the problem, suggestion, development, evaluation and conclusion --- that is iterative rather than linear, and the evaluation step is where the research lives. Without an evaluation plan the artefact is a development project. A usability test with eight users, a benchmark against a named baseline, or an expert review with explicit criteria each turns the artefact into evidence, and the plan states what is evaluated, with whom and how it is measured before the artefact is built.

Hevner et al.\ (2004) and Peffers et al.\ (2007) set out the design science research methodology that underlies the strategy in information systems; von Nordenflycht (2010) shows a related product, a conceptual paper that develops a theory and taxonomy of knowledge-intensive firms without data of its own. A conceptual paper is a study in its own right when it follows a documented method, and so is a systematic review (Session 5); neither appears in Oates' list of six strategies, and both are legitimate bases for an essay and a thesis.

\section{Experiment}

\begin{block}{Experiment}
An experiment refers to a strategy that investigates cause and effect: the researcher manipulates one factor, measures the outcome, and controls everything else, so that a difference in the outcome can be attributed to the manipulation.
\end{block}

Five terms carry the strategy: the independent variable, which is manipulated; the dependent variable, which is measured; control, so that other factors are held constant or randomised away; random assignment of participants to conditions; and the hypothesis, the predicted effect. Two kinds of validity are traded against each other. Internal validity asks whether the treatment caused the effect and is threatened by confounds and selection; external validity asks whether the result holds outside the laboratory and is threatened by artificial tasks and unrepresentative samples. More control means cleaner causal claims and less resemblance to the real world.

Sparrow, Liu and Wegner (2011) show the strategy at its best and at its most exposed. In four laboratory experiments they asked whether having information at our fingertips changes what we remember. In one experiment participants typed trivia statements into a computer and were told that the statements would be saved or erased; recall was tested afterwards. Participants who expected the computer to keep the information recalled less of it, and remembered where it was stored better than what it was. The manipulation is the treatment, recall the outcome, random assignment the comparison --- a clean causal claim. Afterwards, a large replication project of social-science experiments published in \emph{Nature} and \emph{Science} (Camerer et al., 2018) did not reproduce the main effect of the study, which is worth checking before building on it. The question has since returned with generative AI: what do we stop remembering when the assistant remembers for us?

\section{Choosing}

Three criteria decide between strategies, in this order: fit to the question, fit to the resources, and the main risk. A ``how widespread'' question fits a survey and needs access to a population; a ``does the artefact work'' question fits design and creation and needs an evaluation plan and a period of use; a ``does X cause Y'' question fits an experiment and needs control, randomisation and comparable groups. The justification paragraph of Part III has a fixed shape: we choose X because it fits the question in this way and our resources in that way; we rejected Y because of this. The rejected alternatives with reasons are what examiners read for.

The papers read this semester show which strategies our field actually uses. Four surveys, one case study (Mayer et al., 2025), one experiment from psychology (Sparrow et al., 2011), one study of documents at scale (Boodraj et al., 2026), one qualitative workshop study (Woodruff et al., 2024), one systematic review (Enholm et al., 2022) and three conceptual papers --- and not one design and creation study, the strategy most cyber security theses use. The gap is an opportunity: a well-evaluated artefact with a clear question is a contribution in a field that publishes few of them.


\section{Worked examples and tables}

\subsection{The menu}
\begin{center}\small
\small
  \renewcommand{\arraystretch}{1.3}
  \begin{tabular}{@{}p{3.3cm}p{6.9cm}p{2.6cm}@{}}
    \toprule
    \textbf{Strategy} & \textbf{Core idea} & \textbf{Session}\\
    \midrule
    \textbf{Survey} & the same questions to many, systematically & today\\
    \textbf{Design \& creation} & build an IT artefact --- and learn from building and evaluating it & today\\
    \textbf{Experiment} & manipulate a cause, measure the effect, control the rest & today\\
    Case study & one instance, studied in depth and context & session 7\\
    Action research & change something in practice, study the change & session 7\\
    Ethnography & live in the setting, understand its culture & session 7\\
    \bottomrule
  \end{tabular}
  

  A strategy is not a method: within \emph{each} strategy you still choose
  data generation methods (questionnaire, interview, observation, documents)
  --- sessions 8--9.
\end{center}


\subsection{Sampling}
\begin{center}\small
\small
  \renewcommand{\arraystretch}{1.35}
  \begin{tabular}{@{}p{3.4cm}p{10.0cm}@{}}
    \toprule
    \textbf{Population} & everyone your answer should apply to (e.g.\ SMB employees in Norway)\\
    \textbf{Sampling frame} & the actual list you can draw from (e.g.\ employees of 6 partner firms)\\
    \textbf{Probability sampling} & random selection --- allows statistical generalisation\\
    \textbf{Non-probability} & convenience, snowball, purposive --- honest label required; limits generalisation\\
    \textbf{Response rate} & who answered --- and, more dangerous, \emph{who didn't} (non-response bias)\\
    \bottomrule
  \end{tabular}
  

  \begin{alertblock}{Master's reality check}
    You will almost certainly use \textbf{convenience sampling}. That is
    acceptable --- \emph{if} you name it, justify it, and discuss the
    limitation in Part IV. Pretending otherwise is what costs marks.
  \end{alertblock}
\end{center}


\subsection{Survey --- worked example and assessment}
\begin{center}\small
\begin{columns}[T]
    \column{0.52\textwidth}
    \begin{block}{Example (cyber)}
      \textbf{RQ:} which factors are associated with password reuse among
      SMB employees?\\[0.3em]
      Online questionnaire via \textbf{Nettskjema} (UiA's survey
      tool --- data stays inside an approved system), n = 150 from six
      partner firms; validated scales for security fatigue and perceived
      account value; regression analysis.
    \end{block}
    \column{0.45\textwidth}
    \textbf{Strong at:} breadth, patterns, comparability, generalisation
    (with good sampling).
    

    \textbf{Weak at:} depth and \emph{why}; self-report bias
    (people \emph{say} they don't reuse passwords\dots); garbage in, garbage
    out --- question wording is a craft (session 8).
  \end{columns}
\end{center}


\subsection{Design and creation --- the cycle and the evaluation}
\begin{center}\small
\begin{center}
    \resizebox{\textwidth}{!}{%
    \begin{tikzpicture}[node distance=0.3cm and 0.5cm]
      \node[kbox, text width=2.6cm] (a) {\textbf{Awareness}\\ \scriptsize recognise \& articulate the problem};
      \node[kbox, text width=2.6cm, right=of a] (b) {\textbf{Suggestion}\\ \scriptsize tentative design idea};
      \node[kbox, text width=2.6cm, right=of b] (c) {\textbf{Development}\\ \scriptsize build the artefact};
      \node[kbox, text width=2.6cm, right=of c] (d) {\textbf{Evaluation}\\ \scriptsize does it work? how well? for whom?};
      \node[kbox, text width=2.6cm, right=of d] (e) {\textbf{Conclusion}\\ \scriptsize consolidate what was learned};
      \draw[karrow] (a) -- (b); \draw[karrow] (b) -- (c);
      \draw[karrow] (c) -- (d); \draw[karrow] (d) -- (e);
      \draw[garrow, dashed] (d.north) to[out=140,in=40] node[above,font=\scriptsize,kgrey]{iterate} (b.north);
    \end{tikzpicture}}
  \end{center}
  
  Not a waterfall --- the cycle is \textbf{fluid and iterative}: evaluating
  the artefact deepens your understanding of the problem, which changes the
  design, and so on.
  \hfill {\color{kgrey}\small Based on Vaishnavi \& Kuechler, in Oates et al.\ (2022), Ch.\ 8}
\end{center}

\begin{center}\small
When is building a system \textbf{research} --- and when is it just a
  development project?
  

  \small
  \renewcommand{\arraystretch}{1.15}
  \begin{tabular}{@{}p{6.2cm}p{6.6cm}@{}}
    \toprule
    \textbf{Development project} & \textbf{Design \& creation research}\\
    \midrule
    goal: a working system & goal: knowledge, demonstrated via the artefact\\
    requirements from a client & problem grounded in literature\\
    testing: does it work? & evaluation: how well, for whom, compared to what --- with a method\\
    report: what we built & report: what we \emph{learned}, transferable beyond this system\\
    \bottomrule
  \end{tabular}
  

  \begin{alertblock}{For your essay}
    \small If your group builds something, the \textbf{evaluation plan} is
    where the research lives: decide \emph{now} how the artefact will be
    evaluated (usability test? benchmark? expert review?) --- and against what.
  \end{alertblock}
\end{center}

\begin{center}\small
\begin{columns}[T]
    \column{0.52\textwidth}
    \begin{block}{Example (frontend)}
      \textbf{RQ:} can a component library enforce WCAG accessibility by
      default --- and how does it affect developer effort?\\[0.3em]
      Build the library (development); evaluate with 12 developers solving
      standard tasks with vs.\ without it: task time, WCAG violations,
      short interviews.
    \end{block}
    \column{0.45\textwidth}
    \textbf{Strong at:} relevance, motivation (you build!), natural fit with
    your master's thesis, visible product \emph{and} knowledge.
    

    \textbf{Weak at:} scope creep --- building consumes the semester and
    evaluation is left with too little time. Plan the evaluation \textbf{first}.
  \end{columns}
\end{center}

\begin{center}\small
The evaluation is where the research lives --- three common shapes:
  
  \footnotesize
  \renewcommand{\arraystretch}{1.35}
  \begin{tabular}{@{}p{2.8cm}p{5.4cm}p{4.7cm}@{}}
    \toprule
    \textbf{Shape} & \textbf{Example} & \textbf{Produces}\\
    \midrule
    Usability test & 8 users complete 5 tasks with the artefact; time, errors, SUS & user-facing quality evidence\\
    Benchmark & the new API gateway vs.\ the old one on the same load profile & performance evidence, controlled\\
    Expert review & 3 security professionals assess the threat-model notation against criteria & judgement-based evidence for constructs/methods\\
    \bottomrule
  \end{tabular}
  

  Pick the shape that matches the \textbf{claim} your question makes ---
  usable? faster? sound? --- and state it in the essay, with numbers.
\end{center}


\subsection{Experiment --- validity and a worked example}
\begin{center}\small
\begin{columns}[T]
    \column{0.48\textwidth}
    \textbf{Internal validity}\\
    \emph{Did the IV really cause the effect?}
    \begin{itemize}\setlength{\itemsep}{0.3em}
      \item threats: confounds, selection bias, learning effects, the group
            that knew it was being watched
      \item defences: randomisation, control group, counterbalancing
    \end{itemize}
    \column{0.48\textwidth}
    \textbf{External validity}\\
    \emph{Does the effect survive outside the lab?}
    \begin{itemize}\setlength{\itemsep}{0.3em}
      \item threats: artificial tasks, student-only samples, lab pressure
      \item defences: realistic tasks, field experiments, replication
    \end{itemize}
  \end{columns}
  
  \begin{alertblock}{The eternal trade-off}
    More control $\rightarrow$ cleaner causal claim $\rightarrow$ less like
    the real world. There is no free lunch --- only a documented, justified
    position on the trade-off.
  \end{alertblock}
\end{center}

\begin{center}\small
\begin{columns}[T]
    \column{0.52\textwidth}
    \begin{block}{Example (cyber)}
      \textbf{RQ:} does a redesigned browser warning reduce click-through on
      phishing links?\\[0.3em]
      240 participants, randomly assigned to old vs.\ new warning; identical
      simulated e-mail task; DV: click-through rate; hypothesis stated in
      advance; effect reported with confidence intervals.
    \end{block}
    \column{0.45\textwidth}
    \textbf{Strong at:} causal claims --- the only strategy that earns the
    word \emph{because}.
    

    \textbf{Weak at:} realism; recruiting enough participants; many
    interesting variables (breaches, team culture) simply cannot be
    manipulated ethically.
  \end{columns}
\end{center}


\subsection{Question--strategy fit}
\begin{center}\small
\small
  \renewcommand{\arraystretch}{1.35}
  \begin{tabular}{@{}p{5.6cm}p{3.2cm}p{4.2cm}@{}}
    \toprule
    \textbf{If your question sounds like\dots} & \textbf{\dots consider} & \textbf{\dots you will need}\\
    \midrule
    how widespread / which factors are associated & survey & access to enough respondents\\
    can X be built --- and how well does it work & design \& creation & build skills + an evaluation plan\\
    does A cause B & experiment & control, randomisation, participants\\
    \bottomrule
  \end{tabular}
  

  And remember session 4: if none of these fits, wait for next session --- case
  study, action research and ethnography cover the \emph{depth-in-context}
  questions.
  

  \begin{alertblock}{Essay Part III in one sentence}
    Name the strategy, show it fits the \textbf{question}, and show it fits
    your \textbf{resources} --- with the alternatives you rejected and why.
  \end{alertblock}
\end{center}

\begin{center}\small
\scriptsize
  \renewcommand{\arraystretch}{1.15}
  \begin{tabular}{@{}p{4.0cm}p{3.0cm}p{6.3cm}@{}}
    \toprule
    \textbf{Paper} & \textbf{Strategy} & \textbf{Data and analysis}\\
    \midrule
    Bilo\v{s} \& Budimir (2024) & Survey & 694 questionnaires; factor analysis, hierarchical regression\\
    Huy et al.\ (2024) & Survey (with a developmental interview phase) & 671 questionnaires; CB-SEM; 20 interviews to build the use scale\\
    Grassini et al.\ (2024) & Survey & 104 Norwegian students; PLS-SEM (session 8)\\
    Jo \& Park (2024) & Survey & 351 workers via a polling panel; PLS-SEM\\
    Sparrow et al.\ (2011) & Experiment & Four laboratory experiments; random assignment; recall as outcome\\
    Mayer et al.\ (2025) & Case study & One consultancy; semi-structured interviews (session 7)\\
    Boodraj et al.\ (2026) & Documents as data & 161{,}000 job postings analysed with Python (session 9)\\
    Woodruff et al.\ (2024) & Qualitative --- participatory workshops & 54 knowledge workers, seven industries; thematic analysis (session 9)\\
    Enholm et al.\ (2022) & Systematic literature review & 43 papers; concept matrix (session 5)\\
    von Nordenflycht (2010) & Conceptual --- theory and taxonomy & No data; argument from the literature\\
    Sapkota et al.\ (2025); Venkatesh (2022) & Conceptual review; research agenda & No data of their own\\
    \bottomrule
  \end{tabular}
  

  \small One experiment, and it comes from psychology; not one design and
  creation study --- the strategy most cyber security theses use. Two
  strategies Oates does not list, the \textbf{systematic review} and the
  \textbf{conceptual paper}, are studies in their own right when they follow a
  documented method.
\end{center}


\section*{Terms from this session}
\begin{center}\small
\footnotesize
  \renewcommand{\arraystretch}{1.18}
  \begin{tabular}{@{}p{3.6cm}p{9.9cm}@{}}
    \toprule
    \textbf{Research strategy} & the overall plan for answering the question; Oates lists six\\
    \textbf{Survey} & standardised data from many, to find generalisable patterns\\
    \textbf{Sampling frame} & the actual list respondents are drawn from\\
    \textbf{Convenience sampling} & taking who is available --- legitimate if named and its limits discussed\\
    \textbf{Design \& creation} & building an IT artefact and learning through its construction and evaluation\\
    \textbf{Artefact} & construct, model, method or instantiation produced by design \& creation\\
    \textbf{Experiment} & manipulating an independent variable and measuring the effect under control\\
    \textbf{Internal validity} & confidence that the manipulation, not something else, caused the effect\\
    \textbf{External validity} & confidence that the effect holds outside the study setting\\
    \bottomrule
  \end{tabular}
\end{center}



\section*{Questions from the reading}

  \begin{enumerate}\setlength{\itemsep}{0.4em}
    \item What distinguishes a \textbf{strategy} from a \textbf{data generation method}?
          \answer{The strategy is the \textbf{overall plan} (six on the menu); a data generation method is how the data is produced \emph{within} it --- questionnaire, interview, observation, documents.}
    \item In a survey: what is a \textbf{sampling frame}, and why does it limit your claims?
          \answer{The actual \textbf{list} the sample is drawn from --- your claims generalise to the frame's population, not beyond it.}
    \item Name the \textbf{five steps} of design \& creation.
          \answer{Awareness $\rightarrow$ suggestion $\rightarrow$ development $\rightarrow$ evaluation $\rightarrow$ conclusion --- iterative, not a waterfall.}
    \item In an experiment: what do \textbf{internal} and \textbf{external} validity mean?
          \answer{Internal: the manipulation --- not a confound --- caused the effect. External: the effect survives outside the experimental setting.}
  \end{enumerate}

\section*{Questions for discussion}

  \begin{enumerate}\setlength{\itemsep}{0.25em}
    \item \textbf{Score the three} --- survey, design and creation, experiment: for your research question, rate each on fit to question, fit to resources and main risk. Which one leads, and why?
          \answer{Write the answer as the justification paragraph: \emph{we choose X because it fits the question in this way and our resources in that way; we rejected Y because \dots}. The rejected alternatives with reasons are what examiners read for.}
    \item \textbf{If you build something} --- what is the evaluation plan, in three lines: what, with whom, measured how?
          \answer{Without an evaluation plan the artefact is a development project, not design and creation research. A usability test with eight users, a benchmark against a named baseline, or an expert review with explicit criteria each turns the artefact into evidence.}
    \item \textbf{Your included papers} --- which strategies did the studies from your search use for similar questions?
          \answer{That is evidence for your own choice: if the field answers questions like yours with surveys, choosing an experiment needs an argument; if it never uses experiments, an experiment may be a contribution in itself --- provided the controls are feasible.}
  \end{enumerate}

\section*{Key points}
\begin{itemize}
  \item Survey: breadth over depth; sampling decides what can be generalised
  \item Design and creation: knowledge demonstrated through an artefact --- the evaluation plan is where the research lives
  \item Experiment: cause and effect, bought with control; internal and external validity
  \item Choosing: fit to question, fit to resources, main risk --- and defend the choice
\end{itemize}

\section*{Sources}
\begin{itemize}\small
  \item Bilo\v{s}, A., \& Budimir, B. (2024). Understanding the adoption dynamics of ChatGPT among Generation Z: Insights from a modified UTAUT2 model. \emph{Journal of Theoretical and Applied Electronic Commerce Research}, \emph{19}(2), 863--879. \url{https://doi.org/10.3390/jtaer19020045}
  \item Enholm, I.\,M., Papagiannidis, E., Mikalef, P., \& Krogstie, J. (2022). Artificial intelligence and business value: A literature review. \emph{Information Systems Frontiers}, \emph{24}, 1709--1734. \url{https://doi.org/10.1007/s10796-021-10186-w}
  \item Grassini, S., Aasen, M.\,L., \& M\o gelvang, A. (2024). Understanding university students' acceptance of ChatGPT: Insights from the UTAUT2 model. \emph{Applied Artificial Intelligence}, \emph{38}(1), 2371168. \url{https://doi.org/10.1080/08839514.2024.2371168}
  \item Huy, L.\,V., Nguyen, H.\,T.\,T., Vo-Thanh, T., Thinh, N.\,H.\,T., \& Tran, T.\,T.\,D. (2024). Generative AI, why, how, and outcomes: A user adoption study. \emph{AIS Transactions on Human-Computer Interaction}, \emph{16}(1), 1--27. \url{https://doi.org/10.17705/1thci.00198}
  \item Jo, H., \& Park, D.-H. (2024). AI in the workplace: Examining the effects of ChatGPT on information support and knowledge acquisition. \emph{International Journal of Human--Computer Interaction}, \emph{40}(23), 8091--8106. \url{https://doi.org/10.1080/10447318.2023.2278283}
  \item von Nordenflycht, A. (2010). What is a professional service firm? Toward a theory and taxonomy of knowledge-intensive firms. \emph{Academy of Management Review}, \emph{35}(1), 155--174.
  \item Sapkota, R., Roumeliotis, K.\,I., \& Karkee, M. (2025). AI agents vs.\ agentic AI: A conceptual taxonomy, applications and challenges. \emph{Information Fusion}, \emph{126}, 103599. \url{https://doi.org/10.1016/j.inffus.2025.103599}
  \item Woodruff, A., Shelby, R., Kelley, P.\,G., Rousso-Schindler, S., Smith-Loud, J., \& Wilcox, L. (2024). How knowledge workers think generative AI will (not) transform their industries. In \emph{Proceedings of the CHI Conference on Human Factors in Computing Systems (CHI '24)}. ACM. \url{https://doi.org/10.1145/3613904.3642700}
  \item Boodraj, M., Fairbanks, J., Davies, C., \& Akcakaya, O. (2026). Beyond the firewall: The rise of the cybersecurity project manager. In \emph{AMCIS 2026 Proceedings} (Emergent Research Forum paper). AIS. \url{https://aisel.aisnet.org/amcis2026/sig_itpm/sig_itpm/2}
  \item Camerer, C.\,F., Dreber, A., Holzmeister, F., Ho, T.-H., Huber, J., Johannesson, M., et al.\ (2018). Evaluating the replicability of social science experiments in \emph{Nature} and \emph{Science} between 2010 and 2015. \emph{Nature Human Behaviour}, \emph{2}(9), 637--644. \url{https://doi.org/10.1038/s41562-018-0399-z}
  \item Mayer, A.-S., Baygi, R.\,M., \& Buwalda, R. (2025). Generation AI: Job crafting by entry-level professionals in the age of generative AI. \emph{Business \& Information Systems Engineering}, \emph{67}(5), 595--613. \url{https://doi.org/10.1007/s12599-025-00959-x}
  \item Sparrow, B., Liu, J., \& Wegner, D.\,M. (2011). Google effects on memory: Cognitive consequences of having information at our fingertips. \emph{Science}, \emph{333}(6043), 776--778. \url{https://doi.org/10.1126/science.1207745}
  \item Oates, B.\,J., Griffiths, M., \& McLean, R. (2022). \emph{Researching information systems and computing} (2nd ed.). SAGE. --- Chapters 7--9 (Surveys; Design and creation; Experiments).
  \item Laudon, K.\,C., \& Laudon, J.\,P. (2022). \emph{Management information systems: Managing the digital firm} (17th ed., Global ed.). Pearson. --- Chapter 13 (Building information systems) --- system development as practice, in contrast to design and creation as research.
\end{itemize}

\chapter{Research Strategies 2: Case Study, Action Research, Ethnography}
\label{ch:7}
\sessionhead{7}{Thursday 22 October 2026 \;$\cdot$\; reading: Oates, Chapters 10--12}

The second group of strategies studies a phenomenon in its setting: the case study, action research and ethnography. Each answers different questions and carries a different burden of justification --- above all the question of what can be claimed beyond the case. The chapter uses a current case study, an example from S\'ami research and a workshop study as illustrations, and shows how strategies can be combined, with a worked case from my own research.


\section{Case study}

\begin{block}{Case study}
A case study refers to the study of one instance of the thing being investigated --- an organisation, a department, a system, a project --- in depth, in its natural setting, using multiple sources of evidence.
\end{block}

Depth is the strategy's strength and generalisation its burden. A case study cannot generalise statistically to a population, and it does not try to; it generalises analytically, by testing or refining a theory or an explanation that can then be applied elsewhere. Every case study in an essay therefore needs one sentence that says what the case is a case \emph{of}. Cases are selected for a reason --- typical, extreme, critical, revelatory --- and the reason is stated. The data come from interviews, documents, observations and, where available, quantitative records, and their combination is what Oates calls triangulation.

Mayer, Baygi and Buwalda (2025) show a current case study in a setting close to this course. Their case was one large consultancy that had integrated generative AI into its work, chosen because the phenomenon is visible there early and in depth. Their data were semi-structured interviews, first with entry-level professionals and then with senior consultants and managers, about 45 minutes each, transcribed verbatim, with the interview guide published in an appendix. Their analysis used job crafting as the lens --- task, relational and cognitive crafting --- and a fourth form emerged from the data, signal crafting. Their generalisation is analytical: the extension of job-crafting theory travels to other knowledge-work settings; the numbers do not.

\section{Action research}

\begin{block}{Action research}
Action research refers to a strategy in which the researcher works with practitioners to change a real situation and, through the cycle of planning, acting, observing and reflecting, produces knowledge about both the change and the practice.
\end{block}

The strategy differs from design and creation in its object: design and creation produces an artefact and evaluates it; action research changes a practice and learns from the change, with the practitioners as participants rather than subjects. The cycle is iterative, the researcher is part of the situation, and the knowledge claim concerns what happened when the practice was changed in this setting and what that suggests for others. The risk is the same as the strength: involvement makes the researcher a stakeholder, and the essay must say how that was handled. A security team that introduces a new incident-triage routine with a researcher and studies what changes is doing action research; a researcher who builds a triage tool and measures its effect is doing design and creation.

\section{Ethnography}

\begin{block}{Ethnography}
Ethnography refers to the study of a culture --- a group of people and their shared practices, meanings and rules --- from the inside, through prolonged participation and observation in the group's own setting.
\end{block}

Ethnography is time-hungry and insight-rich. It shows what people actually do with technology in their own setting, security practice included, rather than what they report doing. Its data are field notes, observations, conversations and documents; its analysis is interpretive; and its claims are about meaning and practice in this setting. Digital ethnography extends the strategy to online communities, and netnography to communities that exist only online.

Coppélie Cocq's research on Sámi communities illustrates the strategy and one of its ethical lessons. Sámi communities were early adopters of digital communication --- to connect across a large territory and to be heard --- and Cocq's studies follow Sámi use of social media, forums and digital archives over years. Online and offline communities overlap so much that the ``online community'' cannot be studied as a separate, virtual place; the ethnographer's toolbox has to be adapted, with observation across channels, long engagement, and participants as co-interpreters. The ethics combine Indigenous research methodologies with digital media ethics: consent and benefit for the community, Indigenous data sovereignty, and a critical view of how data are harvested, categorised and shared (Cocq, 2022). For a cyber security thesis the lesson is that ``the data are public'' is not an argument: access is not consent, and communities have rights in their data (Session 10).

Some studies sit between the strategies and say so. Woodruff et al.\ (2024) asked how knowledge workers expect generative AI to affect their industries and answered with participatory research workshops in seven industries, 54 participants in three US cities, in which participants discussed scenarios and produced material together while the researchers analysed the discourse thematically. The dominant narrative across groups was generative AI as a tool for menial work under human review, not the disruption projected in the media, and the study names four social forces it may amplify: deskilling, dehumanisation, disconnection and disinformation. It is neither an ethnography, since there is no immersion over time, nor a case study, since there is no single bounded case; a workshop study borrows from both, and the paper says so. Name the strategy you actually use --- and its limits.

\section{Combining strategies}

Strategies can be combined when the question has two parts that need different evidence, for instance a case study to understand a practice followed by a survey to check its breadth. One strategy chosen for a clear reason beats two chosen for safety, and a combination must be justified as a combination. A combination must be justified as a combination, and the justification names what each strategy contributes that the other cannot.

\section{The decision}

All six strategies can be placed on one table with the same four columns: the question they answer, the data they generate, the knowledge they claim, and their main risk. The decision for the essay is one paragraph in Part III --- chosen strategy, fit to question, fit to resources, rejected alternatives with reasons --- and one sentence about what can be claimed beyond the study. A strategy without named data generation methods is a plan without a data plan, and Part IV cannot be written from it; Sessions 8 and 9 supply the methods.


\section{Worked examples and tables}

\subsection{Generalisation and assessment of the case study}
\begin{center}\small
\small
  \begin{columns}[T]
    \column{0.48\textwidth}
    \textbf{Statistical generalisation}\\
    from a \emph{sample} to a \emph{population}\\
    {\color{kgrey}the survey's and experiment's logic --- a case study
    (n = 1) cannot do this, and should not pretend to.}
    \column{0.48\textwidth}
    \textbf{Analytical generalisation}\\
    from a \emph{case} to a \emph{theory or concept}\\
    {\color{kgrey}``this case shows the mechanism X can occur under
    conditions Y'' --- knowledge that transfers to comparable settings.}
  \end{columns}
  
  \begin{block}{Example (cyber)}
    \textbf{RQ:} how does a small Norwegian software company recover
    operational trust after a ransomware incident?\\[0.2em]
    One company, four months: interviews with staff and management, incident
    documents, observation of the changed routines. The product is a
    \emph{mechanism description} others can learn from --- not a percentage.
  \end{block}
\end{center}

\begin{center}\small
\begin{columns}[T]
    \column{0.48\textwidth}
    \textbf{Strong at}
    \begin{itemize}\setlength{\itemsep}{0.3em}
      \item real-life complexity kept intact
      \item \emph{how} and \emph{why} questions
      \item multiple data sources, natural triangulation
      \item fits master's theses with \textbf{one client company}
    \end{itemize}
    \column{0.48\textwidth}
    \textbf{Weak at}
    \begin{itemize}\setlength{\itemsep}{0.3em}
      \item no statistical generalisation
      \item access is everything --- and can be withdrawn
      \item risk of drowning in unstructured data
      \item boundary discipline: what is \emph{in} the case, what is not?
    \end{itemize}
  \end{columns}
  
  For most of you, the client of your \textbf{master's thesis} is a natural
  case --- the strategy question is whether your research question is about
  \emph{that setting} or about something the setting merely hosts.
\end{center}

\begin{center}\small
\small
  \renewcommand{\arraystretch}{1.35}
  \begin{tabular}{@{}p{6.2cm}p{6.6cm}@{}}
    \toprule
    \textbf{Single case: one SMB after ransomware} & \textbf{Multiple: three SMBs after incidents}\\
    \midrule
    depth: four months inside one recovery & breadth: what repeats across recoveries?\\
    access: one willing gatekeeper & access: three gatekeepers --- three times the risk\\
    claim: how it happened \emph{here}, mechanism described & claim: the mechanism recurs (or varies) across settings\\
    master's-feasible: yes & feasible only with shallow data per case\\
    \bottomrule
  \end{tabular}
  

  Rule of thumb at master's scale: \textbf{one deep case beats three shallow
  ones} --- unless comparison \emph{is} the question.
\end{center}


\subsection{Action research --- example, assessment, and the common confusion}
\begin{center}\small
\begin{block}{Example (cyber)}
    \textbf{RQ:} how can alert triage in a small security team be improved,
    and what explains the improvement?\\[0.2em]
    Cycle 1: introduce a severity-tagging rule with the team; measure triage
    time and missed alerts for four weeks; reflect with the team.
    Cycle 2: adjust the rule based on what was learned; measure again.
    The knowledge product: \emph{which} intervention worked, \emph{under
    which conditions}, and \emph{why} --- documented across cycles.
  \end{block}
  
  \begin{columns}[T]
    \column{0.48\textwidth}
    \textbf{Strong at:} relevance, practitioner buy-in, studying change
    itself; natural fit when your master's thesis \emph{changes} something
    at the client.
    \column{0.48\textwidth}
    \textbf{Weak at:} the researcher is \emph{inside} the change ---
    rigour demands documented cycles and honest reflection, not success
    stories. Time: at least two cycles.
  \end{columns}
\end{center}

\begin{center}\small
\small
  \renewcommand{\arraystretch}{1.35}
  \begin{tabular}{@{}p{4.2cm}p{4.4cm}p{4.4cm}@{}}
    \toprule
    & \textbf{Design \& creation} & \textbf{Action research}\\
    \midrule
    centre of gravity & the \textbf{artefact} & the \textbf{practice / organisation}\\
    what is evaluated & how well the artefact works & how the practice changed\\
    people involved & users evaluate & practitioners \textbf{co-drive} the change\\
    typical output & artefact + evaluation knowledge & documented change across cycles\\
    \bottomrule
  \end{tabular}
  

  They can combine: build a tool (design \& creation) \emph{and} study its
  introduction into the team's practice over two cycles (action research).
  If you combine, say so --- and be clear which strategy answers which
  (sub-)question.
\end{center}


\subsection{Ethnography --- example and assessment}
\begin{center}\small
\begin{block}{Example (development culture)}
    \textbf{RQ:} how does an open-source project's community actually
    enforce code-review norms --- beyond what the written guidelines say?\\[0.2em]
    Three months inside the project: observing pull-request discussions,
    participating, field notes, informal chats with maintainers. Finding the
    \emph{unwritten} rules is exactly what only immersion can do.
  \end{block}
  
  \begin{columns}[T]
    \column{0.48\textwidth}
    \textbf{Strong at:} the taken-for-granted; unwritten rules; questions
    where asking people directly gets you the official story, not the real
    one.
    \column{0.48\textwidth}
    \textbf{Weak at:} time (usually beyond one master's semester unless the
    setting is online and access exists); ethics of observation (session 10);
    generalisation.
  \end{columns}
\end{center}

\begin{center}\small
\begin{block}{RQ: how does a security operations team really decide which alerts matter?}
    Six weeks embedded in the SOC (with written access agreement): observing
    shift handovers, sitting with analysts, field notes on the informal
    shortcuts --- ``that rule always cries wolf'', ``if it's from that
    subnet, escalate''.
  \end{block}
  
  \begin{itemize}\setlength{\itemsep}{0.35em}
    \item The written runbook says one thing; the \textbf{lived triage
          culture} does another --- only presence over time reveals the gap
    \item Contrast with session 9's methods: an interview would return the
          runbook version; observation returns practice
    \item Ethics preview: overt observation, team informed, no
          individual-level reporting --- session 10 formalises this
  \end{itemize}
\end{center}


\subsection{Combining strategies}
\begin{center}\small
\footnotesize
  \renewcommand{\arraystretch}{1.35}
  \begin{tabular}{@{}p{5.0cm}p{8.3cm}@{}}
    \toprule
    \textbf{Combination} & \textbf{Who answers what}\\
    \midrule
    Design \& creation + case study & D\&C builds and evaluates the review bot (main RQ); the case study examines its introduction into the client team's practice (SQ)\\
    Survey + interviews (mixed) & survey establishes how widespread password reuse is (SQ1); interviews explain the decisions behind it (SQ2)\\
    \bottomrule
  \end{tabular}
  

  Two rules for combinations: (1) each strategy is \textbf{assigned} to a
  (sub-)question --- no strategy without a job; (2) cost counts double ---
  two strategies means two data plans, two analyses, one deadline.
\end{center}


\subsection{Mismatch clinic and the fit of all six}
\begin{center}\small
\small
  \renewcommand{\arraystretch}{1.4}
  \begin{tabular}{@{}p{6.4cm}p{6.4cm}@{}}
    \toprule
    \textbf{Pairing} & \textbf{Why it fails}\\
    \midrule
    ``How do admins \emph{experience} audits?'' + survey & experience questions need depth; closed items flatten them --- interviews fit\\
    ``Does caching reduce latency?'' + case study & a causal claim needs manipulation and control --- experiment fits\\
    ``How widespread is MFA adoption?'' + ethnography & a prevalence question needs breadth --- survey fits; immersion is overkill\\
    \bottomrule
  \end{tabular}
  

  In each case the strategy is respectable --- it is just answering
  \textbf{someone else's question}. This is the single most common Part III
  error.
\end{center}

\begin{center}\small
\footnotesize
  \renewcommand{\arraystretch}{1.3}
  \begin{tabular}{@{}p{5.9cm}p{3.1cm}p{4.0cm}@{}}
    \toprule
    \textbf{If your question sounds like\dots} & \textbf{\dots consider} & \textbf{\dots you will need}\\
    \midrule
    how widespread / which factors are associated & survey & enough respondents\\
    can X be built --- how well does it work & design \& creation & build skills + evaluation plan\\
    does A cause B & experiment & control, randomisation\\
    how / why does it happen in \emph{this} setting & case study & deep access to one setting\\
    how can this practice be improved --- and why does it work & action research & a client willing to change; $\geq 2$ cycles\\
    what are the unwritten rules of this group & ethnography & time, immersion, access\\
    \bottomrule
  \end{tabular}
  

  Strategies can be \textbf{combined} --- but every combination must be
  justified per (sub-)question, and each extra strategy costs time you do
  not have twice.
\end{center}


\section*{Terms from this session}
\begin{center}\small
\footnotesize
  \renewcommand{\arraystretch}{1.2}
  \begin{tabular}{@{}p{3.8cm}p{9.7cm}@{}}
    \toprule
    \textbf{Case study} & in-depth study of one instance in its natural setting, multiple data sources\\
    \textbf{Analytical generalisation} & generalising from a case to a theory or mechanism, not to a population\\
    \textbf{Action research} & researcher and practitioners change practice and study the change in cycles\\
    \textbf{Plan--act--observe--reflect} & one action-research cycle; rigour requires at least two\\
    \textbf{Ethnography} & extended immersion to understand a group's culture\\
    \textbf{Thick description} & an account rich enough to convey \emph{why} things are done as they are\\
    \textbf{Virtual ethnography} & ethnography of an online setting (project, community, platform)\\
    \textbf{Reflexivity} & documenting the researcher's own influence on the setting\\
    \bottomrule
  \end{tabular}
\end{center}



\section*{Questions from the reading}

  \begin{enumerate}\setlength{\itemsep}{0.4em}
    \item What kind of \textbf{generalisation} can a case study legitimately claim?
          \answer{\textbf{Analytical} generalisation --- from the case to a theory or mechanism --- never statistical generalisation to a population.}
    \item What is the \textbf{dual aim} of action research?
          \answer{Improving the \textbf{practice} \emph{and} contributing to \textbf{knowledge} --- both, demonstrated across documented cycles.}
    \item What makes \textbf{ethnography} different from a very long case study?
          \answer{The focus on \textbf{culture} --- shared meanings and taken-for-granted norms --- reached through immersion, with the researcher as instrument and documented reflexivity.}
    \item What kind of data does an ethnographer primarily produce?
          \answer{\textbf{Field notes} and, ultimately, \textbf{thick description} --- an account rich enough to convey why things are done as they are.}
  \end{enumerate}

\section*{Questions for discussion}

  \begin{enumerate}\setlength{\itemsep}{0.25em}
    \item \textbf{The final call} --- single strategy or a justified combination? Rewrite your justification paragraph: chosen strategy, fit to question, fit to resources, rejected alternatives with reasons.
          \answer{A combination is justified when the question has two parts that need different evidence (for example a case study to understand, then a survey to check breadth). One strategy chosen for a clear reason beats two chosen for safety.}
    \item \textbf{Generalisation} --- if you choose a case study, what can you claim beyond the case?
          \answer{Not statistical generalisation to a population, but analytical generalisation: the case tests or refines a theory or explanation that can be applied elsewhere. Say which, and say what the case is a case \emph{of}.}
    \item \textbf{Data generation methods} --- which methods does your strategy imply: questionnaire, interviews, observation, documents?
          \answer{List them now; sessions 8 and 9 make them concrete. A strategy without named methods is a plan without a data plan --- Part IV of the essay cannot be written from it.}
  \end{enumerate}

\section*{Key points}
\begin{itemize}
  \item Case study: one instance in depth --- and the generalisation question it must answer
  \item Action research: change practice and generate knowledge from the change --- not design and creation
  \item Ethnography: a culture from the inside --- time-hungry, insight-rich
  \item All six strategies on one table: one decision, one justification
\end{itemize}

\section*{Sources}
\begin{itemize}\small
  \item Oates, B.\,J., Griffiths, M., \& McLean, R. (2022). \emph{Researching information systems and computing} (2nd ed.). SAGE. --- Chapters 10--12 (Case studies; Action research; Ethnography).
  \item Cocq, C. (2022). Revisiting the digital humanities through the lens of Indigenous studies --- or how to question the cultural blindness of our technologies and practices. \emph{Journal of the Association for Information Science and Technology}, \emph{73}(2), 333--344.
  \item Woodruff, A., Shelby, R., Kelley, P.\,G., Rousso-Schindler, S., Smith-Loud, J., \& Wilcox, L. (2024). How knowledge workers think generative AI will (not) transform their industries. In \emph{Proceedings of the CHI Conference on Human Factors in Computing Systems (CHI '24)}. ACM. \url{https://doi.org/10.1145/3613904.3642700}
  \item Mayer, A.-S., Baygi, R.\,M., \& Buwalda, R. (2025). Generation AI: Job crafting by entry-level professionals in the age of generative AI. \emph{Business \& Information Systems Engineering}, \emph{67}(5), 595--613. \url{https://doi.org/10.1007/s12599-025-00959-x}
  \item Laudon, K.\,C., \& Laudon, J.\,P. (2022). \emph{Management information systems: Managing the digital firm} (17th ed., Global ed.). Pearson. --- Chapter 3 (Information systems, organizations, and strategy) --- the organisational context of case studies.
\end{itemize}

\part{Data and ethics}
\chapter{Quantitative Data Collection and Analysis}
\label{ch:8}
\sessionhead{8}{Thursday 29 October 2026 \;$\cdot$\; reading: Oates, Chapters 15 and 17}

Numbers require design decisions before the first questionnaire is sent. This chapter covers the questionnaire as an instrument, the levels of measurement that decide which statistics are permitted, descriptive and inferential analysis, and the two words examiners look for: validity and reliability. Because most adoption studies in our field are surveys analysed with structural equation modelling, the chapter presents UTAUT2 as the workhorse theory, a survey study in nine steps, and the evaluation checklist for PLS-SEM.


\section{Questionnaires}

\begin{block}{Questionnaire}
A questionnaire refers to a pre-defined set of questions, assembled in a fixed order, that respondents answer without the researcher present, so that the same data are collected from every respondent in the same way.
\end{block}

Question wording is a craft with known failure modes: the double-barrelled item that asks two things at once; the leading item that suggests its answer; the ambiguous term that respondents interpret differently; the item that assumes a behaviour the respondent may not have; and the scale whose points do not match the question. Most first drafts contain at least one of these. Validated scales exist for most constructs in our field --- the UTAUT2 items are in the appendix of Venkatesh, Thong and Xu (2012) --- and reusing a validated scale is not laziness but method: the items have been tested for reliability and validity, and the results become comparable with the literature. Adaptation to a new setting is documented item by item, translation is checked by back-translation, and a pilot with ten to twenty respondents catches the ambiguities that the author cannot see.

\section{Levels of measurement}

The level at which a variable is measured decides which statistics are permitted. Nominal data are categories without order --- operating system, role --- and allow counts and modes. Ordinal data have an order without equal distances --- Likert agreement, severity ratings --- and allow medians and rank-based tests. Interval and ratio data have equal distances, and ratio data a true zero --- response time in seconds, number of incidents --- and allow means, standard deviations and the parametric tests. The level is stated per item, because it decides both the descriptive statistics and the test later on; treating a five-point Likert item as interval is common practice in our field and should be stated as a decision rather than assumed.

\section{Analysis}

Analysis begins with description: frequencies and percentages, means and medians, spread, and visual displays that show the shape of the data before any test is run. Inference follows carefully. A statistical test asks whether an observed difference or association is unlikely under the assumption of no effect; the $p$-value reports that unlikelihood, and it says nothing about the size of the effect or its practical relevance. A test that returns $p = 0.03$ therefore leaves the most important question open: how large is the effect, and does it matter in practice? Effect sizes and confidence intervals are reported with every test, because a small effect in a large sample can be significant and irrelevant at the same time. The analysis plan --- variables, descriptives, the test with its justification, the tooling --- is written before the data arrive, and it names the one comparison the study rests on. A plan written after the data arrive invites the choice of whichever test happens to be significant.

\section{Quality}

Two words summarise what examiners look for in quantitative work. Validity refers to whether the instrument measures what it claims to measure, and reliability to whether it measures consistently. Both have several faces --- content, construct and criterion validity; internal consistency and test--retest reliability --- and structural equation modelling formalises the most important of them into a checklist.

\section{Adoption surveys, UTAUT2 and PLS-SEM}

Most adoption surveys in our field build on one theory. Venkatesh, Thong and Xu (2012) extended the unified theory of acceptance and use of technology to consumers: seven constructs --- performance expectancy, effort expectancy, social influence, facilitating conditions, hedonic motivation, price value and habit --- predict behavioural intention and use, moderated by age, gender and experience. Validated with 1,512 mobile internet users in Hong Kong, the extension raised the explained variance in intention from 56 to 74 per cent. Venkatesh, Thong and Xu (2016) synthesised the theory's use into a multi-level framework --- baseline UTAUT plus individual, technology, task, location and time contexts --- and set out a research agenda. The 2012 paper supplies a validated instrument; the 2016 paper supplies a map of what has already been tested, which is where an extension has to start.

A survey study runs in nine steps: (1) theory and model, with constructs and hypotheses from a validated theory and the model drawn; (2) the instrument, items from validated scales adapted to the setting, at least three per construct; (3) translation and pilot; (4) ethics and SIKT (Session 10); (5) sampling and distribution --- a defined population, a convenience, panel or random sample, Nettskjema for distribution and storage; (6) cleaning --- incomplete responses, straight-lining, attention checks, with what was removed reported; (7) the measurement model; (8) the structural model, with effect sizes and not only $p$-values; and (9) the report, with tables for descriptives, measurement and structural results, and limitations --- sample, self-report, common method bias. Huy et al.\ (2024) document all nine steps in five pages, and their method section is a good template for Part IV of the essay.

Structural equation modelling comes in two forms, and the choice is stated in the essay. Covariance-based SEM tests theory by reproducing the covariance matrix, prefers reflective measures and established theory, and needs larger samples and near-normal data; its software is AMOS, lavaan or Mplus, and Huy et al.\ (2024) used it. Partial least squares SEM aims at prediction and exploration by maximising explained variance, handles complex models with many constructs and formative measures, works with smaller samples and non-normal data, and runs in SmartPLS; Grassini et al.\ (2024) and Jo and Park (2024) used it. Both consist of a measurement model --- do the items measure the constructs? --- and a structural model --- do the constructs relate as hypothesised? Hair, Hult, Ringle and Sarstedt (2022) set out the criteria for PLS-SEM. For the measurement model: indicator loadings of at least 0.708; Cronbach's alpha and composite reliability between 0.70 and 0.95; average variance extracted of at least 0.50; and discriminant validity with HTMT below 0.90, or 0.85 when conservative. For the structural model: variance inflation factors below three, or at most five; path coefficients with bootstrapping over 5,000 to 10,000 subsamples, with significance and confidence intervals; the $R^2$ of the endogenous constructs and $f^2$ effect sizes; and predictive relevance with $Q^2_\text{predict}$. Every one of these numbers appears in the results of Grassini et al.\ (2024) and Jo and Park (2024). A survey paper without a measurement-model table is not assessable, and neither is an essay that plans one without saying which criteria it will apply.

Two PLS-SEM studies of ChatGPT show what the checklist produces. Grassini, Aasen and Møgelvang (2024) asked which UTAUT2 constructs explain students' intention to use ChatGPT; their sample was 104 students at universities in West and Central Norway, with a stop rule set in advance; they analysed with PLS-SEM in SmartPLS 4 and bootstrapping; performance expectancy had the strongest effect, followed by habit. Jo and Park (2024) asked how ChatGPT's perceived intelligence and self-learning affect information support, knowledge acquisition and use at work; their sample was 351 workers aged 20 to 40 in Korea, recruited through a polling company's panel; perceived intelligence and self-learning led to information support and knowledge acquisition, these to utilitarian benefits, and benefits to intention, with intention rather than benefits driving use. The two sampling decisions --- a small convenience sample with a pre-registered stop rule, and a paid panel --- are both defensible if stated and discussed, and neither generalises to ``Norwegian students'' or ``Korean workers'' as populations.


\section{Worked examples and tables}

\subsection{Item wording --- failures and repairs}
\begin{center}\small
\small
  \renewcommand{\arraystretch}{1.35}
  \begin{tabular}{@{}p{3.2cm}p{10.0cm}@{}}
    \toprule
    \textbf{Failure} & \textbf{Example --- and why it breaks}\\
    \midrule
    Leading & ``Don't you agree that MFA is important?'' --- pushes the answer\\
    Double-barreled & ``Is the app fast and easy to use?'' --- fast \emph{or} easy?\\
    Jargon & ``How often do you rotate credentials?'' --- rotate what?\\
    Vague quantifiers & ``Do you often reuse passwords?'' --- often for whom?\\
    Memory overreach & ``How many phishing mails did you get last year?'' --- nobody knows\\
    \bottomrule
  \end{tabular}
  

  \begin{alertblock}{The one defence that catches all five}
    \textbf{Pilot the questionnaire} on 3--5 people who match your
    respondents --- watch them fill it in, ask where they hesitated. Report
    the pilot in Part IV; it is evidence of method quality.
  \end{alertblock}
\end{center}

\begin{center}\small
\footnotesize
  \renewcommand{\arraystretch}{1.35}
  \begin{tabular}{@{}p{6.4cm}p{6.6cm}@{}}
    \toprule
    \textbf{Broken} & \textbf{Repaired}\\
    \midrule
    Don't you agree MFA is important? & How important is MFA for your daily work? (1--5)\\
    Is the app fast and easy to use? & The app responds quickly. (1--5) \; + \; The app is easy to use. (1--5)\\
    How often do you rotate credentials? & How often do you change your work passwords?\\
    Do you often reuse passwords? & For how many of your work accounts do you use the same password? (none / some / most / all)\\
    How many phishing mails last year? & In the past month, how many suspicious e-mails do you recall receiving?\\
    \bottomrule
  \end{tabular}
  

  Each repair does one thing: removes the push, splits the pair, drops the
  jargon, anchors the quantifier, or shrinks the recall window.
\end{center}


\subsection{Norwegian practicalities and validated scales}
\begin{center}\small
\begin{itemize}\setlength{\itemsep}{0.5em}
    \item \textbf{Tool:} Nettskjema (via UiA) keeps data inside an
          approved system --- relevant the moment you collect anything
          person-related.
    \item \textbf{SIKT:} if you collect personally identifiable data (even
          e-mail addresses, even indirectly identifying combinations), the
          project may need SIKT notification --- session 10 covers this; plan
          for it \emph{now}.
    \item \textbf{Anonymity vs.\ confidentiality:} anonymous = you
          \emph{cannot} link answers to persons; confidential = you can but
          won't. Say which one you promise --- and mean it.
    \item \textbf{Response reality:} expect 10--30\,\% response rates for
          cold e-mail surveys. Plan the denominator accordingly.
  \end{itemize}
\end{center}

\begin{center}\small
\begin{itemize}\setlength{\itemsep}{0.45em}
    \item \textbf{SUS} (System Usability Scale): 10 fixed items, alternating
          direction, scored to a 0--100 usability value --- decades of
          benchmark data to compare against.
    \item \textbf{UTAUT-style acceptance items:} constructs like performance
          expectancy or effort expectancy, each measured with 3--4 Likert
          items averaged into a scale score.
    \item In your essay, reuse looks like one sentence: \emph{``Usability
          will be measured with the SUS (Brooke, 1996); security fatigue with
          the scale from [source], adapted to Norwegian.''}
  \end{itemize}
  
  \begin{alertblock}{Adaptation warning}
    Translating or trimming a validated scale changes it --- say what you
    changed, and pilot the adapted version.
  \end{alertblock}
\end{center}


\subsection{Sampling scenarios}
\begin{center}\small
\footnotesize
  \renewcommand{\arraystretch}{1.35}
  \begin{tabular}{@{}p{4.2cm}p{4.4cm}p{4.4cm}@{}}
    \toprule
    \textbf{Scenario} & \textbf{What it is} & \textbf{What you may claim}\\
    \midrule
    fellow students answer your survey & convenience sample & patterns among IT students here --- not ``employees''\\
    all 60 staff at the client & census of one organisation & solid claims about this organisation; generalise analytically, not statistically\\
    link posted in a developer forum & self-selected volunteers & exploratory signals; heavy self-selection bias --- name it\\
    \bottomrule
  \end{tabular}
  

  None of these is forbidden --- each licenses a \textbf{different
  sentence} in your conclusions. Write the sentence you can afford.
\end{center}


\subsection{Levels of measurement and descriptives}
\begin{center}\small
\small
  \renewcommand{\arraystretch}{1.35}
  \begin{tabular}{@{}p{2.1cm}p{4.5cm}p{3.1cm}p{2.9cm}@{}}
    \toprule
    \textbf{Level} & \textbf{Example} & \textbf{Allowed} & \textbf{Typical summary}\\
    \midrule
    Nominal & OS: Windows / macOS / Linux & count, compare & mode, frequencies\\
    Ordinal & Likert 1--5; severity low--high & rank & median, distribution\\
    Interval & temperature in °C & differences & mean, SD\\
    Ratio & response time in ms; count of alerts & everything, incl.\ ratios & mean, SD, ratios\\
    \bottomrule
  \end{tabular}
  

  Why it matters: the level decides the \textbf{legal operations}. A mean of
  nominal codes is nonsense; a mean of Likert items is debated --- the safe
  master's answer is: report medians and distributions for single Likert
  items, means for validated multi-item scales.
  \hfill {\color{kgrey}\small Based on Oates et al.\ (2022), Ch.\ 17}
\end{center}

\begin{center}\small
\begin{itemize}\setlength{\itemsep}{0.45em}
    \item \textbf{Centre:} mean (interval/ratio), median (ordinal, or skewed
          data), mode (nominal)
    \item \textbf{Spread:} standard deviation, range, quartiles --- a mean
          without spread hides more than it shows
    \item \textbf{Pictures before tests:} histograms and boxplots reveal
          skew, outliers and data-entry errors that summary numbers hide
  \end{itemize}
  
  \begin{block}{Worked example --- the phishing experiment (session 6)}
    Click-through: warning A 18\,\% (n = 120), warning B 9\,\% (n = 120).
    Descriptives already tell the story --- the difference is 9 percentage
    points. The question the next slide answers: could that difference be
    chance?
  \end{block}
\end{center}

\begin{center}\small
\small
  Survey (n = 84, client staff), three scale scores (1--5):
  
  \renewcommand{\arraystretch}{1.3}
  \begin{tabular}{@{}p{4.6cm}cccc@{}}
    \toprule
    \textbf{Scale} & \textbf{Mean} & \textbf{Median} & \textbf{SD} & \textbf{n}\\
    \midrule
    Security fatigue & 3.6 & 4 & 0.9 & 84\\
    Perceived account value & 2.8 & 3 & 1.2 & 84\\
    Reported reuse & 3.1 & 3 & 0.7 & 82\\
    \bottomrule
  \end{tabular}
  

  How to read it aloud: fatigue is high and fairly consistent (small SD);
  account value is mid and \textbf{divided} (large SD --- two camps?);
  two people skipped the reuse items (report the missing data, don't hide
  it). This paragraph \emph{is} descriptive analysis --- before any test.
\end{center}

\begin{center}\small
\small The same fatigue scores, plotted per group before any test:
  \begin{center}
    \begin{tikzpicture}[scale=0.8]
      % axis
      \draw[kgrey,->] (0,0) -- (10.5,0) node[below left,font=\scriptsize]{score};
      \foreach \x/\l in {1/1,3/2,5/3,7/4,9/5}{\draw[kgrey] (\x,0)--(\x,-0.12) node[below,font=\scriptsize]{\l};}
      % trained boxplot
      \draw[thick,kristiania] (3.4,1.0) rectangle (5.8,1.7);
      \draw[thick,kristiania] (4.6,1.0)--(4.6,1.7);
      \draw[thick,kristiania] (2.4,1.35)--(3.4,1.35); \draw[thick,kristiania] (5.8,1.35)--(7.0,1.35);
      \node[left,font=\scriptsize] at (2.2,1.35) {trained};
      % untrained boxplot with outlier
      \draw[thick,kgrey] (5.0,2.3) rectangle (8.2,3.0);
      \draw[thick,kgrey] (6.8,2.3)--(6.8,3.0);
      \draw[thick,kgrey] (3.8,2.65)--(5.0,2.65); \draw[thick,kgrey] (8.2,2.65)--(9.4,2.65);
      \fill[kgrey] (1.2,2.65) circle (0.07);
      \node[left,font=\scriptsize] at (2.2,3.1) {untrained};
    \end{tikzpicture}
  \end{center}
  What the picture says before any p-value: the untrained group sits higher
  \emph{and} spreads wider; one outlier (the dot) would drag a mean --- check
  whether it is a data-entry error before testing. Two minutes of plotting
  prevents the most embarrassing class of analysis mistakes.
\end{center}


\subsection{Inference, tests and effect size}
\begin{center}\small
\small
  \begin{itemize}\setlength{\itemsep}{0.4em}
    \item \textbf{The logic:} could a difference (or association) this large
          plausibly arise by chance if there were no real effect? The
          \textbf{p-value} is the probability of data at least this extreme
          \emph{under that no-effect assumption} --- nothing more.
    \item \textbf{What p is not:} not the probability the hypothesis is
          true; not a measure of importance. Always report the \textbf{effect
          size} (how big?) next to the p-value (how surprising?).
    \item \textbf{Correlation $\neq$ causation:} only the experiment's
          design --- randomisation, control --- licenses the word
          \emph{because}. A survey correlation never does.
  \end{itemize}
  
  \renewcommand{\arraystretch}{1.25}
  \begin{tabular}{@{}p{5.6cm}p{3.4cm}@{}}
    \toprule
    \textbf{Master's-feasible situation} & \textbf{Test}\\
    \midrule
    compare two group means & t-test\\
    compare frequencies across categories & chi-square\\
    association of two numeric variables & correlation\\
    \bottomrule
  \end{tabular}
  \hfill {\color{kgrey}\scriptsize Tools: JASP (free, point-and-click), Excel, R/Python}
\end{center}

\begin{center}\small
Essay Part IV wants the plan, not the results. A complete quantitative
  analysis plan states:
  
  \begin{enumerate}\setlength{\itemsep}{0.4em}
    \item \textbf{Variables:} name each one, its measurement level, and how
          it is operationalised (which items? which log field?)
    \item \textbf{Descriptives:} which summaries and figures, per variable
    \item \textbf{Tests (if any):} which comparison or association, which
          test, and \emph{why that test fits the data type}
    \item \textbf{Tooling:} the software you will use
  \end{enumerate}
  
  \begin{block}{Example line}
    ``Password-reuse (self-report, 5-point Likert, three items averaged) will
    be described with medians and distributions and compared between trained
    and untrained employees with a t-test on the scale mean; analysis in
    JASP.''
  \end{block}
\end{center}

\begin{center}\small
\footnotesize
  \renewcommand{\arraystretch}{1.35}
  \begin{tabular}{@{}p{6.6cm}p{2.4cm}p{3.8cm}@{}}
    \toprule
    \textbf{Situation} & \textbf{Test} & \textbf{Because}\\
    \midrule
    fatigue scale mean: trained vs.\ untrained employees & t-test & two groups, one interval-treated scale score\\
    clicked / didn't click $\times$ warning A / B & chi-square & two categorical variables, frequencies\\
    fatigue score vs.\ reported reuse score & correlation & two numeric variables, association not causation\\
    \bottomrule
  \end{tabular}
  

  The reasoning column is what Part IV wants --- the test name alone is a
  guess; the \emph{because} shows you matched test to data type.
\end{center}

\begin{center}\small
\small
  \renewcommand{\arraystretch}{1.4}
  \begin{tabular}{@{}p{5.4cm}p{3.2cm}p{3.6cm}@{}}
    \toprule
    \textbf{Study} & \textbf{p-value} & \textbf{Effect}\\
    \midrule
    warning redesign, n = 240 & p = .03 & 9 percentage points fewer clicks\\
    warning redesign, n = 24{,}000 & p $<$ .001 & 0.4 percentage points fewer clicks\\
    \bottomrule
  \end{tabular}
  

  The second study is \emph{more} ``significant'' and \emph{less}
  important: with enough data, trivial differences become statistically
  detectable. The decision-relevant number is the \textbf{effect size} ---
  report both, interpret the effect.
\end{center}

\begin{center}\small
\begin{columns}[T]
    \column{0.48\textwidth}
    \textbf{Validity}\\
    \emph{does the item measure what it claims?}
    \begin{itemize}\setlength{\itemsep}{0.3em}
      \item ``How secure do you feel?'' does not measure how secure anyone
            \emph{is}
      \item defence: validated instruments; pilot; triangulate self-report
            with behaviour where possible
    \end{itemize}
    \column{0.48\textwidth}
    \textbf{Reliability}\\
    \emph{would it measure the same tomorrow?}
    \begin{itemize}\setlength{\itemsep}{0.3em}
      \item ambiguous wording $\rightarrow$ answers vary randomly
      \item defence: clear items, multi-item scales, report internal
            consistency for scales you reuse
    \end{itemize}
  \end{columns}
  
  In Part IV's \textbf{limitations} discussion, name the validity threats
  you \emph{cannot} remove (self-report bias, convenience sample) --- naming
  them honestly is rewarded; hiding them is found.
\end{center}


\subsection{UTAUT2 and PLS-SEM}
\begin{center}\small
\small
  \begin{columns}[T]
    \column{0.55\textwidth}
    \textbf{UTAUT2 (2012)} --- seven constructs predict \textbf{behavioural
    intention} and \textbf{use}:
    \begin{itemize}\setlength{\itemsep}{0.1em}
      \item performance expectancy, effort expectancy, social influence, facilitating conditions
      \item hedonic motivation, price value, habit (added for consumers)
      \item moderators: age, gender, experience
    \end{itemize}
    Validated with 1{,}512 mobile internet users in Hong Kong; explained variance in intention rose from 56 to 74 per cent.
    \column{0.43\textwidth}
    \textbf{The synthesis (2016)} --- a multi-level framework: baseline UTAUT
    plus individual, technology, task, location and time contexts; a review
    of how the theory has been extended, and a research agenda.
    

    \begin{block}{\small What this gives you}
      \small A validated instrument (the items are in the appendix of the 2012
      paper) and a map of what has already been tested --- which is where your
      extension has to start.
    \end{block}
  \end{columns}
  
  {\color{uiagrey}\small Every ChatGPT study we have read this semester builds on this line: Bilo\v{s} \& Budimir, Huy et al., Grassini et al., Jo \& Park.}
\end{center}

\begin{center}\small
\small
  \renewcommand{\arraystretch}{1.3}
  \begin{tabular}{@{}p{3.0cm}p{5.0cm}p{5.0cm}@{}}
    \toprule
    & \textbf{PLS-SEM (partial least squares)} & \textbf{CB-SEM (covariance-based)}\\
    \midrule
    Aim & Prediction and exploration; maximises explained variance & Theory testing; reproduces the covariance matrix\\
    Model & Complex models, many constructs, formative measures & Reflective measures, established theory\\
    Data & Works with smaller samples and non-normal data & Needs larger samples and near-normal data\\
    Software & SmartPLS & AMOS, lavaan (R), Mplus\\
    Used in & Grassini et al.\ (2024); Jo \& Park (2024) & Huy et al.\ (2024)\\
    \bottomrule
  \end{tabular}
  

  Both belong to structural equation modelling: a \textbf{measurement model}
  (do the items measure the constructs?) and a \textbf{structural model} (do
  the constructs relate as hypothesised?). Choose by aim and data --- and say
  why in the essay.
\end{center}

\begin{center}\small
\small
  \begin{columns}[T]
    \column{0.5\textwidth}
    \textbf{Measurement model (reflective)}
    \begin{itemize}\setlength{\itemsep}{0.15em}
      \item Indicator loadings $\geq 0.708$
      \item Internal consistency: Cronbach's $\alpha$ and composite reliability between 0.70 and 0.95
      \item Convergent validity: AVE $\geq 0.50$
      \item Discriminant validity: HTMT below 0.90 (conservative: 0.85)
    \end{itemize}
    \column{0.47\textwidth}
    \textbf{Structural model}
    \begin{itemize}\setlength{\itemsep}{0.15em}
      \item Collinearity: VIF below 3 (at most 5)
      \item Path coefficients with bootstrapping (5{,}000--10{,}000 subsamples): significance and confidence intervals
      \item $R^2$ of the endogenous constructs; $f^2$ effect sizes
      \item Predictive relevance: $Q^2_\text{predict}$
    \end{itemize}
  \end{columns}
  
  \begin{alertblock}{Report all of it}
    Every number above appears in the results section of Grassini et al.\
    (2024) and Jo \& Park (2024). A survey paper without a measurement-model
    table is not assessable --- neither is an essay that plans one without
    saying which criteria it will apply.
  \end{alertblock}
\end{center}

\begin{center}\small
\footnotesize
  \renewcommand{\arraystretch}{1.3}
  \begin{tabular}{@{}p{2.2cm}p{5.3cm}p{5.5cm}@{}}
    \toprule
    & \textbf{Grassini, Aasen \& M\o gelvang (2024)} & \textbf{Jo \& Park (2024)}\\
    \midrule
    Question & Which UTAUT2 constructs explain students' intention to use ChatGPT, and their use? & How do ChatGPT's perceived intelligence and self-learning affect information support, knowledge acquisition and use at work?\\
    Sample & 104 students at universities in West and Central Norway; stop rule set in advance (date or 100 participants) & 351 workers aged 20--40, Korea, recruited through a polling company's panel\\
    Analysis & PLS-SEM in SmartPLS 4, bootstrapping & PLS-SEM; measurement then structural model\\
    Result & Performance expectancy strongest, then habit; the adapted model explains a large share of intention & Perceived intelligence and self-learning $\rightarrow$ information support and knowledge acquisition $\rightarrow$ utilitarian benefits $\rightarrow$ intention; intention, not benefits, drives use\\
    \bottomrule
  \end{tabular}
  

  \small Note the two sampling decisions: a small convenience sample with a
  pre-registered stop rule, and a paid panel. Both are defensible if they are
  stated and their limits discussed --- neither generalises to ``Norwegian
  students'' or ``Korean workers'' as populations.
\end{center}


\section*{Terms from this session}
\begin{center}\small
\footnotesize
  \renewcommand{\arraystretch}{1.18}
  \begin{tabular}{@{}p{3.8cm}p{9.7cm}@{}}
    \toprule
    \textbf{Likert scale} & ordered response scale (e.g.\ 1--5 agreement); single items are ordinal\\
    \textbf{Validated instrument} & published, tested item set for a construct (SUS, UTAUT items)\\
    \textbf{Pilot} & trial run of the questionnaire on a few matching respondents\\
    \textbf{Measurement level} & nominal / ordinal / interval / ratio --- decides legal operations\\
    \textbf{Descriptive statistics} & summaries of the data at hand: centre, spread, distributions\\
    \textbf{p-value} & probability of data at least this extreme, assuming no real effect\\
    \textbf{Effect size} & how large the difference or association is --- report it with p\\
    \textbf{Validity (instrument)} & the item measures what it claims to measure\\
    \textbf{Reliability (instrument)} & the item measures consistently\\
    \bottomrule
  \end{tabular}
\end{center}



\section*{Questions from the reading}

  \begin{enumerate}\setlength{\itemsep}{0.4em}
    \item According to Oates, what makes a \textbf{good questionnaire question}?
          \answer{Five criteria (Peterson, in Ch.\ 15): \textbf{brief, relevant, unambiguous, specific, objective} --- our five failure modes are these, inverted.}
    \item What are the \textbf{four levels of measurement} --- and why do they matter?
          \answer{Nominal, ordinal, interval, ratio --- the level decides which operations and which statistics are legal on the data.}
    \item Why should you report \textbf{spread}, not just the mean?
          \answer{Two very different distributions can share a mean --- spread (SD, quartiles) shows whether respondents agree or are split.}
    \item What does a \textbf{p-value} tell you --- and what does it \emph{not} tell you?
          \answer{p = the probability of data at least this extreme \emph{if there were no real effect}. It is not the probability the hypothesis is true, and says nothing about the size or importance of the effect.}
  \end{enumerate}

\section*{Questions for discussion}

  \begin{enumerate}\setlength{\itemsep}{0.25em}
    \item \textbf{Your items} --- draft three questionnaire items for your research question and state each item's level of measurement. Check them against the five failure modes.
          \answer{Most first drafts contain a double-barrelled item or a leading one. State the level of measurement per item; it decides which descriptive statistics and which test are permitted later.}
    \item \textbf{The analysis plan} --- variables, descriptives, test (if any) with justification, tooling: write it before the data arrive. What is the one comparison your study rests on?
          \answer{Naming the one comparison forces the design: two groups or one, paired or independent, which measure. A plan written after the data arrive invites the choice of whichever test happens to be significant.}
    \item \textbf{Effect size} --- your test returns $p = 0.03$. What do you still not know?
          \answer{How large the effect is and whether it matters in practice. Report the effect size and a confidence interval with every test; a small effect in a large sample can be significant and irrelevant at the same time.}
  \end{enumerate}

\section*{Key points}
\begin{itemize}
  \item Questionnaires: question wording is a craft; validated scales exist --- UTAUT2 is the workhorse
  \item Levels of measurement decide what you may compute
  \item A survey study in nine steps; PLS-SEM or CB-SEM, chosen by aim and data
  \item The PLS-SEM checklist: loadings, reliability, AVE, HTMT; VIF, bootstrapped paths, $R^2$, $f^2$, $Q^2$
\end{itemize}

\section*{Sources}
\begin{itemize}\small
  \item Grassini, S., Aasen, M.\,L., \& M\o gelvang, A. (2024). Understanding university students' acceptance of ChatGPT: Insights from the UTAUT2 model. \emph{Applied Artificial Intelligence}, \emph{38}(1), 2371168. \url{https://doi.org/10.1080/08839514.2024.2371168}
  \item Hair, J.\,F., Hult, G.\,T.\,M., Ringle, C.\,M., \& Sarstedt, M. (2022). \emph{A primer on partial least squares structural equation modeling (PLS-SEM)} (3rd ed.). SAGE. --- Chapter 1.
  \item Huy, L.\,V., Nguyen, H.\,T.\,T., Vo-Thanh, T., Thinh, N.\,H.\,T., \& Tran, T.\,T.\,D. (2024). Generative AI, why, how, and outcomes: A user adoption study. \emph{AIS Transactions on Human-Computer Interaction}, \emph{16}(1), 1--27. \url{https://doi.org/10.17705/1thci.00198}
  \item Jo, H., \& Park, D.-H. (2024). AI in the workplace: Examining the effects of ChatGPT on information support and knowledge acquisition. \emph{International Journal of Human--Computer Interaction}, \emph{40}(23), 8091--8106. \url{https://doi.org/10.1080/10447318.2023.2278283}
  \item Venkatesh, V., Thong, J.\,Y.\,L., \& Xu, X. (2012). Consumer acceptance and use of information technology: Extending the unified theory of acceptance and use of technology. \emph{MIS Quarterly}, \emph{36}(1), 157--178. \url{https://doi.org/10.2307/41410412}
  \item Venkatesh, V., Thong, J.\,Y.\,L., \& Xu, X. (2016). Unified theory of acceptance and use of technology: A synthesis and the road ahead. \emph{Journal of the Association for Information Systems}, \emph{17}(5), 328--376. \url{https://doi.org/10.17705/1jais.00428}
  \item Oates, B.\,J., Griffiths, M., \& McLean, R. (2022). \emph{Researching information systems and computing} (2nd ed.). SAGE. --- Chapters 15 (Questionnaires) and 17 (Quantitative data analysis).
\end{itemize}

\chapter{Qualitative Data Collection and Analysis}
\label{ch:9}
\sessionhead{9}{Thursday 5 November 2026 \;$\cdot$\; reading: Oates, Chapters 13--14, 16 and 18}

Qualitative data --- interviews, observations, documents --- demand the same rigour as numbers, with different tools. This chapter covers the interview guide, systematic and participant observation, documents as data, and thematic analysis from transcript to theme. The argument for the number of participants is made with information power rather than saturation, and published studies show what a defensible qualitative data plan looks like.


\section{Interviews}

\begin{block}{Interview}
An interview refers to a conversation between researcher and participant that has a purpose, an agenda and a documented guide, and in which the researcher generates data by asking and listening.
\end{block}

Interviews come in three degrees of structure. A structured interview asks the same fixed questions of everyone and is close to a questionnaire read aloud; a semi-structured interview follows a guide of open questions with room for probes and reordering; an unstructured interview starts from a topic and follows the participant. Most master's theses use the semi-structured form, and the guide is its instrument: an opening that settles the participant, two or three core questions that address the research question, probes that follow the material --- ``what happened then?'' --- and a closing that asks what was missed. Good questions are open, concrete and not leading: ``tell me about the last time you \dots'' asks for an episode and produces data, whereas ``do you think X is important?'' produces a yes. The practicalities are known and still neglected: a pre-tested guide, consent obtained before recording, a quiet setting, audio recording and verbatim transcription, and a note after each interview on what surprised the interviewer.

Mayer, Baygi and Buwalda (2025) show what a published interview study reports. They interviewed entry-level professionals first, and then senior consultants and managers once the first interviews raised questions about how seniors perceive the change --- the sample followed the analysis. The interviews were semi-structured, about 45 minutes, audio-recorded and transcribed verbatim, and the guide is published in the appendix. Three known forms of job crafting came out of the data and one new form, signal crafting, which the guide had not asked for. The parts to copy are two: the sampling decision is explained as a consequence of the analysis, and the guide is published. Both belong in Part IV of the essay.

\section{Observations and documents}

\begin{block}{Observation}
Observation refers to watching and recording what people actually do, as opposed to what they say they do, either systematically with a predefined schedule or as a participant in the setting.
\end{block}

Systematic observation uses a schedule of categories and counts occurrences; participant observation immerses the researcher in the setting and records field notes. Both raise the question of the observer's effect, which the essay addresses rather than ignores. Documents are the third source: policies, incident reports, meeting minutes, logs, job postings and software repositories are data that already exist, and the method questions are the same as for interviews --- which corpus, why, how selected and how coded. Boodraj et al.\ (2026) show documents at scale: they asked which skills and competencies are most in demand for cyber security project managers, took 5,841 cyber security project manager and 155,436 IT project manager job postings from 2014 to 2025 from a labour market database, analysed them with Python-based text analytics, and plan a market basket analysis of co-occurring skills; their preliminary finding is that employers ask for dual-domain certifications and for interpersonal and team skills.

\section{Analysis}

\begin{block}{Thematic analysis}
Thematic analysis refers to a method for identifying, analysing and reporting patterns --- themes --- in qualitative data, through a documented sequence of familiarisation, coding, theme development, review and definition.
\end{block}

The sequence runs from the transcript to the theme, and every step is visible. Familiarisation means reading the material more than once; coding means labelling segments with what they are about; themes group codes into patterns that answer the research question; reviewing tests the themes against the whole data set; and defining names each theme and states what it captures. Two people who code the same transcript independently will disagree on a third of the codes, and that is the normal starting point rather than a problem: the disagreements are discussed, the code book refined, and the procedure reported --- which is what quality in qualitative analysis looks like, not a percentage. Software such as NVivo supports the bookkeeping; it does not do the analysis.

\section{The qualitative data plan}

Part IV of the essay asks for a data plan with five elements: participants, method, handling, analysis and limitations. The number of participants is argued rather than asserted. ``Saturation'' is tied to grounded theory and used inconsistently; Malterud, Siersma and Guassora (2016) propose information power instead. The more relevant information the sample holds, the fewer participants are needed, and five dimensions decide it: a narrow study aim needs fewer participants than a broad one; a dense sample, whose participants hold exactly the experience sought, needs fewer than a sparse one; a study that applies established theory needs fewer than one that does not; a strong dialogue --- clear, focused, articulate --- needs fewer than a weak one; and case analysis in depth needs fewer than cross-case analysis. ``Eight interviews with SOC analysts, a narrow aim, a dense sample and UTAUT as a lens'' is an argument; ``until saturation'' is not, and the argument is revisited during data collection.

The honest limitations of most qualitative theses are access and self-report: a small purposive sample from one organisation, and accounts of behaviour rather than behaviour. Naming them is judgement; hiding them is a weakness.


\section{Worked examples and tables}

\subsection{The interview guide, practicalities and probing}
\begin{center}\small
For the question \emph{``what does `secure enough' mean to sysadmins in
  small businesses?''}:
  
  \small
  \begin{enumerate}\setlength{\itemsep}{0.35em}
    \item \textbf{Opening (easy, factual):} ``Can you walk me through a
          normal day in your role?''
    \item \textbf{Core (open, experience-near):} ``Tell me about the last
          time you decided something was secure enough to ship.''
    \item \textbf{Probe (follow the answer):} ``You said `good enough for
          us' --- what makes it good enough \emph{here}?''
    \item \textbf{Contrast:} ``Was there a time you were overruled on a
          security decision? What happened?''
    \item \textbf{Closing:} ``Is there anything about this topic I should
          have asked and didn't?''
  \end{enumerate}
  
  Rules visible in the example: \textbf{open} questions, \textbf{concrete
  episodes} (not opinions in general), \textbf{no leading}, probes prepared
  but flexible.
\end{center}

\begin{center}\small
\begin{itemize}\setlength{\itemsep}{0.45em}
    \item \textbf{How many?} Until \textbf{saturation}: new interviews stop
          producing new themes. At master's scale, 6--12 semi-structured
          interviews of 30--60 minutes is a defensible plan --- \emph{state
          the saturation logic}, don't just state a number.
    \item \textbf{Recording:} always ask consent, always record (notes lie);
          audio only. Use the \textbf{Nettskjema diktafon app} --- audio goes
          straight to approved storage instead of your phone's cloud backup
          (session 10 explains why that matters).
    \item \textbf{Transcription:} verbatim for the parts you analyse; budget
          4--6 hours of transcription per interview hour --- this is where
          qualitative time actually goes.
    \item \textbf{Pilot} the guide once, exactly like a questionnaire.
  \end{itemize}
\end{center}

\begin{center}\small
\footnotesize
  \renewcommand{\arraystretch}{1.35}
  \begin{tabular}{@{}p{6.3cm}p{6.7cm}@{}}
    \toprule
    \textbf{Broken} & \textbf{Repaired}\\
    \midrule
    Don't you find security policies annoying? & How do the security policies affect your daily work?\\
    Do you follow the password policy? (yes/no) & Tell me about the last time the password policy got in your way.\\
    Why is your team bad at reviews? & How do reviews typically go in your team?\\
    What do you think about security in general? & Describe a recent security decision you were part of.\\
    Would you say training helps? & What, if anything, changed after the last training?\\
    \bottomrule
  \end{tabular}
  

  The repairs share a pattern: \textbf{open} instead of yes/no,
  \textbf{episodes} instead of opinions, \textbf{neutral} instead of loaded.
\end{center}

\begin{center}\small
\small
  \begin{itemize}\setlength{\itemsep}{0.45em}
    \item[] \textbf{P:} ``We just click through the warnings, honestly.''
    \item[] \textbf{Weak follow-up:} ``So you ignore security?''
          {\color{kgrey}(judgement --- shuts the door)}
    \item[] \textbf{Probe 1 --- elaborate:} ``Walk me through the last time
          that happened.''
    \item[] \textbf{Probe 2 --- clarify:} ``When you say \emph{we} --- is
          that the whole team?''
    \item[] \textbf{Probe 3 --- contrast:} ``Was there ever a warning you
          didn't click through? What was different?''
  \end{itemize}
  
  Probes are prepared \textbf{types}, not prepared sentences: elaborate,
  clarify, contrast, and --- hardest --- \textbf{silence}. Two seconds of it
  produces more data than any question.
\end{center}


\subsection{Observations, field notes and documents}
\begin{center}\small
\begin{block}{Observation}
    Systematically watching and recording what people \textbf{actually do}
    in the setting --- overt or covert, participant or non-participant ---
    captured in structured \textbf{field notes}.
  \end{block}
  
  \begin{block}{Documents}
    Existing materials analysed as data: \textbf{found documents} (incident
    reports, policies, commit messages, pull-request threads, meeting
    minutes) or \textbf{researcher-generated} ones (diaries you ask
    participants to keep).
  \end{block}
  
  Why they matter: interviews get the \emph{account}; observation gets the
  \emph{behaviour}; documents get the \emph{trace}. Where the three
  disagree, the interesting finding usually lives.
  \hfill {\color{kgrey}\small Oates et al.\ (2022), Ch.\ 14, 16}
\end{center}

\begin{center}\small
\begin{block}{RQ: how does a development team actually handle security warnings?}
    \begin{itemize}\setlength{\itemsep}{0.35em}
      \item \textbf{Interviews:} developers \emph{say} they review every
            warning before merging.
      \item \textbf{Documents:} the CI logs and pull requests show warnings
            suppressed with \texttt{ignore} flags in a third of merges.
      \item \textbf{Observation:} in sprint pressure, warnings are triaged
            in seconds --- by habit, not policy.
    \end{itemize}
  \end{block}
  
  The finding is the \textbf{gap between account, trace and behaviour} ---
  reachable only because three methods were combined. One method would have
  reported either compliance or cynicism; the combination explains
  \emph{when} each is true.
\end{center}

\begin{center}\small
\small
  \renewcommand{\arraystretch}{1.4}
  \begin{tabular}{@{}p{6.4cm}p{6.4cm}@{}}
    \toprule
    \textbf{Descriptive (what happened)} & \textbf{Interpretive (what I make of it)}\\
    \midrule
    10:14 --- build fails on security scan; A groans, adds \texttt{ignore} flag without opening the finding; nobody comments & suppression is routinised --- no discussion needed, so the norm is settled\\
    10:40 --- B asks A ``real one?''; A shrugs: ``same as always'' & a shared classification exists (``real ones'' vs.\ noise) that no document defines\\
    \bottomrule
  \end{tabular}
  

  Keep the two columns \textbf{separate}: the left is evidence, the right is
  analysis-in-progress. Mixing them is how observations quietly become
  opinions.
\end{center}

\begin{center}\small
\begin{block}{Corpus: 6 months of pull requests in the client's main repository}
    Extract per PR: review turnaround, number of comments, whether the
    security checklist was mentioned, who approved.
  \end{block}
  
  \begin{itemize}\setlength{\itemsep}{0.4em}
    \item \textbf{Found} documents: produced for work, not for you --- so
          they show behaviour without observer effects
    \item But they answer only what they record: turnaround yes, the
          \emph{reasoning} no --- pair them with interviews
    \item Practicalities: written permission for repository access; strip
          author names before analysis; note the corpus boundary (which
          repos, which period) in Part IV
  \end{itemize}
\end{center}


\subsection{Thematic analysis --- process, coding, quality}
\begin{center}\small
\begin{center}
    \resizebox{0.88\textwidth}{!}{%
    \begin{tikzpicture}
      \node[gbox, text width=2.5cm, minimum height=1.2cm, anchor=west] (fam) at (0,0)
        {\textbf{1.\ Familiarise}\\ \scriptsize read all transcripts, first impressions};
      % the cycle: 2 -> 3 -> 4 -> 2
      \node[kbox, text width=2.5cm, minimum height=1.1cm] (code) at (5.6,1.5)
        {\textbf{2.\ Code}\\ \scriptsize label segments, close to the words};
      \node[kbox, text width=2.5cm, minimum height=1.1cm] (thm) at (8.6,-0.2)
        {\textbf{3.\ Build themes}\\ \scriptsize group codes};
      \node[kbox, text width=2.5cm, minimum height=1.1cm] (rev) at (5.6,-1.9)
        {\textbf{4.\ Review}\\ \scriptsize check themes against \emph{all} data};
      \node[gbox, text width=2.6cm, minimum height=1.2cm, anchor=west] (wr) at (11.6,-0.2)
        {\textbf{5.\ Define \& write}\\ \scriptsize name themes, evidence with quotes};
      \draw[karrow] (fam.east) -- (code.west);
      \draw[karrow] (code.east) to[bend left=25] (thm.north);
      \draw[karrow] (thm.south) to[bend left=25] (rev.east);
      \draw[karrow] (rev.west) to[bend left=30]
        node[left,font=\scriptsize,kristiania]{recode} (code.west);
      \draw[garrow] (thm.east) -- node[above,font=\scriptsize,kgrey]{when stable} (wr.west);
      \node[font=\scriptsize\color{kgrey}] at (6.7,-0.2) {cycle 2--4};
    \end{tikzpicture}}
  \end{center}
  
  Steps 2--4 are a \textbf{cycle}, not a sequence: reviewing themes sends you
  back to recode, often several times. You exit to writing only when new
  passes stop changing the themes. Inductive by default: the codes come
  \textbf{from the data}, not from a list made in advance.
  \hfill {\color{kgrey}\small Based on Oates et al.\ (2022), Ch.\ 18}
\end{center}

\begin{center}\small
\small
  \renewcommand{\arraystretch}{1.4}
  \begin{tabular}{@{}p{7.4cm}p{2.7cm}p{2.6cm}@{}}
    \toprule
    \textbf{Transcript (sysadmin, small business)} & \textbf{Code} & \textbf{Theme}\\
    \midrule
    ``Honestly, I patch when something breaks. There's no time otherwise.'' & reactive patching & security as time negotiation\\
    ``The boss asks what security costs, never what it saves.'' & cost framing by management & security as time negotiation\\
    ``After the incident next door, suddenly everything was urgent.'' & incident as trigger & external events drive priority\\
    \bottomrule
  \end{tabular}
  

  Notice: codes stay \textbf{close to the words}; themes \textbf{interpret}
  across codes. In the write-up, every theme is evidenced with 1--2
  \textbf{verbatim quotes} --- the qualitative equivalent of reporting your
  numbers.
\end{center}

\begin{center}\small
\small
  Interviews with frontend developers on accessibility work:
  
  \footnotesize
  \renewcommand{\arraystretch}{1.3}
  \begin{tabular}{@{}p{5.2cm}p{3.4cm}p{4.0cm}@{}}
    \toprule
    \textbf{Codes (close to the words)} & \textbf{Grouped} & \textbf{Theme}\\
    \midrule
    ``no ticket, no time''; ``sprint has no a11y column'' & work not planned & accessibility is invisible work\\
    ``I learned it from Maria''; ``one person owns it'' & knowledge in persons & expertise is a single point of failure\\
    ``audit panic in March''; ``then it faded'' & event-driven attention & compliance is episodic, not continuous\\
    \bottomrule
  \end{tabular}
  

  Note the movement left to right: quotes $\rightarrow$ codes $\rightarrow$
  interpretation. Every theme can be traced back to words someone actually
  said --- that traceability is your quality argument.
\end{center}

\begin{center}\small
\begin{itemize}\setlength{\itemsep}{0.45em}
    \item \textbf{Transparency:} document the coding process --- how codes
          arose, how themes were built. A small \textbf{code book} (code,
          definition, example quote) in the appendix does this.
    \item \textbf{Second coder:} two people code a subset independently and
          discuss to consensus --- the procedure used in the review you saw
          in session 5. At master's scale, even coding \emph{two} transcripts in
          pairs strengthens the method section.
    \item \textbf{Evidence, not anecdote:} themes carried by quotes from
          \emph{several} participants; note disconfirming cases instead of
          hiding them.
    \item \textbf{Tools:} a table in Word does the job at 6--12 interviews;
          dedicated software (e.g.\ NVivo) only pays off at larger scale.
  \end{itemize}
\end{center}


\subsection{Information power and the data plan}
\begin{center}\small
\small ``Saturation'' is tied to grounded theory and is used inconsistently.
  Malterud et al.\ propose \textbf{information power}: the more relevant
  information the sample holds, the fewer participants are needed. Five
  dimensions decide it:
  
  \footnotesize
  \renewcommand{\arraystretch}{1.15}
  \begin{tabular}{@{}p{3.4cm}p{4.6cm}p{4.6cm}@{}}
    \toprule
    \textbf{Dimension} & \textbf{Fewer participants} & \textbf{More participants}\\
    \midrule
    Study aim & Narrow & Broad\\
    Sample specificity & Dense: participants hold exactly the experience sought & Sparse\\
    Established theory & Applied & Not applied\\
    Quality of dialogue & Strong: clear, focused, articulate & Weak\\
    Analysis strategy & Case analysis (in-depth, few) & Cross-case analysis (many)\\
    \bottomrule
  \end{tabular}
  

  \small For your essay: argue the sample size with these five, before data
  collection, and revisit it during collection. ``Eight interviews with SOC
  analysts, a narrow aim, a dense sample and UTAUT as a lens'' is an argument;
  ``until saturation'' is not.
\end{center}

\begin{center}\small
\small
  \begin{enumerate}\setlength{\itemsep}{0.3em}
    \item \textbf{Participants:} who, how recruited, how many --- with the
          saturation logic stated
    \item \textbf{Method:} interview type and guide (appendix), or
          observation protocol, or document corpus and access
    \item \textbf{Data handling:} recording, transcription, storage ---
          consistent with consent and SIKT (session 10)
    \item \textbf{Analysis:} thematic analysis, the five steps, coding
          approach, quality measures (code book, second coder)
    \item \textbf{Limitations:} researcher influence, sample scope,
          self-report vs.\ behaviour --- named, not hidden
  \end{enumerate}
  \begin{block}{Example line}
    \small ``Eight semi-structured interviews (30--45 min) with sysadmins
    recruited via the client's network, recorded and transcribed verbatim,
    analysed thematically with an inductively developed code book; two
    transcripts double-coded.''
  \end{block}
\end{center}


\section*{Terms from this session}
\begin{center}\small
\footnotesize
  \renewcommand{\arraystretch}{1.18}
  \begin{tabular}{@{}p{3.8cm}p{9.7cm}@{}}
    \toprule
    \textbf{Semi-structured interview} & interview guide plus freedom to follow up --- the master's default\\
    \textbf{Interview guide} & the planned questions and probes; goes in the appendix\\
    \textbf{Saturation} & the point where new data stops producing new themes\\
    \textbf{Field notes} & structured written record of observations\\
    \textbf{Found documents} & existing materials analysed as data (logs, policies, PR threads)\\
    \textbf{Code} & a short label for a meaningful data segment, close to the words\\
    \textbf{Theme} & an interpretation that groups codes across the data\\
    \textbf{Code book} & table of codes with definitions and example quotes\\
    \textbf{Second coder} & independent coding of a subset, compared to consensus\\
    \bottomrule
  \end{tabular}
\end{center}



\section*{Questions from the reading}

  \begin{enumerate}\setlength{\itemsep}{0.4em}
    \item \textbf{Structured} vs.\ \textbf{semi-structured} interviews --- what exactly changes?
          \answer{Structured: fixed questions in fixed order --- comparability. Semi-structured: a guide, but order, omissions and follow-ups are free --- focus \emph{plus} discovery.}
    \item What belongs in good \textbf{field notes}?
          \answer{What was observed --- when, where, who, what happened --- kept \textbf{separate} from your interpretations, which are noted but marked as such.}
    \item \textbf{Found} vs.\ \textbf{researcher-generated} documents --- one example of each?
          \answer{Found: exist independently of the research (policies, logs, PR threads). Researcher-generated: produced \emph{for} the research (e.g.\ diaries you ask participants to keep).}
    \item What are the steps that lead from a \textbf{transcript} to a \textbf{theme}?
          \answer{Familiarise with the transcripts $\rightarrow$ code close to the words $\rightarrow$ group codes into candidate themes $\rightarrow$ review against all data $\rightarrow$ define and write, evidenced with quotes.}
  \end{enumerate}

\section*{Questions for discussion}

  \begin{enumerate}\setlength{\itemsep}{0.25em}
    \item \textbf{Your interview guide} --- opening, two core questions, one probe, closing: draft it for your research question. Is each question open, concrete and not leading?
          \answer{``Tell me about the last time you \dots'' asks for an episode and produces data; ``Do you think X is important?'' produces a yes. Probes (``what happened then?'') are where the material usually is.}
    \item \textbf{Coding a fragment} --- two people code the same transcript independently. They disagree on a third of the codes. Is that a problem?
          \answer{It is the normal starting point and the reason for a second coder. Discuss the disagreements, refine the code book, and report the procedure --- that is what quality in qualitative analysis looks like, not a percentage.}
    \item \textbf{Your qualitative data plan} --- participants, method, handling, analysis, limitations: write the paragraph. Which limitation is the honest one?
          \answer{Usually access and self-report: a small purposive sample from one organisation, and accounts of behaviour rather than behaviour. Name them; naming a limitation is judgement, hiding it is a weakness.}
  \end{enumerate}

\section*{Key points}
\begin{itemize}
  \item Interviews: a conversation with a purpose and a documented guide
  \item Observations and documents: what people do, and what their work leaves behind
  \item Thematic analysis: from transcript to theme, every step visible
  \item The qualitative data plan --- the same rigour as the numbers
\end{itemize}

\section*{Sources}
\begin{itemize}\small
  \item Malterud, K., Siersma, V.\,D., \& Guassora, A.\,D. (2016). Sample size in qualitative interview studies: Guided by information power. \emph{Qualitative Health Research}, \emph{26}(13), 1753--1760. \url{https://doi.org/10.1177/1049732315617444}
  \item Woodruff, A., Shelby, R., Kelley, P.\,G., Rousso-Schindler, S., Smith-Loud, J., \& Wilcox, L. (2024). How knowledge workers think generative AI will (not) transform their industries. In \emph{Proceedings of the CHI Conference on Human Factors in Computing Systems (CHI '24)}. ACM. \url{https://doi.org/10.1145/3613904.3642700}
  \item Boodraj, M., Fairbanks, J., Davies, C., \& Akcakaya, O. (2026). Beyond the firewall: The rise of the cybersecurity project manager. In \emph{AMCIS 2026 Proceedings} (Emergent Research Forum paper). AIS. \url{https://aisel.aisnet.org/amcis2026/sig_itpm/sig_itpm/2}
  \item Mayer, A.-S., Baygi, R.\,M., \& Buwalda, R. (2025). Generation AI: Job crafting by entry-level professionals in the age of generative AI. \emph{Business \& Information Systems Engineering}, \emph{67}(5), 595--613. \url{https://doi.org/10.1007/s12599-025-00959-x}
  \item Oates, B.\,J., Griffiths, M., \& McLean, R. (2022). \emph{Researching information systems and computing} (2nd ed.). SAGE. --- Chapters 13--14, 16 and 18 (Interviews; Observations; Documents; Qualitative data analysis).
\end{itemize}

\chapter{Ethics in Research}
\label{ch:10}
\sessionhead{10}{Thursday 12 November 2026 \;$\cdot$\; reading: Oates, Chapters 4--5}

Research ethics is not paperwork: it is the part of the design that decides whether the study may be done at all, and how. This chapter covers personal data and the law, informed consent, anonymity versus confidentiality, the SIKT notification and the data management plan, and three special zones --- public data, security research and AI tools. It ends with the references every student should keep at hand, including the University of Agder's own guidelines.


\section{The law}

\begin{block}{Personal data}
Personal data refers to any information relating to an identified or identifiable person --- directly, such as a name or e-mail address, or indirectly, such as an IP address, a voice recording, or a job title in a small organisation.
\end{block}

Personal data is broader than it looks, and the law obliges the researcher. E-mail addresses, IP addresses, voice recordings and free-text answers are personal data; so are job titles in a small organisation, because they identify indirectly. The General Data Protection Regulation and the Norwegian Personal Data Act set the frame: a legal basis for processing, information to the persons concerned, minimisation, secure storage, and deletion when the purpose is fulfilled. If anything on the data inventory of a study is personal data, the SIKT notification follows.

\section{Participants' rights}

\begin{block}{Informed consent}
Informed consent refers to a participant's agreement to take part, given freely after being told what the study is for, what participation involves, how data are stored and deleted, that withdrawal is possible at any time, and whom to contact.
\end{block}

Consent is documented, usually with a form that states purpose, procedure, storage and deletion, the right to withdraw, and contact details. Two promises are often confused. Confidentiality means that the researcher knows who said what and does not reveal it; anonymity means that no one, including the researcher, can link the data to a person. Confidentiality can be promised; anonymity only if the data will never be linkable --- which is rarely true of interviews, and never true of a recorded voice. The client problem follows: examiners sign no non-disclosure agreement, and client-confidential details, especially vulnerabilities, do not belong in the graded essay. Anonymisation is planned before data collection, not after.

\section{SIKT in practice}

SIKT --- the Norwegian agency for shared services in education and research --- assesses the processing of personal data in research. The notification is submitted before data collection, and the assessment takes weeks, which are planned into the timeline. It asks what is collected, from whom, on what legal basis, how it is stored and when it is deleted; the data management plan answers the same questions for the researcher's own use. Nettskjema, developed and run by the University of Oslo and available to UiA users, keeps survey and consent data in Norway and has consent text and anonymous forms built in. Four worked cases in the session --- a survey with e-mail addresses, an analysis of public GitHub data, interviews at a client, and a phishing simulation --- show that the answers differ by case, and that the paperwork is the design, not an appendix to it.

\section{Special zones}

Three areas need extra care. Public data are not free of ethics: forum posts, GitHub commits and Discord messages are publicly visible, but the authors expect an audience of peers rather than researchers, and quoting them verbatim can identify them; the Sámi example of Session 7 makes the same point about communities. Security research raises two further questions --- authorisation, because testing systems one is not authorised to test is an offence regardless of intent, and responsible disclosure, because a vulnerability found in a thesis has an owner who must learn of it first. AI tools raise the newest questions: never paste participant data into an external service, declare which tools were used and for what, and treat the tool's output as text to be checked rather than evidence.

\section{References to keep at hand}

Six references belong on the desk of every student writing an essay. The University of Agder's ethical guidelines and its guidelines for the use of artificial intelligence state what the institution expects of its students and staff. The national guidelines of the research ethics committees (NESH and NENT) at forskningsetikk.no state what the research community expects. SIKT's personal-data service is where the notification is submitted. Nettskjema is where survey and consent data are collected. Lastly, the Montréal Declaration for a Responsible Development of Artificial Intelligence (2018) sets out ten principles --- well-being, autonomy, privacy, solidarity, democratic participation, equity, diversity, prudence, responsibility and sustainability --- that frame any study that builds or evaluates AI, and the student's own use of AI tools.


\section{Worked examples and tables}

\subsection{Consent and anonymisation}
\begin{center}\small
\footnotesize Consent is only valid if it is \textbf{informed, voluntary
  and documented}. The information sheet / consent form states:
  \begin{itemize}\setlength{\itemsep}{0.2em}
    \item \textbf{who} is doing the research, for what \textbf{purpose}
    \item \textbf{what participation involves} (time, recording, topics)
    \item what \textbf{data} is collected, how it is \textbf{stored},
          who can access it, and \textbf{when it is deleted}
    \item that participation is \textbf{voluntary} and can be
          \textbf{withdrawn at any time, without reason or consequence}
    \item whether results are reported \textbf{anonymously} or
          \textbf{confidentially} --- and which one you actually promise
    \item \textbf{contact} for questions and complaints
  \end{itemize}
  {\color{kgrey}\footnotesize SIKT provides consent-form templates --- use
  them rather than writing your own from scratch.}
\end{center}

\begin{center}\small
\footnotesize
  \renewcommand{\arraystretch}{1.35}
  \begin{tabular}{@{}p{7.6cm}p{5.2cm}@{}}
    \toprule
    \textbf{Example wording} & \textbf{Which requirement it meets}\\
    \midrule
    ``We are three master's students investigating how developers handle security warnings.'' & who + purpose\\
    ``Participation means one recorded interview of 30--45 minutes about your daily work.'' & what it involves\\
    ``Recordings are stored on university-approved storage, accessible to the group only, and deleted after transcription.'' & data, storage, deletion\\
    ``Participation is voluntary; you may withdraw at any time without giving a reason.'' & voluntariness, withdrawal\\
    ``Results are reported confidentially: quotes are shown to you before use and stripped of identifying detail.'' & the promise --- and how it is kept\\
    \bottomrule
  \end{tabular}
  

  Five sentences cover the checklist --- built on the SIKT template, adapted
  to your study. Attach the full form to Part IV as an appendix.
\end{center}

\begin{center}\small
\small
  \renewcommand{\arraystretch}{1.4}
  \begin{tabular}{@{}p{6.3cm}p{6.5cm}@{}}
    \toprule
    \textbf{As transcribed} & \textbf{As reported}\\
    \midrule
    ``As the only senior backend dev here, I told Anna in March that the
    payment service was wide open.'' & ``One developer described raising a
    serious service vulnerability with management months before it was
    addressed.''\\
    \bottomrule
  \end{tabular}
  

  What the repair removed: the \textbf{unique role}, the \textbf{named
  colleague}, the \textbf{date}, and the \textbf{specific system} --- any one
  of which identifies the speaker to an internal reader. The finding
  survives; the person is protected. Let participants review their quotes
  before submission.
\end{center}


\subsection{SIKT, and four worked cases}
\begin{center}\small
\begin{center}
    \resizebox{0.95\textwidth}{!}{%
    \begin{tikzpicture}
      \node[kbox, text width=3.5cm, minimum height=1.5cm, anchor=west] (q1) at (0,0)
        {\textbf{Will you process}\\ \textbf{personal data?}\\ \scriptsize incl.\ audio, e-mail, indirect identification};
      \node[gbox, text width=3.3cm, minimum height=1.0cm, anchor=west] (no) at (4.6,0.9)
        {\textbf{No} $\rightarrow$ no notification\\ \scriptsize document \emph{why} in Part IV};
      \node[kbox, text width=3.3cm, minimum height=1.0cm, anchor=west] (yes) at (4.6,-0.9)
        {\textbf{Yes} $\rightarrow$ notify \textbf{SIKT}\\ \scriptsize before data collection starts};
      \node[kbox, text width=3.6cm, minimum height=1.5cm, anchor=west] (form) at (8.9,-0.9)
        {\textbf{The form asks:}\\ \scriptsize purpose, data types, participants,\\ \scriptsize consent, storage, deletion date};
      \draw[garrow] (q1.east) -- (no.west);
      \draw[karrow] (q1.east) -- (yes.west);
      \draw[karrow] (yes) -- (form);
    \end{tikzpicture}}
  \end{center}
  
  \begin{itemize}\setlength{\itemsep}{0.35em}
    \item \textbf{Timeline:} assessment can take \textbf{weeks} --- submit
          well before you plan to collect. In your essay, put the
          notification into the project plan.
    \item \textbf{Storage:} Nettskjema and UiA-approved storage keep
          you inside the rules; private Dropbox and free transcription apps
          do not.
  \end{itemize}
\end{center}

\begin{center}\small
\small
  \textbf{Design:} questionnaire to client staff; e-mail collected for a
  follow-up interview draw.
  \begin{itemize}\setlength{\itemsep}{0.35em}
    \item Data inventory: responses + \textbf{e-mail} $\rightarrow$ personal
          data, direct
    \item SIKT: \textbf{yes} --- notify before launch; deletion date for the
          address list stated
    \item Consent: purpose, voluntariness, storage in Nettskjema, addresses
          deleted after interviews are scheduled
    \item Cheaper alternative worth considering: separate, voluntary sign-up
          form --- the survey itself becomes anonymous
  \end{itemize}
  
  The last bullet is the pattern to learn: \textbf{redesign can remove the
  data problem entirely} --- minimisation by design beats paperwork.
\end{center}

\begin{center}\small
\small
  \textbf{Design:} mine 200 public repositories for dependency-update
  behaviour; no interaction with authors.
  \begin{itemize}\setlength{\itemsep}{0.35em}
    \item Data inventory: commits carry \textbf{usernames and e-mails}
          $\rightarrow$ personal data exists in the raw material
    \item Mitigation: strip identities at ingestion; analyse and report only
          \textbf{aggregates} per repository
    \item SIKT: with identities stripped immediately and nothing
          person-level stored, likely \textbf{no} --- but the reasoning is
          \emph{documented} in Part IV, not assumed
    \item Quoting a named user's commit message in the report would reverse
          this analysis --- don't
  \end{itemize}
\end{center}

\begin{center}\small
\small
  \textbf{Design:} eight interviews with the client's developers about
  security practices.
  \begin{itemize}\setlength{\itemsep}{0.35em}
    \item Personal data: audio, names, role information $\rightarrow$
          \textbf{SIKT yes}; consent per template
    \item The client trap in force: promise \textbf{confidentiality} (not
          anonymity --- you know who said what), aggregate roles, quote
          review by participants
    \item Power relations: the CTO offers to ``send the team'' --- decline;
          recruit directly, stress voluntariness, never report who
          participated
    \item Storage: recordings on approved storage only; deleted after
          transcription; transcripts pseudonymised (P1\dots P8)
  \end{itemize}
\end{center}

\begin{center}\small
\small
  \textbf{Design:} simulated phishing campaign to measure training effect at
  the client.
  \begin{itemize}\setlength{\itemsep}{0.35em}
    \item Deception involved $\rightarrow$ the bar rises: written
          \textbf{management authorisation} with scope (who, when, what
          pretext)
    \item \textbf{No individual-level reporting} to the employer --- click
          rates per department at most, never per person
    \item \textbf{Debriefing}: all recipients informed afterwards what
          happened and why; support contact included
    \item SIKT: click data linked to persons during collection
          $\rightarrow$ notify; anonymise at analysis
  \end{itemize}
  
  If any of the three protections is refused by the client, the honest
  Part IV answer is: this design is not ethically feasible here --- and you
  redesign.
\end{center}


\subsection{The data management plan}
\begin{center}\small
\footnotesize
  \renewcommand{\arraystretch}{1.35}
  \begin{tabular}{@{}p{3.0cm}p{3.3cm}p{3.2cm}p{3.2cm}@{}}
    \toprule
    \textbf{Data} & \textbf{Stored where} & \textbf{Access} & \textbf{Deleted}\\
    \midrule
    audio recordings & approved uni storage & two coders & after transcription\\
    transcripts (pseudonymised) & approved uni storage & group & project end + grading\\
    survey responses & Nettskjema & group & per SIKT notification\\
    key file (name $\leftrightarrow$ P\#) & separate, access-restricted & one person & after quote review\\
    \bottomrule
  \end{tabular}
  

  Four columns answer every storage question SIKT and your consent form ask
  --- copy the structure, fill your rows, attach it to Part IV.
\end{center}


\subsection{Special zones}
\begin{center}\small
\begin{itemize}\setlength{\itemsep}{0.45em}
    \item Forum posts, GitHub commits, Discord messages are \textbf{publicly
          visible} --- but the people who wrote them did not consent to
          becoming \textbf{research subjects}.
    \item Questions to answer in Part IV: do the authors reasonably
          \textbf{expect an audience like you}? Can quotes be
          \textbf{searched} and traced back? Is the community sensitive?
    \item The typing study (session 1) used public GitHub code --- defensible
          because it analysed \textbf{code, not people}, in aggregate.
          Quoting named users' comments would be a different matter.
    \item Virtual ethnography (session 7): observing a community raises the
          overt/covert question --- covert observation needs strong,
          documented justification.
  \end{itemize}
\end{center}

\begin{center}\small
\begin{itemize}\setlength{\itemsep}{0.5em}
    \item \textbf{Authorisation in writing.} Testing systems --- scanning,
          pentesting, phishing simulations --- requires the owner's explicit,
          written permission \emph{with defined scope}. Without it, the same
          activity is an offence, not research.
    \item \textbf{Responsible disclosure.} If your project finds a
          vulnerability: report it to the owner, agree a remediation window,
          and anonymise it in the report. A essay that touches
          vulnerabilities should state this policy in advance.
  \end{itemize}
  
  \begin{block}{Phishing simulations with employees}
    Deception is involved --- so the bar rises: management authorisation,
    debriefing of participants afterwards, no individual-level reporting to
    the employer. State all three in Part IV.
  \end{block}
\end{center}

\begin{center}\small
\begin{alertblock}{The fact that changes your planning}
    Essay reports and master's reports are read by \textbf{external examiners}
    --- and examiners normally sign \textbf{no NDA} with your client.
  \end{alertblock}
  
  \small Consequences, especially for security work:
  \begin{itemize}\setlength{\itemsep}{0.35em}
    \item \textbf{Agree in writing, up front,} what may appear in the graded
          report --- Oates: discuss confidentiality \emph{before} you start,
          via disguise or delayed publication if needed
    \item \textbf{Anonymise the client} where required (``a Norwegian retail
          SMB'') --- the method quality is gradable without the name
    \item \textbf{Pentesting and vulnerabilities:} detailed findings belong
          in a separate \emph{client deliverable}; the graded report carries
          only remediated, generalised descriptions. Appendices do not help
          --- examiners read those too
    \item If the client demands confidentiality that conflicts with
          examination, raise it with the course responsible \textbf{early}
          --- not in submission week
  \end{itemize}
\end{center}

\begin{center}\small
\begin{itemize}\setlength{\itemsep}{0.5em}
    \item \textbf{Never paste participant data} --- transcripts, survey
          exports, log extracts --- into external AI services. That is a
          disclosure to a third party your consent form did not cover.
    \item AI assistance in \textbf{writing and coding} is governed by
          UiA's rules and the \textbf{AI declaration} you submit with
          the essay --- declare what you used and for what.
    \item If your \emph{method} uses AI (e.g.\ AI-assisted coding of
          interview data), that is a methodological choice: describe it,
          justify it, and address its reliability like any other instrument.
  \end{itemize}
\end{center}


\section*{Terms from this session}
\begin{center}\small
\footnotesize
  \renewcommand{\arraystretch}{1.2}
  \begin{tabular}{@{}p{3.8cm}p{9.7cm}@{}}
    \toprule
    \textbf{Personal data} & information relating to an identified or identifiable person, incl.\ indirectly\\
    \textbf{Informed consent} & voluntary, documented agreement based on full information\\
    \textbf{Right to withdraw} & participants may leave at any time, without reason or consequence\\
    \textbf{Anonymity} & answers cannot be linked to persons --- not even by the researcher\\
    \textbf{Confidentiality} & the link exists but is protected and never revealed\\
    \textbf{SIKT notification} & Norwegian notification of projects processing personal data, before collection\\
    \textbf{Responsible disclosure} & reporting found vulnerabilities to the owner before any publication\\
    \textbf{Data minimisation} & collecting only the data the research question requires\\
    \bottomrule
  \end{tabular}
\end{center}



\section*{Questions from the reading}

  \begin{enumerate}\setlength{\itemsep}{0.4em}
    \item Which \textbf{rights} do people directly involved in research have, according to Oates?
          \answer{Five: the right \textbf{not to participate}, to \textbf{withdraw}, to \textbf{informed consent}, to \textbf{anonymity}, to \textbf{confidentiality} --- and organisations hold them too.}
    \item What is the difference between \textbf{anonymity} and \textbf{confidentiality}?
          \answer{Anonymity protects \textbf{identity} (pseudonyms, disguise --- sometimes even gender, in a small department); confidentiality protects the \textbf{data} and the link --- kept safe, never passed on.}
    \item Why is internet research ethically difficult even when the data is \textbf{publicly visible}?
          \answer{Public visibility is not consent to be \emph{studied}: ownership of posts is unclear, people treat ``public'' groups as private, consent at scale is impractical, and data linkage defeats anonymisation.}
    \item Which \textbf{responsibilities} does Oates assign to the ethical researcher --- beyond following the law?
          \answer{Respect expectations of anonymity and confidentiality; never coerce --- especially from a position of power; obtain informed consent and do not deceive unless there is no alternative and no harm can ensue.}
  \end{enumerate}

\section*{Questions for discussion}

  \begin{enumerate}\setlength{\itemsep}{0.25em}
    \item \textbf{Data inventory} --- list every data type your design collects. Which of them are personal data, directly or indirectly?
          \answer{E-mail addresses, IP addresses, voice recordings and free-text answers are personal data; so are job titles in a small organisation, because they identify indirectly. If anything on the list is personal, the SIKT decision follows.}
    \item \textbf{Consent} --- what must your consent form say, and what may you promise: anonymity or confidentiality?
          \answer{Purpose, what participation involves, storage and deletion, the right to withdraw, and who to contact. Promise confidentiality; promise anonymity only if you will never be able to link data to a person --- which is rarely true of interviews.}
    \item \textbf{Your ethics paragraph} --- write the Part IV paragraph from the inventory, the SIKT decision and the risk table. Which risk is specific to security research?
          \answer{Scope: testing systems you are not authorised to test, and reporting vulnerabilities to examiners who have signed no NDA. State the authorisation you have and what stays out of the graded essay.}
  \end{enumerate}

\section*{Key points}
\begin{itemize}
  \item Personal data is broader than it looks --- and the law obliges you
  \item Informed consent, withdrawal, anonymity versus confidentiality
  \item SIKT: when, how, and where it sits in your timeline; the data management plan
  \item Special zones: public data, security research, AI tools
\end{itemize}

\section*{Sources}
\begin{itemize}\small
  \item Oates, B.\,J., Griffiths, M., \& McLean, R. (2022). \emph{Researching information systems and computing} (2nd ed.). SAGE. --- Chapters 4--5 (The digital researcher; Participants and research ethics).
  \item Universit\'e de Montr\'eal. (2018). \emph{Montr\'eal Declaration for a Responsible Development of Artificial Intelligence}. \url{https://montrealdeclaration-responsibleai.com}
  \item Laudon, K.\,C., \& Laudon, J.\,P. (2022). \emph{Management information systems: Managing the digital firm} (17th ed., Global ed.). Pearson. --- Chapter 4 (Ethical and social issues in information systems).
\end{itemize}

\part{Writing}
\chapter{Producing the Essay}
\label{ch:11}
\sessionhead{11}{Thursday 19 November 2026 \;$\cdot$\; reading: Oates, Chapter 21}

The essay is an argument, not a form. This chapter presents it as a paper --- introduction, theoretical foundations, method, findings, discussion, conclusion --- with the question each section answers, and then works through the craft of reading and synthesising literature with SQ3R and the concept matrix. The describe--justify pattern, the checklist part by part, the recurring failures and the toolchain complete the chapter. It is the chapter to keep open while writing.


\section{The essay as argument}

The essay is an argument with a golden thread: the research question comes first, and every choice after it traces back to that question. Where the thread breaks is where the examiner's first question will come from, and the typical break is a strategy chosen in Part III that the question in Part I does not require.

Read as a paper, the essay has six sections, and each answers one question. The introduction answers what the problem is, who has it, why now --- and what the research question is (Part I). The theoretical foundations answer what is already known and through which theory or framework we look (Part II). The method answers how we find out --- strategy, methods, sample, analysis --- and why (Part III). The findings answer what was found or, in a proposal, what the analysis will produce and how it will be reported (Part IV). The discussion answers what the findings mean against the literature and what the limitations are. Lastly, the conclusion answers what the contribution is and, above all, \emph{so what}: who acts differently because of it. This is the structure of the papers read in this course --- Huy et al., Enholm et al., Woodruff et al. --- and of my own papers. A literature review can be the basis: then the method section is the search protocol and the findings are the synthesis.

Boodraj et al.\ (2026) show the shape in six pages. Their introduction sets out project failure rates, cyber threats and the absence of a standard for the competencies of a cyber security project manager --- the gap and the question. Their foundations draw two themes from the literature. Their method is a labour market analysis of 161,000 job postings with Python analytics, with a market basket analysis planned. Their preliminary findings are dual-domain certifications and interpersonal skills in demand, and their discussion asks how organisations recruit, train and integrate these managers. Every section is present, part of the work is still ahead, and the topic is cyber security: this is what the essay looks like when it is done well.

\section{Writing methodology}

\begin{block}{Describe--justify}
The describe--justify pattern refers to writing every methodological choice as a pair: what was chosen (describe), and why this choice fits the question and the resources better than the alternatives (justify), with the source cited at the choice.
\end{block}

``We use semi-structured interviews (describe) because the question asks for practitioners' reasoning, which a questionnaire cannot capture (justify; Oates, Chapter 13).'' Every choice in Parts III and IV carries a because, and the curriculum is cited at the choice, not in a separate list of definitions. Four before--after excerpts in the session show the same transformation in each part: a draft that names and describes becomes a draft that argues. Limitations are written in three sentences --- what the study cannot show, why, and what would be needed to show it --- and are graded as judgement, not confessed as faults. The signposting of a strong Part III opens with the strategy, the reason and the alternative in three sentences, so that a reader knows the argument before the details.

\section{Reading and synthesising}

Writing the foundations depends on reading well, and reading well is a method. SQ3R --- survey, question, read, recite, review --- was published by Robinson in 1946 for people who had to work through dense technical material fast and remember it, and the problem is unchanged: passive, linear reading yields poor retention. Survey takes two minutes: title, first paragraph, every heading, every table, the conclusion, for orientation rather than comprehension. Question turns each heading into a question before the section is read, so that reading has a target. Read hunts for the answers and marks them. Recite closes the paper and says the findings in one's own words; if that fails, the section is read again. Review goes back over the whole paper for its throughline --- the one idea that runs through it. A paper read this way takes forty minutes and leaves a record; a paper read front to back takes longer and leaves an impression. Presthus and Sørum's (2025) survey of consumer perspectives on AI serves as the worked example in the session: an online survey of 306 Norwegians and seven everyday scenarios, eight pages, one idea --- people want AI in addition to humans, not instead.

The record goes into the concept matrix. Webster and Watson (2002), then editors at \emph{MIS Quarterly}, wrote their guide because reviews in the field read like phone books --- an impressive cast of names, no plot. Their answer was to organise the review by concept rather than by author, and to keep a small table while reading: one row per paper, one column per concept, a mark where a paper covers it. The row is added during the review pass while the paper is open; a concept not yet in the table receives a new column. The matrix is then read down a column, and the column is the material for one synthesised section: several papers, one argument. Watson and Webster (2020) take the matrix one step further. A mark becomes a 1, a blank a 0, and the grid is the recipe for a graph in which every 1 is a line between a paper and a concept; papers that meet at a concept form one paragraph, and missing edges mark open questions. In their release 2.0 the edges carry a type --- relates to, causes, moderates --- so that each paper's core model becomes a small graph, coded in Cypher and loaded into a graph database, and the graphs combine across papers into one searchable synthesis. The matrix is binary: a mention is a 1, never a count. The graph is inspiration, not proof.

Applied to the papers of this semester, the matrix has one row for Woodruff et al., Huy et al., Jo and Park, and Enholm et al., and columns for expected impact, mechanism, human role and method. Read by column, the mechanism column already contains a finding: the surveys explain adoption, the workshop study explains its limits, the review explains value. That is a paragraph of the foundations, and it is synthesis rather than summary.

\section{Part by part}

Parts I and II must contain the background and the prevailing knowledge, the problem with its owner and cost, the aims and objectives, one main research question with sub-questions and its justification, and the contribution to practice and, where applicable, to academia; Part II adds the keyword table, the search log with databases, strings and dates, the inclusion and exclusion criteria, a PRISMA-style flow, and a synthesis in blocks. Parts III and IV must contain the strategy named and argued with fit to question, fit to resources and rejected alternatives; the data generation methods, sample and instrument; the analysis plan; the ethics paragraph with data inventory, consent, SIKT and storage; and the limitations. The formalities are about ten pages excluding front page and references, English or Norwegian, APA 7 referencing, and an honest and specific AI declaration --- UiA's rules on the use of artificial intelligence apply.

A rough page budget spends the ten pages deliberately: about two pages on the introduction, problem and question, about two and a half on relevant work, about three on the choice of methods, and about two and a half on data, analysis, limitations and ethics. Drafts overspend on reproduction --- definitions, background --- and underspend on judgement --- justification, limitations --- which is the exact opposite of what is graded. Tables carry information cheaply.

\section{Naming the contribution}

The conclusion states the form of the contribution in Presthus and Munkvold's (2016) vocabulary --- a confirmation, extension or contradiction of a named model; a taxonomy; a set of guidelines; a rich case description --- and for whom: a contribution to practice names a decision that becomes better informed, a contribution to theory names a gap that becomes smaller. The decision is taken before the data are collected, and the conclusion is written first as a target and revised when the findings are in. The so-what test is simple: cover the findings, and read the conclusion. If it still tells a security manager, a developer or a policymaker what to do differently on Monday, the essay has a contribution; if it only restates the findings, it has a summary.

\section{Failure modes}

Seven failures recur. Textbook methodology: pages of definitions and no argument. Unjustified strategy: the strategy named but not argued, with no ``we rejected \dots\ because''. Shopping-list literature: paper by paper, no synthesis. Broken thread: a question that Part III does not answer. Afterthought ethics: a paragraph copied from a template, without a data inventory. Missing limitations: none stated, which reads as none noticed. Vague plan: ``interviews will be conducted'' with no numbers, no guide, no analysis. All seven are avoidable before submission, and the checklist of this chapter is the instrument.

\section{The toolchain}

One rule beats all tool debates: decide the tools now and use the same ones until submission. Zotero, free, collects every paper with its metadata from a browser button and cites while you write in Word and LaTeX; one note per paper holds the seven items you mark --- problem and gap, question, method, findings, admitted limitations, the so-what, and one sentence worth quoting. Word with the Zotero plug-in is the low-friction default; LaTeX, in Overleaf or locally, if you know it. Grammarly or Word's editor serve grammar and clarity only, never content, and are declared if used. draw.io draws the research process and the model as figures, in BPMN notation where a process is described. Excel holds the search log and the concept matrix, one sheet each. Nettskjema collects the data; SmartPLS runs PLS-SEM; JASP handles descriptive and inferential statistics; NVivo supports qualitative coding.


\section{Worked examples and tables}

\subsection{What the examiners look for}
\begin{center}\small
\small
  \renewcommand{\arraystretch}{1.4}
  \begin{tabular}{@{}p{5.4cm}p{7.6cm}@{}}
    \toprule
    \textbf{Essay part} & \textbf{The examiners look for}\\
    \midrule
    I --- introduction, problem, research questions & One question, a clear hierarchy, justified relevance\\
    II --- relevant work & A documented, repeatable search; synthesis, not a list\\
    III --- choice of methods & Choices \emph{argued}: fit to question and resources, alternatives rejected with reasons\\
    IV --- data, analysis, limitations, ethics & Judgement: honest limitations, concrete ethics and privacy, a realistic plan\\
    \bottomrule
  \end{tabular}
  
  Weights: the written essay counts 40\,\% of the course grade, its
  presentation and discussion in the oral exam 20\,\%. Note where judgement
  lives: Parts III and IV. That is where grades separate.
\end{center}

\begin{center}\small
\begin{center}
    \resizebox{0.96\textwidth}{!}{%
    \begin{tikzpicture}
      \node[kbox, text width=2.2cm, minimum height=1.0cm, anchor=west] (a) at (0,0)    {\textbf{Problem}};
      \node[kbox, text width=2.2cm, minimum height=1.0cm, anchor=west] (b) at (2.7,0)  {\textbf{Question}};
      \node[kbox, text width=2.2cm, minimum height=1.0cm, anchor=west] (c) at (5.4,0)  {\textbf{Strategy}};
      \node[kbox, text width=2.2cm, minimum height=1.0cm, anchor=west] (d) at (8.1,0)  {\textbf{Methods}};
      \node[kbox, text width=2.2cm, minimum height=1.0cm, anchor=west] (e) at (10.8,0) {\textbf{Analysis}};
      \draw[karrow] (a) -- (b); \draw[karrow] (b) -- (c);
      \draw[karrow] (c) -- (d); \draw[karrow] (d) -- (e);
    \end{tikzpicture}}
  \end{center}
  
  Every element must \textbf{answer to the one before it}. The most common
  structural failure is a broken thread: a qualitative question followed by a
  survey, or an analysis plan that ignores half the collected data.
  
  \begin{block}{The reviewer's test --- apply it to your own draft}
    Read only your research question, then only your methods section.
    Could a stranger tell they belong to the same project?
  \end{block}
\end{center}


\subsection{Describe and justify --- worked pairs}
\begin{center}\small
\small
  Every methodological statement has two halves:
  
  \begin{center}
    \begin{tikzpicture}
      \node[gbox, text width=5.6cm, minimum height=1.5cm, anchor=west] (d) at (0,0)
        {\textbf{Describe} (what, precisely)\\ \scriptsize ``Eight semi-structured
         interviews, 30--45 minutes, recorded and transcribed verbatim\dots''};
      \node[kbox, text width=5.6cm, minimum height=1.5cm, anchor=west] (j) at (6.6,0)
        {\textbf{Justify} (why this, here)\\ \scriptsize ``\dots because the question
         asks how sysadmins \emph{experience} security trade-offs, which
         requires depth over breadth (Oates, 2026).''};
      \draw[karrow] (d) -- (j);
    \end{tikzpicture}
  \end{center}
  
  Writing conventions that help:
  \begin{itemize}\setlength{\itemsep}{0.3em}
    \item \textbf{Future tense} --- a essay plans (``we will recruit''),
          it does not pretend to have finished
    \item \textbf{Precision} --- numbers, durations, tools, versions;
          ``several interviews'' is not a plan
    \item \textbf{References at choices} --- the method literature (Oates,
          Creswell) is cited \emph{where the choice is made}, not dumped in
          one paragraph
  \end{itemize}
\end{center}

\begin{center}\small
\footnotesize
  \renewcommand{\arraystretch}{1.4}
  \begin{tabular}{@{}p{5.9cm}p{7.3cm}@{}}
    \toprule
    \textbf{Describe} & \textbf{Justify}\\
    \midrule
    ``An online questionnaire (12 items, Nettskjema) will be sent to all 60 client staff\dots'' & ``\dots because SQ1 asks how widespread reuse is, which requires breadth across the whole organisation.''\\
    ``The component library will be evaluated in a usability test (8 users, 5 tasks, SUS)\dots'' & ``\dots because the research claim concerns developer-facing usability, not runtime performance.''\\
    ``The incident will be studied as a single explanatory case\dots'' & ``\dots because the question asks \emph{how} recovery unfolded in context --- depth in one setting over breadth.''\\
    \bottomrule
  \end{tabular}
  

  Note the mechanical link: every justification points back at the
  \textbf{question's wording}. That is the golden thread, sentence by
  sentence.
\end{center}

\begin{center}\small
\small
  \begin{block}{Where the references belong}
    ``Since the question asks how developers \emph{experience} accessibility
    work, a qualitative approach is indicated (Creswell \& Creswell, 2018).
    We therefore use semi-structured interviews, which combine a focused
    guide with room for follow-up (Oates et al., 2026, Ch.\ 13). Analysis
    follows the thematic-analysis steps of Braun and Clarke (2006).''
  \end{block}
  
  \begin{itemize}\setlength{\itemsep}{0.4em}
    \item Three citations, each attached to the \textbf{decision it
          licenses} --- approach, method, analysis
    \item Contrast with the failure mode: ten references parked in one
          introductory sentence, then an unreferenced methods section
    \item This is also how the examiners verify you \emph{used} the
          curriculum, which the essay text explicitly asks for
  \end{itemize}
\end{center}


\subsection{Before and after --- one excerpt per part}
\begin{center}\small
\small
  \renewcommand{\arraystretch}{1.4}
  \begin{tabular}{@{}p{6.3cm}p{6.5cm}@{}}
    \toprule
    \textbf{Draft} & \textbf{Repaired}\\
    \midrule
    ``Cyber security is very important nowadays. Many companies get hacked. We want to look at passwords.'' & ``Despite mandatory annual training, the client reports recurring password reuse across systems. Reuse turns one credential leak into many; for an SMB without a security team, that multiplication is the main breach risk. This essay therefore asks which factors are associated with reuse among the client's employees.''\\
    \bottomrule
  \end{tabular}
  

  The repair adds: an \textbf{owner} (the client), a \textbf{cost}
  (multiplication of leaks), and a direct handover to the \textbf{question}.
  Zero words spent on ``nowadays''.
\end{center}

\begin{center}\small
\small
  \renewcommand{\arraystretch}{1.4}
  \begin{tabular}{@{}p{6.3cm}p{6.5cm}@{}}
    \toprule
    \textbf{Draft} & \textbf{Repaired}\\
    \midrule
    ``Smith (2021) surveyed employees. Lee (2022) did an experiment. Hansen (2023) interviewed admins.'' & ``Across survey studies, security fatigue consistently predicts reuse (Smith, 2021; Lee, 2022). The one interview study complicates this: admins describe reuse as a deliberate time trade-off, not exhaustion (Hansen, 2023). No study examines SMBs, where both mechanisms plausibly intensify --- the gap this essay addresses.''\\
    \bottomrule
  \end{tabular}
  

  Same three sources --- but the repaired version runs the \textbf{four
  moves}: establish, complicate, gap, therefore.
\end{center}

\begin{center}\small
\small
  \renewcommand{\arraystretch}{1.4}
  \begin{tabular}{@{}p{6.3cm}p{6.5cm}@{}}
    \toprule
    \textbf{Draft} & \textbf{Repaired}\\
    \midrule
    ``A case study is an in-depth investigation of one instance in its natural setting (Oates, 2026). Case studies can be exploratory, descriptive or explanatory\dots'' [continues for a page] & ``We choose a single explanatory case study because the question asks \emph{how} trust was rebuilt in this specific recovery (Oates, 2026). A survey was rejected: no population of comparable recoveries is accessible. An experiment was rejected: the event cannot be manipulated.''\\
    \bottomrule
  \end{tabular}
  

  The textbook page becomes three sentences of \textbf{choices} --- and the
  citation moves to where the choice is made.
\end{center}

\begin{center}\small
\small
  \renewcommand{\arraystretch}{1.4}
  \begin{tabular}{@{}p{6.3cm}p{6.5cm}@{}}
    \toprule
    \textbf{Draft} & \textbf{Repaired}\\
    \midrule
    ``Some interviews will be conducted and analysed. Ethical guidelines will be followed.'' & ``Eight semi-structured interviews (30--45 min) with developers recruited directly (not via management), recorded, transcribed, thematically analysed with a code book; two transcripts double-coded. Audio and role data are personal data: SIKT will be notified before recruitment; confidentiality (not anonymity) is promised and enforced by quote review.''\\
    \bottomrule
  \end{tabular}
  

  Every vague noun acquired a \textbf{number, a procedure, or a decision}.
  That is the whole difference between the drafts the examiners mark down and
  the ones they don't.
\end{center}


\subsection{Must-haves and the page budget}
\begin{center}\small
\small
  \begin{columns}[T]
    \column{0.48\textwidth}
    \textbf{Part I}
    \begin{itemize}\setlength{\itemsep}{0.25em}
      \item background: prevailing knowledge, key challenges
      \item problem with an \textbf{owner and a cost}
      \item aim, objectives
      \item \textbf{one} main research question; sub-questions in clear
            hierarchy
      \item contribution to practice (and academia, if applicable);
            why the project is important and interesting
    \end{itemize}
    \column{0.48\textwidth}
    \textbf{Part II}
    \begin{itemize}\setlength{\itemsep}{0.25em}
      \item keyword table; \textbf{one} final search string
      \item databases named and justified
      \item search log (iterations, hit counts)
      \item inclusion/exclusion criteria
      \item results in a table; \textbf{PRISMA figure}
      \item what the included papers taught you --- feeding Part III
    \end{itemize}
  \end{columns}
\end{center}

\begin{center}\small
\small
  \begin{columns}[T]
    \column{0.48\textwidth}
    \textbf{Part III}
    \begin{itemize}\setlength{\itemsep}{0.25em}
      \item strategy named --- and \textbf{argued}: fit to question, fit to
            resources
      \item rejected alternatives, with reasons
      \item data generation methods, concretely (guide/items in appendix)
      \item consistency with the approach implied by the question
    \end{itemize}
    \column{0.48\textwidth}
    \textbf{Part IV}
    \begin{itemize}\setlength{\itemsep}{0.25em}
      \item data plan: what, from whom, how much, when
      \item analysis plan (session 8/9 templates)
      \item \textbf{limitations} --- named, not hidden
      \item ethics: personal data inventory, consent, \textbf{SIKT}
            decision documented
      \item expected contributions, revisited
    \end{itemize}
  \end{columns}
  
  {\color{kgrey}\small Formalities: about 10 pages excl.\ front page and references; English or Norwegian; APA 7 referencing; AI declaration submitted and honest. And remember session 10: examiners
  sign no NDA --- client-confidential details (esp.\ vulnerabilities) do not
  belong in the graded report.}
\end{center}

\begin{center}\small
\small
  \renewcommand{\arraystretch}{1.35}
  \begin{tabular}{@{}p{6.2cm}p{2.6cm}p{4.0cm}@{}}
    \toprule
    \textbf{Part} & \textbf{Rough budget} & \textbf{Where drafts overspend}\\
    \midrule
    I --- introduction, problem, question & $\sim$2 pages & General background prose\\
    II --- relevant work & $\sim$2.5 pages & Summarising paper by paper\\
    III --- choice of methods & $\sim$3 pages & Textbook definitions\\
    IV --- data, analysis, limitations, ethics & $\sim$2.5 pages & Boilerplate ethics\\
    \bottomrule
  \end{tabular}
  
  A guide, not a rule --- but the pattern is stable: drafts overspend on
  \textbf{reproduction} (definitions, background) and underspend on
  \textbf{judgement} (justification, limitations) --- the exact opposite of
  what is graded. Tables (search log, matrix, PRISMA) carry information cheaply.
\end{center}


\subsection{Signposting, limitations and the AI declaration}
\begin{center}\small
\begin{block}{The first paragraph of Part III, as a working example}
    ``This part justifies the methodological design. Section 3.1 argues for
    a mixed-methods approach based on the two sub-questions; 3.2 details the
    survey (instrument, sample, procedure); 3.3 details the interview study;
    3.4 states how the two strands will be integrated.''
  \end{block}
  
  \begin{itemize}\setlength{\itemsep}{0.4em}
    \item Four sentences, and the examiner now reads with a map --- nothing
          surprises, nothing is searched for
    \item The same pattern in miniature at every section start: \emph{what
          this section does, in one sentence}
    \item Cost: $\sim$50 words per part. Return: the reader never wonders
          whether something is missing or merely elsewhere
  \end{itemize}
\end{center}

\begin{center}\small
\begin{block}{From a strong Part IV}
    ``Three limitations bound our claims. First, the sample is a census of
    one organisation; findings generalise analytically, not statistically.
    Second, reuse is self-reported; social desirability likely deflates it,
    so our estimates are conservative. Third, the first author works at the
    client; to reduce interviewer effects, interviews are conducted by the
    other members and coding is done in pairs.''
  \end{block}
  
  What makes it strong: each limitation is \textbf{named}, its
  \textbf{direction of bias} is reasoned about, and where possible a
  \textbf{mitigation} is attached. No limitation is apologised for ---
  bounded claims are what research is.
\end{center}

\begin{center}\small
\small
  \renewcommand{\arraystretch}{1.4}
  \begin{tabular}{@{}p{5.8cm}p{7.0cm}@{}}
    \toprule
    \textbf{Vague (looks evasive)} & \textbf{Specific (looks honest)}\\
    \midrule
    ``AI tools were used to some extent during the work.'' & ``ChatGPT was used to brainstorm synonym candidates for the search string (Part II) and to check grammar in the final draft. All sources were located, read and selected by the group; no participant data was entered into any AI tool; no text was generated for Parts I--IV.''\\
    \bottomrule
  \end{tabular}
  

  The specific version names \textbf{tools, tasks, and boundaries}. In an
  essay that explicitly values transparency and traceability, specificity is
  self-protection, not confession.
\end{center}


\subsection{The concept matrix with this semester's papers}
\begin{center}\small
\small One row per paper, one column per concept; the cells hold what the
  paper says about the concept, not whether it mentions it.
  
  \footnotesize
  \renewcommand{\arraystretch}{1.15}
  \begin{tabular}{@{}p{2.7cm}p{2.3cm}p{2.5cm}p{2.4cm}p{2.1cm}@{}}
    \toprule
    \textbf{Paper} & \textbf{Expected impact} & \textbf{Mechanism} & \textbf{Human role} & \textbf{Method}\\
    \midrule
    Woodruff et al.\ (2024) & Not disruptive; amplifies four forces & Deskilling, dehumanisation, disconnection, disinformation & Reviewer of AI output & Workshops, 54\\
    Huy et al.\ (2024) & Adoption driven by expectancy and fit & UTAUT2 + TTF paths & User; curiosity moderates & Survey, 671\\
    Jo \& Park (2024) & Information support, knowledge acquisition & Perceived intelligence, self-learning & Worker aged 20--40 & Survey, 351\\
    Enholm et al.\ (2022) & Business value, first- and second-order & Enablers, inhibitors, typologies & Organisation & SLR, 43\\
    \bottomrule
  \end{tabular}
  

  \small Read the matrix by \textbf{column}: the ``mechanism'' column already
  contains a finding --- surveys explain adoption, the workshop study explains
  its limits, the review explains value. That is a paragraph of section 2:
  synthesis, not summary (Webster \& Watson, 2002; Watson \& Webster, 2020).
\end{center}


\subsection{The seven recurring failures}
\begin{center}\small
\footnotesize
  \renewcommand{\arraystretch}{1.28}
  \begin{tabular}{@{}p{4.6cm}p{8.7cm}@{}}
    \toprule
    \textbf{Failure} & \textbf{What it looks like --- and the fix}\\
    \midrule
    Textbook methodology & Part III defines ``what a case study is'' for a page --- replace definitions with \emph{your} justified choices\\
    Unjustified strategy & strategy named, never argued --- add fit-to-question, fit-to-resources, rejected alternatives\\
    Shopping-list literature & paper-by-paper summary --- rewrite by concept, end each block facing your question\\
    Broken thread & question and methods from different projects --- run the reviewer's test\\
    Afterthought ethics & one generic paragraph --- data inventory, consent, SIKT decision, concretely\\
    Missing limitations & the plan admits no weaknesses --- name convenience sampling, self-report, scope\\
    Vague plan & ``some interviews will be conducted'' --- numbers, durations, recruitment, tools\\
    \bottomrule
  \end{tabular}
\end{center}


\section*{Terms from this session}
\begin{center}\small
\footnotesize
  \renewcommand{\arraystretch}{1.25}
  \begin{tabular}{@{}p{3.8cm}p{9.7cm}@{}}
    \toprule
    \textbf{Golden thread} & the chain problem $\rightarrow$ question $\rightarrow$ strategy $\rightarrow$ methods $\rightarrow$ analysis, each answering to the previous\\
    \textbf{Describe--justify pattern} & every methodological statement: what precisely, then why this here\\
    \textbf{Signposting} & telling the reader where they are and what comes next\\
    \textbf{Word budget} & deliberate allocation of the 10 pages across the four parts\\
    \textbf{Limitations section} & the study's named weaknesses --- rewarded when honest, found when hidden\\
    \textbf{AI declaration} & the submitted statement of what AI assistance was used, and for what\\
    \bottomrule
  \end{tabular}
\end{center}



\section*{Questions from the reading}

  \begin{enumerate}\setlength{\itemsep}{0.4em}
    \item What belongs in a \textbf{methodology chapter}, according to Oates?
          \answer{Your chosen \textbf{combination of strategy and data generation methods --- and why}; Oates notes it is often the easiest chapter to draft early: you know what you did and why.}
    \item What is the difference between \textbf{describing} a method and \textbf{justifying} it?
          \answer{Describing states what was done; justifying builds the \textbf{argument}: what question, how answered, what evidence, and \emph{so what} --- a thesis \emph{is} an argument.}
    \item What makes academic writing \textbf{criticisable} --- in the good sense from session 1?
          \answer{Full documentation opens the work to \textbf{scrutiny}: readers can check the process --- which is session 1's criticisability, delivered.}
    \item What does Oates say about \textbf{structure and signposting} for the reader?
          \answer{The conventional structure (introduction $\rightarrow$ review $\rightarrow$ methodology $\rightarrow$ results $\rightarrow$ discussion $\rightarrow$ limitations $\rightarrow$ conclusions), meaningful headings, and \textbf{signposts} telling the reader what came and what comes next.}
  \end{enumerate}

\section*{Questions for discussion}

  \begin{enumerate}\setlength{\itemsep}{0.25em}
    \item \textbf{Self-assessment} --- score your draft against the part-by-part checklist. Which part is weakest, and what is the one sentence it is missing?
          \answer{In most drafts Part III is weakest: the strategy is named but not argued. The missing sentence usually begins with ``we rejected \dots\ because''.}
    \item \textbf{Describe and justify} --- take one methodological choice from your draft and write it as a describe--justify pair, citing the curriculum at the choice.
          \answer{``We use semi-structured interviews (describe) because the question asks for practitioners' reasoning, which a questionnaire cannot capture (justify; Oates, Ch.\ 13).'' Every choice in Part III and IV carries a because.}
    \item \textbf{Limitations} --- write the limitations paragraph in three sentences: what the study cannot show, why, and what would be needed to show it.
          \answer{Limitations are graded as judgement, not confessed as faults. A limitation that comes with the design that would remove it (a larger sample, a second site, a longer period) reads as a researcher's sentence, not a student's.}
  \end{enumerate}

\section*{Key points}
\begin{itemize}
  \item The essay as a paper: introduction, theoretical foundations, method, findings, discussion, conclusion --- each answers one question
  \item Reading with SQ3R, marking seven things, one note per paper; the concept matrix turns notes into synthesis
  \item Describe and justify: every choice carries a because, cited at the choice
  \item Name the contribution and pass the so-what test; seven recurring failures --- all avoidable before submission
\end{itemize}

\section*{Sources}
\begin{itemize}\small
  \item Enholm, I.\,M., Papagiannidis, E., Mikalef, P., \& Krogstie, J. (2022). Artificial intelligence and business value: A literature review. \emph{Information Systems Frontiers}, \emph{24}, 1709--1734. \url{https://doi.org/10.1007/s10796-021-10186-w}
  \item Huy, L.\,V., Nguyen, H.\,T.\,T., Vo-Thanh, T., Thinh, N.\,H.\,T., \& Tran, T.\,T.\,D. (2024). Generative AI, why, how, and outcomes: A user adoption study. \emph{AIS Transactions on Human-Computer Interaction}, \emph{16}(1), 1--27. \url{https://doi.org/10.17705/1thci.00198}
  \item Jo, H., \& Park, D.-H. (2024). AI in the workplace: Examining the effects of ChatGPT on information support and knowledge acquisition. \emph{International Journal of Human--Computer Interaction}, \emph{40}(23), 8091--8106. \url{https://doi.org/10.1080/10447318.2023.2278283}
  \item Presthus, W., \& Munkvold, B.\,E. (2016). How to frame your contribution to knowledge? A guide for junior researchers in information systems. \emph{NOKOBIT}, \emph{24}(1). Paper presented at NOKOBIT 2016, Bergen.
  \item Watson, R.\,T., \& Webster, J. (2020). Analysing the past to prepare for the future: Writing a literature review a roadmap for release 2.0. \emph{Journal of Decision Systems}, \emph{29}(3), 129--147. \url{https://doi.org/10.1080/12460125.2020.1798591}
  \item Webster, J., \& Watson, R.\,T. (2002). Analyzing the past to prepare for the future: Writing a literature review. \emph{MIS Quarterly}, \emph{26}(2), xiii--xxiii.
  \item Woodruff, A., Shelby, R., Kelley, P.\,G., Rousso-Schindler, S., Smith-Loud, J., \& Wilcox, L. (2024). How knowledge workers think generative AI will (not) transform their industries. In \emph{Proceedings of the CHI Conference on Human Factors in Computing Systems (CHI '24)}. ACM. \url{https://doi.org/10.1145/3613904.3642700}
  \item Boodraj, M., Fairbanks, J., Davies, C., \& Akcakaya, O. (2026). Beyond the firewall: The rise of the cybersecurity project manager. In \emph{AMCIS 2026 Proceedings} (Emergent Research Forum paper). AIS. \url{https://aisel.aisnet.org/amcis2026/sig_itpm/sig_itpm/2}
  \item Presthus, W., \& S\o rum, H. (2025). Consumer perspectives on the increasing use of artificial intelligence: It is helpful, scary or so wrong? \emph{Procedia Computer Science}, \emph{256}, 158--165. \url{https://doi.org/10.1016/j.procs.2025.02.108}
  \item Robinson, F.\,P. (1946). \emph{Effective study}. Harper \& Brothers.
  \item Oates, B.\,J., Griffiths, M., \& McLean, R. (2022). \emph{Researching information systems and computing} (2nd ed.). SAGE. --- Chapter 21 (Presentation of the research).
  \item Laudon, K.\,C., \& Laudon, J.\,P. (2022). \emph{Management information systems: Managing the digital firm} (17th ed., Global ed.). Pearson. --- Chapter 14 (Making the business case for information systems and managing projects) --- rigour and relevance seen from the organisation.
\end{itemize}

\chapter{Summing-up}
\label{ch:12}
\sessionhead{12}{Thursday 26 November 2026 \;$\cdot$\; reading: the whole book}

The final chapter assembles the course in one picture. It flashes back through the twelve sessions, shows the running example end to end, diagnoses the mismatches between question and strategy that most often cost marks, and sets out what separates the grades. It ends where the course began: with the question \emph{how do they know?} --- and with the three places this pipeline reappears after the course.


\section{The course in one picture}

Twelve sessions form one pipeline: from problem and purpose to research question, from the question to the literature, from the literature to a strategy, from the strategy to methods of generating and analysing data, and from the data to ethics and the report. The pipeline is the essay's table of contents, the method backbone of the master's thesis, and the skeleton of any empirical work done later. Every box has been filled at least once; the essay is the pipeline written down.

The flashbacks follow the sessions. Session 2 asked one question three ways and showed that the approach must be justified. Session 3 supplied the 6Ps, one line each. Session 4 named the five failure modes of a research question --- too broad, trivial yes or no, not researchable, double-barrelled, assumption baked in --- and their repairs. Session 5 built the review chain: keyword table, string and log, criteria and PRISMA flow, concept matrix, the four moves of synthesis. Sessions 6 and 7 placed the six strategies on one table, one line each. Sessions 8 and 9 built the two data plans, quantitative and qualitative, with the same rigour. Session 10 stated the ethics questions: what is collected, from whom, on what legal basis, stored how, deleted when. Session 11 turned the whole into an argument.

\section{One study, end to end}

The running example assembles itself from the building blocks. A problem in a security operations centre --- analysts ignore alerts --- becomes a purpose, a main question with two sub-questions, a documented literature search, a case study strategy with justified alternatives, semi-structured interviews with an information-power argument for their number, a thematic analysis plan, an ethics paragraph with data inventory and SIKT decision, and stated limitations. Read in that order it is the essay; read backwards it is the master's thesis waiting for its data.

The mismatch clinic diagnoses the most common failure in Part III. A survey used to understand why administrators ignore alerts asks a why-question with a breadth strategy. An experiment with no control group has no comparison and therefore no cause. An ethnography of three weeks lacks the immersion the strategy needs. Each is a mismatch between question and strategy, and each is visible in the justification paragraph --- or in its absence.

\section{The essay and the oral exam}

The facts of submission are few: about ten pages excluding front page and references, the four parts clearly marked; English or Norwegian, as a PDF; the written essay counts 40\,\% and its presentation and discussion 20\,\%; APA 7 references with the curriculum cited at the choices; an honest and specific AI declaration; and the deadline, portal and oral exam dates as announced. The oral exam is the essay presented and its choices defended in discussion --- the same describe--justify pairs, spoken. Frequently asked questions have short answers: tables count towards the ten pages and are used for density, not evasion; there is no reference quota, but every claim that needs support has it; the research question may still change if Parts II to IV are updated to match; in a group, each member must be able to defend every part; and a pending SIKT decision is stated in Part IV with its submission date, because the plan is graded, not the approval letter.

What separates the grades is judgement rather than reproduction. In Parts I and II, reproduction means background and literature reported correctly, and judgement means one question with a hierarchy and a synthesis that argues. In Parts III and IV, reproduction means strategy and methods named and described, and judgement means choices justified, alternatives rejected with reasons, limitations and ethics concrete. In the oral exam, reproduction is the essay presented, and judgement is the choices defended under questioning. The climb from a passing to a strong essay is mostly judgement, which is why Sessions 10 and 11 were disproportionately about Parts III and IV.

\section{After the course}

This pipeline reappears in three places. First, in the master's thesis: the method backbone built here --- question, background, justified methods, evaluation, ethics --- is its method chapter, and should be reused deliberately. Secondly, in further studies: a doctorate is this pipeline, scaled up, and its vocabulary and checkpoints are now familiar. Thirdly, in working life: post-incident reviews, A/B tests, security assessments and vendor-claim reviews all rest on the question \emph{how do they know?} and on the ability to design the small study that answers it.


\section{Flashbacks and tables}

\subsection{The pipeline}
\begin{center}\small
\begin{center}
    \resizebox{0.97\textwidth}{!}{%
    \begin{tikzpicture}
      \node[kbox, text width=2.0cm, minimum height=1.15cm, anchor=west] (a) at (0,0)
        {\textbf{Problem \&}\\ \textbf{purpose}\\ \scriptsize S3};
      \node[kbox, text width=2.0cm, minimum height=1.15cm, anchor=west] (b) at (2.5,0)
        {\textbf{Research}\\ \textbf{question}\\ \scriptsize S4};
      \node[kbox, text width=2.0cm, minimum height=1.15cm, anchor=west] (c) at (5.0,0)
        {\textbf{Literature}\\ \textbf{review}\\ \scriptsize S5};
      \node[kbox, text width=2.0cm, minimum height=1.15cm, anchor=west] (d) at (7.5,0)
        {\textbf{Strategy}\\ \scriptsize S6--7};
      \node[kbox, text width=2.0cm, minimum height=1.15cm, anchor=west] (e) at (10.0,0)
        {\textbf{Methods \&}\\ \textbf{analysis}\\ \scriptsize S8--9};
      \node[kbox, text width=2.0cm, minimum height=1.15cm, anchor=west] (f) at (12.5,0)
        {\textbf{Ethics \&}\\ \textbf{report}\\ \scriptsize S10--11};
      \draw[karrow] (a)--(b); \draw[karrow] (b)--(c);
      \draw[karrow] (c)--(d); \draw[karrow] (d)--(e); \draw[karrow] (e)--(f);
      \node[gbox, text width=6.0cm, minimum height=0.8cm] (base) at (7.2,-1.7)
        {\textbf{Approach \& worldview} (S2) --- underneath every step};
      \draw[garrow] (base.north) -- (7.2,-0.62);
    \end{tikzpicture}}
  \end{center}
  
  {\color{uiared}\small\textbf{You are here:} every box has been filled at
  least once --- the essay is the pipeline written down.}\\[2pt]
  The same pipeline is your essay's table of contents, your master's thesis's
  method backbone, and --- later --- the skeleton of any empirical work you do.
\end{center}


\subsection{Flashbacks}
\begin{center}\small
\footnotesize
  \renewcommand{\arraystretch}{1.35}
  \begin{tabular}{@{}p{3.1cm}p{9.9cm}@{}}
    \toprule
    Quantitative & survey, n = 300 app users: which measured factors predict abandonment?\\
    Qualitative & 10 interviews: what does abandoning the app \emph{mean} to users?\\
    Mixed & QUANT $\rightarrow$ qual: measure the pattern, then explain it\\
    \bottomrule
  \end{tabular}
  

  One phenomenon, three legitimate studies --- the \textbf{question wording}
  decides which. If your Part I question and Part III strategy imply
  different rows of this table, fix one of them this session.
\end{center}

\begin{center}\small
\small
  \renewcommand{\arraystretch}{1.35}
  \begin{tabular}{@{}p{3.0cm}p{10.0cm}@{}}
    \toprule
    Purpose & the problem, its owner, its cost --- and the so-that\\
    Products & what new knowledge (and possibly artefact) comes out\\
    Process & the pipeline: question $\rightarrow$ \dots\ $\rightarrow$ report\\
    Participants & whose time and data you borrow --- and their rights\\
    Paradigm & the worldview underneath the choices\\
    Presentation & the essay now; the master's thesis next\\
    \bottomrule
  \end{tabular}
  

  Final use of the 6Ps: as a completeness check --- read your draft and tick
  each P. An unticked P is a missing section, not a style issue.
\end{center}

\begin{center}\small
\small
  \begin{itemize}\setlength{\itemsep}{0.45em}
    \item \textbf{Too broad} --- a career, not a study
    \item \textbf{Trivial yes/no} --- answerable without data
    \item \textbf{Not researchable} --- a value judgement in disguise
    \item \textbf{Double-barreled} --- two studies wearing one question mark
    \item \textbf{Assumption baked in} --- presumes the effect it should test
  \end{itemize}
  
  Final check: read your main question aloud and try to accuse it of each
  failure. If one accusation sticks even slightly, the repair patterns from
  session 4 are in the session 4 slides.
\end{center}

\begin{center}\small
\begin{center}
    \resizebox{0.95\textwidth}{!}{%
    \begin{tikzpicture}
      \node[kbox, text width=2.4cm, minimum height=0.95cm, anchor=west] (a) at (0,0)   {keyword table};
      \node[kbox, text width=2.4cm, minimum height=0.95cm, anchor=west] (b) at (3.4,0) {string + log};
      \node[kbox, text width=2.4cm, minimum height=0.95cm, anchor=west] (c) at (6.8,0) {criteria + PRISMA};
      \node[kbox, text width=2.4cm, minimum height=0.95cm, anchor=west] (d) at (10.2,0) {concept matrix};
      \node[kbox, text width=2.4cm, minimum height=0.95cm, anchor=west] (e) at (13.6,0){four moves};
      \draw[karrow] (a)--(b); \draw[karrow] (b)--(c); \draw[karrow] (c)--(d); \draw[karrow] (d)--(e);
    \end{tikzpicture}}
  \end{center}
  
  Five artefacts, and Part II assembles itself around them. The most common
  gap at this stage: the matrix exists but the \textbf{moves} were never
  written --- the review lists instead of arguing. One writing session fixes
  it.
\end{center}

\begin{center}\small
\footnotesize
  \renewcommand{\arraystretch}{1.3}
  \begin{tabular}{@{}p{3.4cm}p{9.6cm}@{}}
    \toprule
    Survey & breadth: patterns across many, sampling decides the claim\\
    Design \& creation & learning via making --- the research lives in the evaluation\\
    Experiment & the only strategy that earns the word \emph{because}\\
    Case study & depth in one setting; analytical generalisation\\
    Action research & change practice in cycles, with the practitioners\\
    Ethnography & immersion; the unwritten rules\\
    \bottomrule
  \end{tabular}
  

  Part III check: your paragraph should read as \emph{one chosen, others
  rejected with reasons} --- not as this table copied out.
\end{center}

\begin{center}\small
\small
  \begin{columns}[T]
    \column{0.48\textwidth}
    \textbf{Quantitative (S8)}
    \begin{itemize}\setlength{\itemsep}{0.25em}
      \item variables + measurement levels
      \item descriptives per variable
      \item test with a \emph{because}
      \item tool named
    \end{itemize}
    \column{0.48\textwidth}
    \textbf{Qualitative (S9)}
    \begin{itemize}\setlength{\itemsep}{0.25em}
      \item participants + saturation logic
      \item guide/protocol in appendix
      \item thematic analysis, code book
      \item second-coder procedure
    \end{itemize}
  \end{columns}
  
  Both templates end the same way: \textbf{limitations, named}. If your
  Part IV has numbers but no plan, or a plan but no limits, today's sprint
  starts there.
\end{center}

\begin{center}\small
\small
  \begin{enumerate}\setlength{\itemsep}{0.45em}
    \item What data do we collect --- and which of it is \textbf{personal},
          directly or indirectly?
    \item \textbf{SIKT}: needed? Decision documented either way?
    \item Consent: informed, voluntary, withdrawable --- and which promise:
          anonymous or confidential?
    \item The \textbf{client trap}: can an internal reader identify a
          speaker?
    \item Special zones touched? Public data, deception, security testing,
          AI tools?
  \end{enumerate}
  
  Five questions, five sentences of answers --- that is a complete Part IV
  ethics discussion at master's scale.
\end{center}


\subsection{The running example, assembled}
\begin{center}\small
\footnotesize
  \renewcommand{\arraystretch}{1.3}
  \begin{tabular}{@{}p{3.2cm}p{10.2cm}@{}}
    \toprule
    Problem (S3) & employees reuse passwords despite training --- SMBs carry the breach risk\\
    Purpose (S3) & to identify which factors drive reuse among SMB employees, so training can target them\\
    Question (S4) & which factors are associated with password reuse among SMB employees?\\
    Sub-questions & SQ1 how widespread is reuse (descriptive); SQ2 how do employees explain it (qualitative)\\
    Literature (S5) & string \texttt{("password reuse" OR "credential reuse") AND (employee* OR staff) AND (SMB OR SME)}; log, criteria, PRISMA; matrix shows fatigue converging, risk perception thin\\
    \bottomrule
  \end{tabular}
  

  Note how each row \textbf{consumes} the one above --- the golden thread in
  action.
\end{center}

\begin{center}\small
\footnotesize
  \renewcommand{\arraystretch}{1.3}
  \begin{tabular}{@{}p{3.2cm}p{10.2cm}@{}}
    \toprule
    Strategy (S6--7) & survey (SQ1) + interviews (SQ2) = mixed, explanatory sequential; case study rejected: question targets factors, not one setting\\
    Data (S8) & questionnaire, n $\approx$ 150 via six partner firms; validated fatigue scale; Nettskjema\\
    Data (S9) & 8 semi-structured interviews, saturation logic; thematic analysis, code book, two transcripts double-coded\\
    Analysis (S8) & descriptives; t-test trained vs.\ untrained; effect size reported; JASP\\
    Ethics (S10) & personal data: yes (e-mail, audio) $\rightarrow$ SIKT; consent per template; confidentiality, client-trap mitigations\\
    Limitations & convenience sample, self-report; named in Part IV\\
    \bottomrule
  \end{tabular}
  

  This table \emph{is} a essay in miniature. Your essay is this,
  written out --- with the justifications attached.
\end{center}


\subsection{Mismatch clinic}
\begin{center}\small
\small For each: where does the thread break, and what is the cheapest fix?
  
  \begin{enumerate}\setlength{\itemsep}{0.5em}
    \item Question: ``how do users \emph{experience} the new login flow?''
          --- Method: telemetry funnel analysis only.
    \item Question: ``does the caching layer reduce p95 latency?'' ---
          Strategy: interviews with the platform team.
    \item Question: ``which factors predict design-system adoption?'' ---
          Data: 5 interviews in one team.
  \end{enumerate}
  
  {\color{kgrey}\small (1) experience needs voices --- add interviews or
  reword to behaviour; (2) a causal system question needs the experiment
  from session 6; (3) \emph{predict} implies breadth --- survey, or reword to
  \emph{how does adoption happen here}.}
\end{center}


\subsection{Submission, FAQ and grades}
\begin{center}\small
\begin{itemize}\setlength{\itemsep}{0.45em}
    \item \textbf{Format:} about 10 pages excl.\ front page and references; the
          four parts clearly marked
    \item \textbf{Language:} English or Norwegian; PDF
    \item \textbf{Weights:} written essay 40\,\%; presentation and discussion
          in the oral exam 20\,\%
    \item \textbf{References:} APA 7; cite the curriculum where your choices
          lean on it
    \item \textbf{AI declaration:} submitted with the essay --- honest and specific
    \item \textbf{Deadline, portal and oral exam dates:} to be announced ---
          check now, not the night before
  \end{itemize}
  
  \begin{block}{The oral exam in one line}
    You present the essay and defend its choices in discussion: the same
    describe--justify pairs, spoken.
  \end{block}
\end{center}

\begin{center}\small
\small
  \renewcommand{\arraystretch}{1.35}
  \begin{tabular}{@{}p{6.0cm}p{7.0cm}@{}}
    \toprule
    \textbf{Question} & \textbf{Answer}\\
    \midrule
    Do tables count toward the 10 pages? & Yes; use tables for density, not for evasion\\
    How many references are enough? & No quota: every claim that needs support has it; the curriculum is cited at the choices\\
    Can we change the research question now? & Yes --- if Parts II--IV are updated to match; a late but consistent question beats an early broken thread\\
    We wrote as a group --- does everyone need to know everything? & Yes: in the oral exam each of you must be able to defend every part\\
    Our SIKT answer is still pending & Say so in Part IV with the submission date --- the plan is graded, not the approval letter\\
    \bottomrule
  \end{tabular}
\end{center}

\begin{center}\small
\small
  \renewcommand{\arraystretch}{1.35}
  \begin{tabular}{@{}p{3.2cm}p{4.6cm}p{4.9cm}@{}}
    \toprule
    & \textbf{Reproduction} & \textbf{Judgement}\\
    \midrule
    \textbf{Parts I--II} & Background and literature reported correctly & One question with a hierarchy; a synthesis that argues\\
    \textbf{Parts III--IV} & Strategy and methods named and described & Choices justified, alternatives rejected with reasons, limitations and ethics concrete\\
    \textbf{Oral exam} & The essay presented & The choices defended under questioning\\
    \bottomrule
  \end{tabular}
  
  Read the columns left to right: the climb from a passing to a strong essay
  is mostly \textbf{judgement} --- justification, limitations, ethics. That is
  why sessions 10--11 were disproportionately about Parts III--IV.
\end{center}



\section*{Questions from the reading}

  \begin{enumerate}\setlength{\itemsep}{0.4em}
    \item Name the \textbf{6Ps} --- and say which two your essay is most likely to under-specify.
          \answer{Purpose, products, process, participants, paradigm, presentation. Products and paradigm are usually implicit in drafts; both must be stated.}
    \item What is the difference between a research \textbf{strategy} and a research \textbf{method}?
          \answer{The strategy is the overall approach to answering the question (survey, design and creation, experiment, case study, action research, ethnography); the methods are the techniques for generating and analysing data (interviews, questionnaires, observation, documents; quantitative or qualitative analysis).}
    \item Which \textbf{evaluation criteria} does Oates attach to a literature review?
          \answer{Purpose stated; sources of appropriate quality; search documented and repeatable; literature critically evaluated and synthesised rather than listed; the review connected to the study's own question.}
    \item What makes a good \textbf{research question}, in Oates' terms?
          \answer{It follows from the purpose, it is precise and researchable with the resources available, it implies a strategy, and it can be answered --- or, for an essay, it can be shown how it would be answered.}
  \end{enumerate}

\section*{Questions for discussion}

  \begin{enumerate}\setlength{\itemsep}{0.25em}
    \item \textbf{Your golden thread} --- read your draft from the research question to the data plan: does every choice trace back to the question?
          \answer{Where the thread breaks is where the examiner's first question will come from. Typical break: a strategy chosen in Part III that the question in Part I does not require.}
    \item \textbf{The mismatch clinic} --- a survey to understand why administrators ignore alerts; an experiment with no control group; an ethnography of three weeks. What is wrong with each?
          \answer{The first asks a why-question with a breadth strategy; the second has no comparison and therefore no cause; the third lacks the immersion the strategy needs. Each is a strategy--question mismatch, the most common failure in Part III.}
    \item \textbf{One sentence for the oral exam} --- if the examiners asked only ``why this method?'', what would you answer?
          \answer{The answer is the describe--justify pair from Part III, spoken: the question asks for X; this strategy produces X; the alternatives would not, because \dots\ If that sentence exists, the oral exam is a conversation, not an interrogation.}
  \end{enumerate}

\section*{Key points}
\begin{itemize}
  \item Twelve sessions, one pipeline: from problem to question to literature to strategy to methods to ethics to essay
  \item The running example assembled end to end --- and the mismatches to avoid
  \item Submission: about 10 pages, four parts, APA 7, AI declaration; dates to be announced
  \item What separates the grades: judgement, not reproduction
  \item Where this reappears: your master's thesis, further studies, working life
\end{itemize}

\section*{Sources}
\begin{itemize}\small
  \item Oates, B.\,J., Griffiths, M., \& McLean, R. (2022). \emph{Researching information systems and computing} (2nd ed.). SAGE. --- the whole book.
\end{itemize}

\end{document}
